\documentclass[11pt]{article}
\usepackage[margin=1in]{geometry}
\IfFileExists{lmodern.sty}{\usepackage[T1]{fontenc}\usepackage{lmodern}}{\DeclareTextCommand{\k}{OT1}[1]{\c{##1}}}% ogonek fallback when Latin Modern is not installed
\usepackage{amsmath,amssymb,amsthm,mathtools}
\usepackage{stmaryrd}
\usepackage{booktabs}
\usepackage{longtable,array}
\usepackage[hidelinks]{hyperref}
\usepackage[capitalise,nameinlink]{cleveref}

% ---------- theorem environments ----------
\newtheorem{theorem}{Theorem}[section]
\newtheorem{lemma}[theorem]{Lemma}
\newtheorem{proposition}[theorem]{Proposition}
\newtheorem{corollary}[theorem]{Corollary}
\newtheorem{fact}[theorem]{Fact}
\theoremstyle{definition}
\newtheorem{definition}[theorem]{Definition}
\newtheorem{remark}[theorem]{Remark}
\newtheorem{example}[theorem]{Example}

% ---------- notation ----------
\newcommand{\Fq}{\mathbb{F}_q}
\newcommand{\K}{\mathbb{K}}
\newcommand{\N}{\mathbb{N}}
\newcommand{\iv}[2]{\llbracket #1, #2 \rrbracket}
\newcommand{\wt}{\mathrm{wt}}
\newcommand{\eq}{\mathrm{eq}}
\newcommand{\RS}{\mathrm{RS}}
\newcommand{\Enc}{\mathrm{Enc}}
\newcommand{\fold}{\mathrm{fold}}
\newcommand{\coef}[2]{[X^{#1}]\,#2}
\newcommand{\sub}{\preceq}
\newcommand{\cmp}[1]{\overline{#1}}
\newcommand{\Cc}{\mathcal{C}}
\newcommand{\Pp}{\mathcal{P}}
\newcommand{\Dd}{\mathcal{D}}
\newcommand{\lo}{\mathrm{lo}}
\newcommand{\res}{\mathrm{res}}
\newcommand{\ev}{\mathrm{ev}}
\newcommand{\F}{\mathbb{F}}
\newcommand{\hi}{\mathrm{hi}}
\newcommand{\Pieval}{\Pi_{\mathrm{eval}}}

\title{The Boolean-Kernel Basis:\\ Native Evaluation-Form FRI Folding over Multiplicative Domains,\\ with an Application to Multilinear Polynomial Commitments}
\author{Abdelali Mkhida\\ Algorizk Labs\\ \texttt{ali.mkhida@algorizk.xyz}}
\date{September 2026}

\begin{document}
\maketitle

\begin{abstract}
Multilinear polynomial commitment schemes in the BaseFold paradigm commit to a Reed--Solomon encoding of a multilinear polynomial and prove evaluations by interleaving a sumcheck with FRI-style folding.
Over smooth multiplicative domains, the classical fold acts on monomial coefficients, so a prover holding the polynomial as a table of values on the Boolean hypercube must first convert it to coefficient form.
We introduce the \emph{Boolean-kernel basis} $K_b(X) = \prod_{k=1}^{n}(X^{2^{k-1}} + b_k)$, $b \in \{0,1\}^n$, of the univariate polynomials of degree $< 2^n$.
Its change of basis to monomials is the Kronecker power of a $2\times 2$ matrix and costs $\tfrac{n}{2}2^n$ additions or subtractions.
In this basis a suitably normalised FRI fold is \emph{exactly} the restriction of one variable of the multilinear extension of the coordinate table; hence folding the Reed--Solomon encoding of $\sum_b f(b)K_b$ evaluates the multilinear extension $\tilde f$ of the table $f$ directly.
We also characterise low degree in kernel coordinates by a sign condition and express univariate evaluation as a scaled multilinear evaluation.
Building on this fold dictionary, we describe a transparent, hash-based multilinear polynomial commitment scheme.
We prove it sound, evaluation binding and round-by-round knowledge sound in the unique-decoding regime, with a self-contained argument whose only external ingredient about codes is the correlated-agreement theorem for Reed--Solomon codes (together with efficient unique decoding, for the extractor). From round-by-round knowledge soundness we derive the relaxed round-by-round knowledge soundness of Chiesa, Di, Hu and Zheng; by their theorem on the BCS transformation for interactive oracle reductions, the non-interactive scheme, in which the commitment is a Merkle root chosen by the adversary, is knowledge sound, for a single opening, against quantum adversaries in the quantum random oracle model, with explicit concrete bounds.
We implement the scheme in Rust over the Goldilocks field and compare it with a coefficient-form baseline sharing all other code: the two encodings are indistinguishable in running time at the resolution of our machine, so the structural property comes at no measurable cost.
The paper is self-contained: every notion used in the proofs is defined in the paper, every result specific to the paper is proved, and the external results are stated in \cref{sec:prelim,sec:proofs}; the post-quantum corollary uses the quantum extraction model of~\cite{CDHZ25}, whose definitions we restate only as far as the statement requires. Companion Lean~4 files formalise the algebraic, probabilistic and soundness results, with the correlated-agreement theorem as the only axiom.
\end{abstract}

\tableofcontents

\section{Introduction}\label{sec:intro}

Many succinct arguments reduce the correctness of a computation to a claim about a multilinear polynomial evaluated at a random point.
This is the case for sumcheck-based systems such as Spartan~\cite{Set20} and HyperPlonk~\cite{CBBZ23}.
Such systems need a \emph{multilinear polynomial commitment scheme}: the prover commits to a multilinear polynomial $\tilde f$ in $n$ variables and later proves that $\tilde f(z) = v$ for a point $z$ chosen by the verifier.
When the scheme is built from hash functions and error-correcting codes, the resulting arguments are transparent and rely on no number-theoretic assumption, which makes them candidates for post-quantum security.

A multilinear polynomial can be given in two forms.
In \emph{table form} (also called evaluation or Lagrange form) it is given by its values $f(b)$ on the Boolean hypercube $B_n = \{0,1\}^n$; this is the form in which provers naturally hold their witnesses.
In \emph{coefficient form} it is given by its monomial coefficients.
The two forms are related by M\"obius inversion on the Boolean lattice.

Code-based schemes in the BaseFold paradigm~\cite{ZCF24} commit to an encoding of $f$ and prove an evaluation by running a sumcheck whose round challenges also drive an FRI-like folding of the codeword~\cite{BBHR18}.
With Reed--Solomon codes on smooth multiplicative subgroups, as in DeepFold~\cite{GLHQTZ25}, the classical FRI fold $U \mapsto U_e + \theta\, U_o$ acts on coefficient vectors.
It computes the multilinear evaluation of the \emph{coefficient} vector: the monomial coefficients of the committed univariate polynomial must be the monomial coefficients of $\tilde f$.
A prover that holds a table must therefore first move to coefficient form, and must keep the two views consistent across the committed word, the sumcheck state and the folds.

This paper asks whether there is a basis of univariate polynomials in which FRI-style folding over a smooth multiplicative domain acts \emph{directly} on the table form.
There is, and it is explicit.

\subsection{Contributions}\label{sec:intro:contrib}

\paragraph{The Boolean-kernel basis (\cref{sec:kernel}).}
For $y \in \Fq^n$ let $K_y(X) = \prod_{k=1}^{n}(X^{2^{k-1}} + y_k)$.
We show that the $N = 2^n$ polynomials $K_b$, $b \in B_n$, which have coefficients in $\{0,1\}$, form a basis of $\F[X]_{<N}$ over every field $\F$.
The change of basis to monomials is the Kronecker power $C^{\otimes n}$ of $C = \bigl(\begin{smallmatrix} 0 & 1 \\ 1 & 1\end{smallmatrix}\bigr)$.
It costs exactly $\tfrac{n}{2}N$ additions or subtractions in either direction, and no multiplications.
In this basis, having degree $< 2^t$ is a sign condition on the coordinates (\cref{thm:lowdeg}).
For all but at most $N-1$ points $\zeta$, evaluating $U$ at $\zeta$ equals, up to an explicit nonzero scalar, evaluating the multilinear extension of its coordinates at a point on an explicit curve (\cref{prop:eval}).

\paragraph{The fold dictionary (\cref{sec:fold}).}
Our main structural result is \cref{thm:fold-dictionary}.
Let $U = \sum_b \lambda(b) K_b$ and $U = U_e(X^2) + X U_o(X^2)$. Then the kernel coordinates of the fold
\[
  \mathrm{pfold}_r(U) \;=\; (2r-1)\,U_e + (1-r)\,U_o
\]
are exactly $(1-r)\lambda(2b') + r\lambda(2b'+1)$, the restriction of the first variable of the table $\lambda$.
Consequently, encoding a table $f$ as the Reed--Solomon codeword of $U_f = \sum_b f(b) K_b$ and folding it $n$ times with challenges $r_1,\dots,r_n$ yields the constant $\tilde f(r_1,\dots,r_n)$ (\cref{cor:codeword-fold}), with no change of representation at any point.
The kernel fold agrees with the classical FRI fold up to a nonzero scalar and a reparametrisation of the challenge (\cref{sec:fold:fri}); mixing the two parametrisations breaks correctness (\cref{ex:pitfall}), a mistake we made in an early implementation.

\paragraph{A multilinear PCS with a self-contained soundness proof (\cref{sec:pcs,sec:sound}).}
We instantiate the BaseFold paradigm with the kernel encoding and prove it sound in the unique-decoding regime (\cref{thm:sound}).
The only external ingredient about codes in the soundness proof is the correlated-agreement theorem for Reed--Solomon codes of Ben-Sasson, Carmon, Ishai, Kopparty and Saraf~\cite{BCIKS23}; the efficiency of the extractor also uses polynomial-time unique decoding of Reed--Solomon codes.
The theorem enters through a one-line observation: the kernel fold of any word is an affine line $u^0 + r\,u^1$ whose generators determine the word (\cref{lem:line}).
The proof tracks the set of passing query chains level by level, and it yields evaluation binding (\cref{cor:binding}).
The same argument gives round-by-round knowledge soundness with an explicit decoding extractor (\cref{thm:rbr}).
The BCS theorem in the quantum random oracle model of Chiesa, Manohar and Spooner~\cite{CMS19} does not cover instances containing a Merkle-committed oracle, such as the commitment. Chiesa, Di, Hu and Zheng~\cite{CDHZ25} extend the BCS transformation to interactive oracle reductions, whose instances may contain such oracles, under a relaxed round-by-round knowledge soundness hypothesis. We derive this hypothesis from \cref{thm:rbr} (\cref{thm:rrbr}); their theorem then makes the non-interactive scheme knowledge sound, for a single opening, against quantum adversaries that choose the commitment, with explicit concrete bounds (\cref{thm:cdhz,cor:pq}); \cref{rem:pq-params} gives parameters for this setting.

\paragraph{Implementation and experiments (\cref{sec:impl,sec:exp}).}
We give a Rust implementation over the Goldilocks field with a quadratic extension for challenges.
We compare it with a coefficient-form baseline that shares all other code.
The two encodings are indistinguishable in running time at the resolution of our machine, as the operation counts predict.
The kernel encoding thus provides its structural property at no measurable cost.

\paragraph{What we do not claim.}
The scheme is not asymptotically or concretely faster than Reed--Solomon instantiations of BaseFold or DeepFold (\cref{rem:basefold}).
Its evaluation protocol is BaseFold's sumcheck-plus-folding, applied to a different encoding.
Our soundness analysis is restricted to the unique-decoding regime, which needs more queries than list-decoding analyses~\cite{Hab24,Zei26}.
Zero knowledge is out of scope, and we analyse one commitment opened at one point.
The post-quantum statement (\cref{cor:pq}) uses the theorem of~\cite{CDHZ25}, which we quote and do not reprove, and it applies to their compiler (salted Merkle trees, their challenge derivation), not to our implementation as it stands (\cref{rem:fs}).
The contribution is the basis and the dictionary: a precise, elementary account of why FRI folding and table-form multilinear restriction are the same operation over multiplicative domains in odd characteristic.

\subsection{Related work}\label{sec:intro:related}

\paragraph{Proximity tests for Reed--Solomon codes.}
FRI~\cite{BBHR18} tests proximity to Reed--Solomon codes by repeated even--odd folding.
Its soundness in the unique-decoding and list-decoding regimes rests on proximity gaps for Reed--Solomon codes~\cite{BCIKS23}.
WHIR~\cite{ACFY24} tests proximity to constrained Reed--Solomon codes with very fast verification and supports multilinear evaluation claims.
Our soundness proof uses only the unique-decoding correlated-agreement theorem of~\cite{BCIKS23}; extending it to the list-decoding regime is the most natural next step.

\paragraph{Multilinear commitments from folding.}
BaseFold~\cite{ZCF24} introduced foldable codes and the interleaving of sumcheck rounds with folding rounds under shared challenges.
Our evaluation protocol follows it directly.
DeepFold~\cite{GLHQTZ25} instantiates the paradigm with Reed--Solomon codes and improves proof size; its folding computes coefficient-form evaluations, and our coefficient-form baseline of \cref{sec:exp} follows the same principle.
Hab\"ock~\cite{Hab24} analyses BaseFold over Reed--Solomon codes up to the Johnson bound.
Zeilberger~\cite{Zei26} proves proximity gaps for multilinear evaluation for all linear codes.
Both results concern the soundness of the folding paradigm; our contribution is orthogonal and concerns the encoding.

\paragraph{Binary fields and the novel polynomial basis.}
FRI-Binius~\cite{DP24} adapts BaseFold to binary towers using the novel polynomial basis of Lin, Chung and Han~\cite{LCH14}.
In that basis the additive FFT, and with it FRI-style folding over additive subspaces, interacts cleanly with multilinear structure.
The Boolean-kernel basis plays an analogous role for smooth \emph{multiplicative} subgroups in \emph{odd} characteristic.
The two constructions are genuinely different. In characteristic $2$ every factor $X^{2^{k-1}} + 1$ equals $(X+1)^{2^{k-1}}$, so $K_b = X^{N-1-b}(X+1)^{b}$ is a Bernstein-type basis; moreover the multiplicative group of a field of characteristic $2$ has odd order, so smooth multiplicative domains do not exist, and the even--odd formulas~\eqref{eq:even-odd} divide by $2$.

\paragraph{Multilinear-to-univariate reductions.}
A different line of work proves multilinear evaluations with a \emph{univariate} polynomial commitment, typically KZG.
Gemini~\cite{BCHO22} reduces a coefficient-form multilinear claim to univariate claims through recursive even--odd splitting.
Zeromorph~\cite{KT24} works in table form and uses quotients and degree checks.
Mercury~\cite{EG25} reaches constant proof size and expresses the evaluation as a coefficient of a product with a tensor-structured polynomial. Samaritan~\cite{GPS25} improves on Gemini in proof size, prover cost and verifier cost.
Mohamed~\cite{Moh25} relates univariate and multilinear sumchecks through the inverse Kronecker substitution.
These schemes treat the univariate commitment as a black box.
In ours, the univariate polynomial is committed through a Reed--Solomon code, and the basis is chosen so that the code's own folding carries the multilinear structure.

\paragraph{Other domains.}
Brakedown~\cite{GLSTW23} and related schemes avoid FFTs by using linear-time encodable codes.
Circle STARKs~\cite{HLP24} obtain smooth folding domains over the Mersenne prime $2^{31}-1$ from the circle curve.
The fold dictionary as stated needs a smooth multiplicative subgroup in odd characteristic.
Whether a Boolean-kernel-type basis exists for circle domains and for additive domains is an open question that we leave to future work.

\subsection{Organisation}\label{sec:intro:org}

\Cref{sec:prelim} fixes notation and collects the facts on multilinear polynomials and Reed--Solomon codes that we use, including the correlated-agreement theorem.
\Cref{sec:proofs} recalls interactive oracle proofs and their probability space, round-by-round (knowledge) soundness, the sumcheck protocol, polynomial commitments, and the definitions and the theorem of~\cite{CDHZ25} on the BCS transformation for interactive oracle reductions, which we use for the non-interactive scheme.
\Cref{sec:kernel} develops the Boolean-kernel basis, and \cref{sec:fold} proves the fold dictionary.
\Cref{sec:pcs} defines the commitment scheme and \cref{sec:sound} proves its soundness.
\Cref{sec:impl,sec:exp} describe the implementation and the measurements, and \cref{sec:concl} lists open problems.
\Cref{app:notation} collects the notation.

\section{Preliminaries}\label{sec:prelim}

This section fixes notation and collects the algebraic facts used throughout.
Every result specific to this paper is proved.
Standard facts from algebra and coding theory are quoted with a reference and labelled \emph{Fact}; the correlated-agreement theorem of \cref{sec:prelim:gaps} is the only external result about codes used in the soundness analysis, and \cref{fact:decode} is used only for the efficiency of the extractor.
\Cref{app:notation} lists the notation in one table.

\subsection{Integers, bits and points}\label{sec:prelim:bits}

\paragraph{Integers.}
$\N = \{0,1,2,\dots\}$.
For integers $a$ and $c$ we write $\iv{a}{c} = \{ i \in \mathbb{Z} : a \le i \le c \}$, which is empty if $a > c$.

\paragraph{Number of variables.}
The letter $n \in \N$ denotes a \emph{number of variables}, and $N = 2^n$.
Unless a hypothesis states otherwise, every statement of \cref{sec:prelim,sec:kernel,sec:fold} holds for every $n \ge 0$, and is applied with various numbers of variables (for instance $n-1$ or $n-j$); when the number of variables must be explicit, we write it as a superscript, as in $K^{(n)}_b$ (\cref{def:kernel}).

\paragraph{Binary decomposition.}
Every $b \in \iv{0}{N-1}$ has a unique \emph{binary decomposition}
\[
  b \;=\; \sum_{k=1}^{n} b_k\, 2^{k-1}, \qquad b_k \in \{0,1\},
\]
and $(b_1,\dots,b_n)$ is its \emph{bit vector}; bit $b_1$ is the least significant one.
Writing $B_n = \{0,1\}^n$ for the \emph{Boolean hypercube}, the map $b \mapsto (b_1,\dots,b_n)$ is a bijection $\iv{0}{N-1} \to B_n$, and we identify $b$ with its bit vector.
For a commutative ring $A$ with $1 \ne 0$, we also regard $B_n$ as a subset of $A^n$, via $\{0,1\} \subseteq A$.
For $b, c \in \iv{0}{N-1}$:
\begin{itemize}
  \item $\wt(b) = \sum_{k=1}^n b_k$ is the \emph{Hamming weight} of $b$;
  \item $\cmp{b} = (N-1) - b$ is the \emph{complement} of $b$, with bits $\cmp{b}_k = 1 - b_k$;
  \item $c \sub b$ means $c_k \le b_k$ for all $k \in \iv{1}{n}$; we then say that $c$ is a \emph{submask} of $b$.
\end{itemize}

\paragraph{Splitting indices.}
Let $t \in \iv{0}{n}$. Every $b \in \iv{0}{N-1}$ can be written uniquely as $b = c + 2^t h$ with $c \in \iv{0}{2^t-1}$ and $h \in \iv{0}{2^{n-t}-1}$, namely $c = b \bmod 2^t$ and $h = \lfloor b/2^t \rfloor$.
The bits are $c_k = b_k$ for $k \in \iv{1}{t}$ and $h_k = b_{t+k}$ for $k \in \iv{1}{n-t}$.
For $t = 1$ we write $b = b_1 + 2b'$, with $b' = \lfloor b/2 \rfloor$ of bit vector $(b_2,\dots,b_n)$.
When a letter is declared to denote an index, its subscripts denote its bits.

\paragraph{Points.}
Let $A$ be a commutative ring and $j \in \iv{0}{n}$.
For $x = (x_1,\dots,x_j) \in A^j$ and $y = (y_1,\dots,y_{n-j}) \in A^{n-j}$, we write $(x,y) = (x_1,\dots,x_j,y_1,\dots,y_{n-j}) \in A^n$.
For $x \in A^j$ and an index $b \in \iv{0}{2^{n-j}-1}$, $(x, b)$ is the point $(x_1,\dots,x_j,b_1,\dots,b_{n-j})$, using the bit vector of $b$ with $n-j$ bits.
For $r = (r_1,\dots,r_m) \in A^m$ and $j \in \iv{0}{m}$, we write $r_{\le j} = (r_1,\dots,r_j)$.

\subsection{Fields and univariate polynomials}\label{sec:prelim:fields}

Throughout the paper, $\Fq$ is a finite field with $q$ elements and \emph{odd characteristic}, so that $2$ is invertible in $\Fq$.
All results of \cref{sec:prelim,sec:kernel,sec:fold,sec:pcs,sec:sound} hold for every such field; in \cref{sec:sound} they are also applied to an extension field $\K \supseteq \Fq$.

For a field $\F$, $\F[X]$ is the ring of univariate polynomials over $\F$.
For $d \in \N$, $\F[X]_{<d}$ is the $\F$-vector space of polynomials of degree less than $d$ (with $\deg 0 = -\infty$); it has dimension $d$.
For $P \in \F[X]$ and $a \in \N$, $\coef{a}{P}$ denotes the coefficient of $X^a$ in $P$.

\begin{lemma}[Root counting]\label{lem:roots}
Let $\F$ be a field. A nonzero polynomial $P \in \F[X]$ of degree $d$ has at most $d$ roots in $\F$.
Consequently, two polynomials of $\F[X]_{<d}$ that agree on $d$ distinct points of $\F$ are equal.
\end{lemma}

\begin{proof}
If $P(a) = 0$, Euclidean division by $X - a$ gives $P = (X-a)Q$ with $\deg Q = d-1$, and every root of $P$ other than $a$ is a root of $Q$, because $\F$ has no zero divisors. Induction on $d$ gives the first statement. For the second, apply it to the difference, which has degree $< d$.
\end{proof}

\begin{fact}[{\cite[Theorem~2.8]{LN97}}]\label{fact:cyclic}
The multiplicative group $\F^{\times}$ of a finite field $\F$ is cyclic.
\end{fact}

\begin{lemma}[Splitting polynomials]\label{lem:splits}
Let $\F$ be a field.
\begin{enumerate}
  \item (Even--odd split.) Let $d \ge 1$. Every $P \in \F[X]_{<2d}$ can be written uniquely as
  \[
    P(X) \;=\; P_e(X^2) + X\,P_o(X^2), \qquad P_e, P_o \in \F[X]_{<d},
  \]
  namely with $\coef{i}{P_e} = \coef{2i}{P}$ and $\coef{i}{P_o} = \coef{2i+1}{P}$ for $i \in \iv{0}{d-1}$.
  \item (Block split.) Let $t \in \N$ and $m \ge 1$. Every $P \in \F[X]_{<2^t m}$ can be written uniquely as
  \[
    P(X) \;=\; \sum_{i=0}^{m-1} X^{2^t i}\, R_i(X), \qquad R_i \in \F[X]_{<2^t},
  \]
  namely with $\coef{a}{R_i} = \coef{2^t i + a}{P}$ for $a \in \iv{0}{2^t-1}$.
\end{enumerate}
Both maps $P \mapsto (P_e, P_o)$ and $P \mapsto (R_0,\dots,R_{m-1})$ are $\F$-linear.
\end{lemma}

\begin{proof}
(1) Each monomial $X^a$ with $a < 2d$ appears on the right-hand side exactly once: as $(X^2)^{i}$ in $P_e(X^2)$ if $a = 2i$, and as $X (X^2)^{i}$ in $X P_o(X^2)$ if $a = 2i+1$. Comparing coefficients gives existence and uniqueness.
(2) Likewise, every $a < 2^t m$ is written uniquely as $a = 2^t i + a'$ with $i \in \iv{0}{m-1}$ and $a' \in \iv{0}{2^t-1}$ (Euclidean division), and $X^a$ appears exactly once, in $X^{2^t i} R_i(X)$. Linearity is clear from the coefficient formulas.
\end{proof}

\subsection{Multilinear polynomials}\label{sec:prelim:ml}

In this subsection, $A$ is a commutative ring.
A polynomial in $A[x_1,\dots,x_n]$ is \emph{multilinear} if its degree in each variable $x_k$ is at most $1$.
The multilinear polynomials form an $A$-module, which we denote by $\mathcal{L}_n(A)$, and we write $\mathcal{L}_n = \mathcal{L}_n(\Fq)$.
Since the monomials form a basis of $A[x_1,\dots,x_n]$, the multilinear monomials $\prod_{k=1}^n x_k^{a_k}$, $a \in \iv{0}{N-1}$, form a basis of $\mathcal{L}_n(A)$; in particular, for $n \ge 1$ every $P \in \mathcal{L}_n(A)$ can be written uniquely as $P = P_0 + x_n P_1$ with $P_0, P_1 \in \mathcal{L}_{n-1}(A)$ in the variables $x_1,\dots,x_{n-1}$.
A polynomial with coefficients in $A$ can be evaluated at the points of $A'^n$ for every commutative ring $A' \supseteq A$.

\begin{definition}[Monomials, equality polynomial, Lagrange polynomials]\label{def:bases}
For $a \in \iv{0}{N-1}$, the \emph{canonical monomial} is
$m_a(x) = \prod_{k=1}^{n} x_k^{a_k}$.
The \emph{equality polynomials} are
\[
  \eq_1(y,z) \;=\; yz + (1-y)(1-z),
  \qquad
  \eq(y,z) \;=\; \prod_{k=1}^{n} \eq_1(y_k,z_k),
\]
in the variables $y = (y_1,\dots,y_n)$ and $z = (z_1,\dots,z_n)$.
For $b \in \iv{0}{N-1}$, the \emph{Lagrange polynomial} is $e_b(x) = \eq(b, x) = \prod_{k=1}^{n} \bigl( b_k x_k + (1-b_k)(1-x_k) \bigr)$.
\end{definition}

\begin{lemma}[Equality polynomial]\label{lem:eq}
\begin{enumerate}
  \item For $b, c \in B_n$, $\eq(b,c) = 1$ if $b = c$, and $\eq(b,c) = 0$ otherwise.
  \item For every commutative ring $A' \supseteq A$ and $z \in A'^n$, the polynomial $x \mapsto \eq(x, z)$ is a multilinear polynomial with coefficients in $A'$, and $\eq(x,z) = \eq(z,x)$.
  \item Let $j \in \iv{0}{n}$, let $x$ and $z$ be points with $j$ coordinates, and let $x'$ and $z'$ be points with $n-j$ coordinates. Then
  \[
    \eq\bigl((x,x'),(z,z')\bigr) \;=\; \eq(x,z)\,\eq(x',z'),
  \]
  where the right-hand side uses $j$ and $n-j$ variables.
  \item $\eq_1(0, r) = 1 - r$ and $\eq_1(1, r) = r$ for every $r \in A$.
\end{enumerate}
\end{lemma}

\begin{proof}
(1) Each factor $\eq_1(b_k,c_k)$ equals $1$ if $b_k = c_k$ and $0$ otherwise.
(2) Each factor $\eq_1(x_k,z_k)$ has degree $\le 1$ in $x_k$ and does not involve the other variables; and $\eq_1$ is symmetric.
(3) and (4) are immediate from the definition.
\end{proof}

\begin{proposition}[Multilinear extension]\label{prop:mle}
A \emph{table} over $A$ is a function $f : \iv{0}{N-1} \to A$.
For every table $f$ there is a unique $\tilde f \in \mathcal{L}_n(A)$ with $\tilde f(b) = f(b)$ for all $b \in B_n$, namely
\[
  \tilde f(x) \;=\; \sum_{b=0}^{N-1} f(b)\, e_b(x) \;=\; \sum_{b=0}^{N-1} f(b)\, \eq(b,x).
\]
It is the \emph{multilinear extension} of $f$.
The family $(e_b)_{b \in \iv{0}{N-1}}$ is a basis of the $A$-module $\mathcal{L}_n(A)$, and $f \mapsto \tilde f$ is an $A$-linear bijection from the tables over $A$ onto $\mathcal{L}_n(A)$.
\end{proposition}

\begin{proof}
By \cref{lem:eq}, the displayed polynomial is multilinear (item~2) and interpolates $f$ (item~1).
For uniqueness, it suffices to show that a multilinear $P$ vanishing on $B_n$ is zero, which we prove by induction on $n$.
For $n = 0$, $P$ is a constant and $B_0$ has one point.
For $n \ge 1$, write $P = P_0 + x_n P_1$ as above.
Setting $x_n = 0$ shows that $P_0$ vanishes on $B_{n-1}$, so $P_0 = 0$; setting $x_n = 1$ then shows that $P_1$ vanishes on $B_{n-1}$, so $P_1 = 0$.
For the basis statement: if $P \in \mathcal{L}_n(A)$, then $P - \sum_b P(b)\, e_b$ is multilinear and vanishes on $B_n$, hence is zero, so the $e_b$ generate $\mathcal{L}_n(A)$; and if $\sum_b a_b e_b = 0$, evaluating at $c \in B_n$ gives $a_c = 0$.
\end{proof}

If $f$ takes values in a subring $A_0 \subseteq A$, the formula of \cref{prop:mle} shows that the multilinear extension of $f$ over $A$ is its multilinear extension over $A_0$, with coefficients viewed in $A$; we write $\tilde f$ in both cases.

\begin{corollary}[Multilinear extension of the equality table]\label{cor:eq-mle}
For every commutative ring $A' \supseteq A$ and $z \in A'^n$, the multilinear extension of the table $b \mapsto \eq(b,z)$ over $A'$ is the polynomial $x \mapsto \eq(x,z)$.
\end{corollary}

\begin{proof}
By \cref{lem:eq}(2) the polynomial $x \mapsto \eq(x,z)$ is multilinear, and it takes the value $\eq(b,z)$ at $b \in B_n$; conclude by the uniqueness in \cref{prop:mle}.
\end{proof}

\begin{definition}[Coefficient form]\label{def:coeff-form}
The \emph{coefficient form} of $\tilde f$ is the family $(\alpha_a)_{a \in \iv{0}{N-1}}$ of its coefficients: $\tilde f = \sum_{a} \alpha_a\, m_a$.
The table $f$ itself is the family of coordinates of $\tilde f$ in the Lagrange basis; we call it the \emph{table form}.
\end{definition}

\begin{lemma}[Table form $\leftrightarrow$ coefficient form]\label{lem:mobius}
With the notation of \cref{def:coeff-form},
\begin{align*}
  f(b) &\;=\; \sum_{\substack{a \in \iv{0}{N-1} \\ a \sub b}} \alpha_a && \text{for all } b \in \iv{0}{N-1}, \\
  \alpha_a &\;=\; \sum_{\substack{b \in \iv{0}{N-1} \\ b \sub a}} (-1)^{\wt(a)-\wt(b)}\, f(b) && \text{for all } a \in \iv{0}{N-1}.
\end{align*}
\end{lemma}

\begin{proof}
For $b \in B_n$ we have $m_a(b) = \prod_{k : a_k = 1} b_k$, which equals $1$ if $a \sub b$ and $0$ otherwise; evaluating $\tilde f = \sum_a \alpha_a m_a$ at $b$ gives the first identity.
For the second, substitute the first into its right-hand side:
\[
  \sum_{b \sub a} (-1)^{\wt(a)-\wt(b)} \sum_{c \sub b} \alpha_c
  \;=\; \sum_{c \sub a} \alpha_c \sum_{c \sub b \sub a} (-1)^{\wt(a)-\wt(b)} .
\]
The inner sum runs over the subsets of the $\wt(a) - \wt(c)$ positions where $a$ has a $1$ and $c$ has a $0$; it equals $(1-1)^{\wt(a)-\wt(c)}$, which is $1$ if $c = a$ and $0$ otherwise.
\end{proof}

\begin{definition}[Restriction]\label{def:restriction}
Let $n \ge 1$, let $f$ be a table over $A$ with $n$ variables, and let $r \in A$.
The \emph{restriction} of $f$ at $r$ is the table $\res_r(f)$ over $A$ with $n-1$ variables defined by
\[
  \res_r(f)(b') \;=\; (1-r)\, f(2b') \;+\; r\, f(2b'+1), \qquad b' \in \iv{0}{N/2-1}.
\]
For $j \in \iv{0}{n}$ and $r = (r_1,\dots,r_j) \in A^j$, the \emph{iterated restriction} $\res_r(f)$ is defined by $\res_{()}(f) = f$ and $\res_{(r_1,\dots,r_j)}(f) = \res_{r_j}\bigl(\res_{(r_1,\dots,r_{j-1})}(f)\bigr)$; it is a table with $n-j$ variables.
For $j = 1$ the two definitions agree, identifying $A$ with $A^1$.
\end{definition}

The restriction is $A$-linear in $f$.

\begin{lemma}[Restriction evaluates the first variables]\label{lem:restriction}
Let $f$ be a table over $A$ with $n$ variables.
\begin{enumerate}
  \item If $n \ge 1$ and $r \in A$, then, as polynomials in $y = (y_1,\dots,y_{n-1})$,
  \[
    \widetilde{\res_r(f)}(y) \;=\; \tilde f(r, y).
  \]
  \item For $j \in \iv{0}{n}$ and $r \in A^j$, $\widetilde{\res_r(f)}(y) = \tilde f(r, y)$ as polynomials in $y = (y_1,\dots,y_{n-j})$.
  In particular $\res_r(f)(b) = \tilde f(r, b)$ for all $b \in \iv{0}{2^{n-j}-1}$, and for $j = n$ the table $\res_r(f)$ has the single entry $\tilde f(r)$.
\end{enumerate}
\end{lemma}

\begin{proof}
(1) Write $b = b_1 + 2b'$. By \cref{lem:eq}(3), $e_b(x_1, y) = \eq_1(b_1, x_1)\, e_{b'}(y)$, where $e_{b'}$ has $n-1$ variables.
Substituting $x_1 = r$, using \cref{lem:eq}(4) and grouping the terms $b_1 = 0$ and $b_1 = 1$,
\[
  \tilde f(r, y) \;=\; \sum_{b'=0}^{N/2-1} \bigl( (1-r) f(2b') + r f(2b'+1) \bigr)\, e_{b'}(y) \;=\; \sum_{b'} \res_r(f)(b')\, e_{b'}(y),
\]
which is $\widetilde{\res_r(f)}(y)$ by \cref{prop:mle}.
(2) Induction on $j$; the case $j = 0$ is trivial. For $j \ge 1$, by (1) applied to $\res_{r_{\le j-1}}(f)$ and the induction hypothesis,
$\widetilde{\res_{r}(f)}(y) = \widetilde{\res_{r_{\le j-1}}(f)}(r_j, y) = \tilde f(r_{\le j-1}, r_j, y)$.
Evaluating at $y = b \in B_{n-j}$ gives the second statement, since the multilinear extension interpolates its table.
\end{proof}

\begin{lemma}[Naturality]\label{lem:natural}
Let $\vartheta : A \to A'$ be a homomorphism of commutative rings, applied entrywise to tables and points and coefficientwise to polynomials.
For every table $f$ over $A$, $\vartheta(\tilde f) = \widetilde{\vartheta \circ f}$; and for $j \in \iv{0}{n}$ and $r \in A^j$, $\vartheta \circ \res_r(f) = \res_{\vartheta(r)}(\vartheta \circ f)$.
\end{lemma}

\begin{proof}
Both $\tilde f = \sum_b f(b) e_b$ and $\res_r(f)$ are given by polynomial formulas with integer coefficients in the entries of $f$, of $r$ and of the variables, and $\vartheta$ commutes with such formulas.
\end{proof}

\begin{remark}[Formal variables and extension fields]\label{rem:formal}
\Cref{lem:restriction} holds over an arbitrary commutative ring $A$.
Taking $A = \Fq[X]$ and $r = X$ shows that, for a table $f$ over $\Fq$, the polynomial $\tilde f(X, y)$ is the multilinear extension, in $y$, of the table $\res_X(f)$ over $\Fq[X]$; applying \cref{lem:natural} to the evaluation homomorphism $\Fq[X] \to \Fq$, $X \mapsto r$, recovers $\res_r(f)$ and $\tilde f(r, y)$.
Taking $A = \K$ covers challenges from an extension field.
\end{remark}

\subsection{Smooth domains and Reed--Solomon codes}\label{sec:prelim:rs}

\begin{definition}[Smooth domain]\label{def:domain}
A \emph{smooth domain} is a subgroup $L \le \Fq^{\times}$ of order $M = 2^{\mu}$ with $\mu \in \N$.
For $j \in \iv{0}{\mu}$, we write $L^{2^j} = \{\xi^{2^j} : \xi \in L\}$ for the image of the $2^j$-th power map.
In particular $L^2$ is the image of the squaring map; it is not a Cartesian power.
\end{definition}

\begin{lemma}[Power maps on smooth domains]\label{lem:power}
Let $\mu \in \N$.
\begin{enumerate}
  \item $\Fq^\times$ has a subgroup of order $2^\mu$ if and only if $2^\mu \mid q-1$; it is then unique and cyclic.
  \item Let $L$ be a smooth domain of order $M = 2^\mu$ and $j \in \iv{0}{\mu}$. The map $\xi \mapsto \xi^{2^j}$ is a group homomorphism from $L$ onto $L^{2^j}$, which is the smooth domain of order $M/2^j$, and every element of $L^{2^j}$ has exactly $2^j$ preimages in $L$.
  \item $(L^{2^j})^{2} = L^{2^{j+1}}$ for $j \in \iv{0}{\mu-1}$.
  \item If $\mu \ge 1$, then $-1 \in L$, $-1 \ne 1$, and the preimages of $\eta = \xi^2 \in L^2$ under squaring are exactly $\xi$ and $-\xi$, which are distinct.
\end{enumerate}
\end{lemma}

\begin{proof}
(1) A subgroup of order $e$ of a group of order $q-1$ exists only if $e \mid q-1$ (Lagrange's theorem).
Conversely, let $e \mid q-1$. By \cref{fact:cyclic}, $\Fq^\times$ is generated by some $\omega$ of order $q-1$, and $\omega^{(q-1)/e}$ generates a cyclic subgroup of order $e$.
Every subgroup of order $e$ is contained in the set of roots of $X^e - 1$, which has at most $e$ elements by \cref{lem:roots}; so there is exactly one such subgroup.
(2) The power map is a homomorphism because $L$ is commutative. If $\omega$ generates $L$, then $\omega^{2^j}$ generates the image and has order $2^{\mu-j}$; by (1) the image is the smooth domain of that order. The preimages of an element form a coset of the kernel, which has order $|L|/|L^{2^j}| = 2^j$.
(3) $(\xi^{2^j})^2 = \xi^{2^{j+1}}$.
(4) The kernel of squaring has order $2$ by (2) and consists of roots of $X^2 - 1 = (X-1)(X+1)$; since the characteristic is odd, $-1 \ne 1$, so the kernel is $\{1,-1\}$. The preimages of $\xi^2$ form the coset $\{\xi, -\xi\}$, and $\xi \ne -\xi$ because $2\xi \ne 0$.
\end{proof}

\begin{definition}[Fibres]\label{def:fibre}
Let $L$ be a smooth domain of order $M \ge 2$. The \emph{fibre} over $\eta \in L^2$ is the set $\{\xi, -\xi\}$ of its two square roots in $L$ (\cref{lem:power}(4)).
The $M/2$ fibres partition $L$.
\end{definition}

\begin{definition}[Reed--Solomon code]\label{def:rs}
Let $\F \supseteq \Fq$ be a finite field, let $L$ be a smooth domain of order $M$, and let $1 \le d \le M$.
The \emph{evaluation map} is $\ev_L : \F[X] \to \F^{L}$, $P \mapsto (P(\xi))_{\xi \in L}$, and the \emph{Reed--Solomon code} of degree bound $d$ on $L$ over $\F$ is
\[
  \RS_{\F}[L,d] \;=\; \ev_L\bigl(\F[X]_{<d}\bigr) \;\subseteq\; \F^{L}.
\]
We write $\RS[L,d]$ when $\F = \Fq$. The \emph{rate} is $\rho = d/M$.
Elements of $\F^L$ are \emph{words}, and elements of $\RS_\F[L,d]$ are \emph{codewords}.
For $w, w' \in \F^L$, the \emph{relative Hamming distance} is $\Delta(w,w') = |\{\xi \in L : w(\xi) \ne w'(\xi)\}|/M$; for a nonempty set $\Cc \subseteq \F^L$ we put $\Delta(w, \Cc) = \min_{c \in \Cc} \Delta(w,c)$.
We say that $w$ and $w'$ \emph{agree} on a set $S \subseteq L$ if $w(\xi) = w'(\xi)$ for all $\xi \in S$.
For finite fields $\F \subseteq \F'$ we regard $\F^L \subseteq \F'^L$ entrywise; then $\ev_L$ commutes with this inclusion and $\RS_{\F}[L,d] \subseteq \RS_{\F'}[L,d]$.
\end{definition}

\begin{lemma}[Hamming distance]\label{lem:metric}
$\Delta$ is a metric on $\F^L$. In particular, if $w$ agrees with $w'$ on at least $\alpha M$ points and $w'$ agrees with $w''$ on at least $\beta M$ points, then $w$ agrees with $w''$ on at least $(\alpha + \beta - 1)M$ points.
\end{lemma}

\begin{proof}
Symmetry, and $\Delta(w,w') = 0$ if and only if $w = w'$, are clear. If $w(\xi) \ne w''(\xi)$, then $w(\xi) \ne w'(\xi)$ or $w'(\xi) \ne w''(\xi)$; counting gives the triangle inequality. The second statement is the triangle inequality rewritten in terms of agreements, $1 - \Delta$.
\end{proof}

\begin{lemma}[Parameters of Reed--Solomon codes]\label{lem:rs-params}
Let $\Cc = \RS_\F[L,d]$ with $|L| = M$.
\begin{enumerate}
  \item The map $\ev_L : \F[X]_{<d} \to \Cc$ is an $\F$-linear bijection. For $c \in \Cc$ we write $P_c \in \F[X]_{<d}$ for the unique polynomial with $\ev_L(P_c) = c$.
Since $\ev_L$ is injective on $\F'[X]_{<M}$ for every finite field $\F' \supseteq \F$, $P_c$ depends only on the word $c$, and not on $d$ or on the field in which $c$ is regarded.
  \item Two distinct codewords agree on at most $d-1$ points of $L$, so $\Delta(c,c') \ge 1 - (d-1)/M > 1 - \rho$ for all distinct $c, c' \in \Cc$.
  \item If $2\delta < 1 - (d-1)/M$, in particular if $\delta \le (1-\rho)/2$, then every word $w \in \F^L$ is within distance $\delta$ of at most one codeword.
\end{enumerate}
\end{lemma}

\begin{proof}
Items 1 and 2 follow from \cref{lem:roots}, since $d \le M$ and the points of $L$ are distinct. For item 3, two codewords within distance $\delta$ of $w$ would be within distance $2\delta < 1 - (d-1)/M$ of each other (\cref{lem:metric}), contradicting item 2. Finally $(1-\rho)/2 = (1 - d/M)/2 < (1 - (d-1)/M)/2$.
\end{proof}

\begin{fact}[Efficient unique decoding~{\cite{WB86,Gao03}}]\label{fact:decode}
There is an algorithm that, given $L$, $d$ and a word $w \in \F^L$, outputs the unique $c \in \RS_\F[L,d]$ with $\Delta(w,c) < \tfrac12\bigl(1 - (d-1)/M\bigr)$ whenever it exists, using a number of operations in $\F$ polynomial in $M$.
Checking the distance of the output to $w$ then decides whether such a codeword exists, at the cost of $O(M)$ further comparisons.
\end{fact}

\begin{lemma}[Decoding over an extension field]\label{lem:descent}
Let $\K \supseteq \Fq$ be a finite field, $L$ a smooth domain of order $M$, $1 \le d \le M$, and $w \in \Fq^L$.
If $c \in \RS_\K[L,d]$ satisfies $\Delta(w,c) < \tfrac12\bigl(1 - (d-1)/M\bigr)$, then $c \in \Fq^L$ and $P_c \in \Fq[X]_{<d}$; in particular $c \in \RS[L,d]$.
\end{lemma}

\begin{proof}
Let $\mathrm{Frob} : \K \to \K$, $x \mapsto x^q$, be the Frobenius map; it is an injective ring homomorphism, since $q$ is a power of the characteristic $p$ and $(x+y)^p = x^p + y^p$.
Its fixed points are the roots of $X^q - X$, of which there are at most $q$ by \cref{lem:roots}; every element of $\Fq$ is one (since $x^{q-1} = 1$ for $x \in \Fq^\times$), so the fixed points are exactly the elements of $\Fq$.
Apply $\mathrm{Frob}$ entrywise to words, and coefficientwise to polynomials.
Since $L \subseteq \Fq$, we have $\mathrm{Frob}(P(\xi)) = \mathrm{Frob}(P)(\xi)$ for $\xi \in L$, so $\mathrm{Frob}(c) = \ev_L(\mathrm{Frob}(P_c)) \in \RS_\K[L,d]$.
Since $\mathrm{Frob}$ is injective and $\mathrm{Frob}(w) = w$, $\Delta(w, \mathrm{Frob}(c)) = \Delta(\mathrm{Frob}(w), \mathrm{Frob}(c)) = \Delta(w,c)$.
By \cref{lem:rs-params}(3), $\mathrm{Frob}(c) = c$, so $c \in \Fq^L$.
Then $\mathrm{Frob}(P_c)$ and $P_c$ both lie in $\K[X]_{<d}$ and evaluate to $c$, so they are equal by \cref{lem:rs-params}(1), and the coefficients of $P_c$ lie in $\Fq$.
\end{proof}

\paragraph{Even--odd split of words.}
Let $L$ be a smooth domain of order $M \ge 2$, and $\F \supseteq \Fq$ a finite field.
For a word $w \in \F^L$, define the words $w_e, w_o \in \F^{L^2}$ by
\begin{equation}\label{eq:even-odd}
  w_e(\xi^2) \;=\; \frac{w(\xi) + w(-\xi)}{2},
  \qquad
  w_o(\xi^2) \;=\; \frac{w(\xi) - w(-\xi)}{2\xi}
  \qquad (\xi \in L).
\end{equation}

\begin{lemma}[Even--odd split of words]\label{lem:split-words}
Let $w \in \F^L$ with $|L| = M \ge 2$.
\begin{enumerate}
  \item The words $w_e$ and $w_o$ are well defined, and $w \mapsto (w_e, w_o)$ is $\F$-linear.
  \item (Locality.) For $\eta \in L^2$, the values $w_e(\eta)$ and $w_o(\eta)$ depend only on the values of $w$ on the fibre over $\eta$.
  \item For every $\xi \in L$,
  \begin{equation}\label{eq:even-odd-inverse}
    w(\xi) \;=\; w_e(\xi^2) + \xi\, w_o(\xi^2), \qquad w(-\xi) \;=\; w_e(\xi^2) - \xi\, w_o(\xi^2).
  \end{equation}
  Consequently, $w$ vanishes on the fibre over $\eta$ if and only if $w_e(\eta) = w_o(\eta) = 0$.
  \item If $1 \le d \le M/2$, $P \in \F[X]_{<2d}$ and $w = \ev_L(P)$, then $w_e = \ev_{L^2}(P_e)$ and $w_o = \ev_{L^2}(P_o)$, where $(P_e, P_o)$ is the even--odd split of $P$ (\cref{lem:splits}); in particular $w_e, w_o \in \RS_\F[L^2, d]$.
\end{enumerate}
\end{lemma}

\begin{proof}
(1) The denominators are nonzero because the characteristic is odd and $\xi \ne 0$. Replacing $\xi$ by $-\xi$ leaves both right-hand sides of \eqref{eq:even-odd} unchanged, and $\xi$ and $-\xi$ are the only square roots of $\xi^2$ in $L$ (\cref{lem:power}(4)), so the words are well defined. Linearity is clear.
(2) is visible in \eqref{eq:even-odd}.
(3) Add and subtract the two formulas of \eqref{eq:even-odd}. If $w_e(\eta) = w_o(\eta) = 0$, then \eqref{eq:even-odd-inverse} gives $w(\pm\xi) = 0$; the converse follows from \eqref{eq:even-odd}.
(4) For $\xi \in L$, $P(\pm\xi) = P_e(\xi^2) \pm \xi P_o(\xi^2)$; substituting into \eqref{eq:even-odd} gives $w_e(\xi^2) = P_e(\xi^2)$ and $w_o(\xi^2) = P_o(\xi^2)$. Moreover $P_e, P_o \in \F[X]_{<d}$, and $d \le M/2 = |L^2|$ by \cref{lem:power}(2).
\end{proof}

\subsection{Proximity gaps}\label{sec:prelim:gaps}

The external result used in our soundness analysis is the correlated-agreement theorem of Ben-Sasson, Carmon, Ishai, Kopparty and Saraf in the unique-decoding regime.
In their notation, $\RS[\F, D, k]$ is the code of polynomials of degree at most $k$ evaluated on a set $D \subseteq \F$, with rate $(k+1)/|D|$.
The statement below is theirs with $D = L$ and $d = k+1$.
It holds over any finite field $\F$; we apply it with $\F \supseteq \Fq$ and $L \subseteq \Fq^\times \subseteq \F$.

\begin{theorem}[{\cite[Theorem~4.1]{BCIKS23}}]\label{thm:bciks}
Let $\F \supseteq \Fq$ be a finite field, $\Cc = \RS_\F[L,d]$ with $|L| = M$ and rate $\rho = d/M$, and $0 < \delta \le (1-\rho)/2$.
Let $u^0, u^1 \in \F^{L}$ and
\[
  S \;=\; \bigl\{\, r \in \F \;:\; \Delta(u^0 + r\, u^1,\, \Cc) \le \delta \,\bigr\}.
\]
If $|S| > M$, then there are a set $L' \subseteq L$ with $|L'| \ge (1-\delta)M$ and codewords $\hat c^0, \hat c^1 \in \Cc$ such that $u^0$ agrees with $\hat c^0$, and $u^1$ agrees with $\hat c^1$, on $L'$.
\end{theorem}

\section{Interactive Oracle Proofs and Polynomial Commitments}\label{sec:proofs}

This section recalls the proof-system notions used in \cref{sec:pcs,sec:sound}: interactive oracle proofs, their probability space and soundness notions, the sumcheck protocol, polynomial commitment schemes, and the compilation of interactive oracle proofs into non-interactive arguments.
Everything in \cref{sec:proofs:iop,sec:proofs:rbr,sec:proofs:sumcheck,sec:proofs:pcs} is proved.
\Cref{sec:proofs:bcs} concerns the compiled, non-interactive scheme; it restates the definitions of Chiesa, Di, Hu and Zheng~\cite{CDHZ25} that we need and quotes their theorem (\cref{thm:cdhz}), and nothing in \cref{sec:proofs:iop,sec:proofs:rbr,sec:proofs:sumcheck,sec:proofs:pcs} or in \cref{sec:kernel,sec:fold,sec:pcs} depends on it.

\subsection{Interactive oracle proofs}\label{sec:proofs:iop}

A \emph{relation} $\mathcal{R}$ is a set of pairs $(\mathtt{x},\mathtt{w})$ of an \emph{instance} $\mathtt{x}$ and a \emph{witness} $\mathtt{w}$; its \emph{language} is $L_{\mathcal{R}} = \{\mathtt{x} : \exists\, \mathtt{w},\ (\mathtt{x},\mathtt{w}) \in \mathcal{R}\}$.
An instance may have an \emph{oracle part}: a word to which the verifier has query access only.

\begin{definition}[Interactive oracle proof]\label{def:iop}
A \emph{public-coin interactive oracle proof} (IOP)~\cite{BCS16} with $\nu \ge 1$ rounds and \emph{challenge sets} $\mathsf{Ch}_1,\dots,\mathsf{Ch}_\nu$ (finite and nonempty) is a pair $(\mathsf{P},\mathsf{V})$ of a prover and a verifier.
On a common instance $\mathtt{x}$ (and a witness given to $\mathsf{P}$ only), the interaction proceeds as follows. For $i = 1, \dots, \nu$:
\begin{enumerate}
  \item $\mathsf{P}$ sends a message $\mathsf{m}_i$. Parts of $\mathsf{m}_i$ may be designated as \emph{oracles}: $\mathsf{V}$ does not read them, but may later query individual entries.
  \item $\mathsf{V}$ replies with a \emph{challenge} $\mathsf{c}_i \in \mathsf{Ch}_i$.
\end{enumerate}
After the last round, $\mathsf{V}$ reads the non-oracle parts of the instance and of the messages, queries entries of the oracles (including the oracle part of the instance) at positions that are determined by these non-oracle parts and by the challenges, and outputs \emph{accept} or \emph{reject} as a deterministic function of everything it has read.
\end{definition}

A \emph{partial transcript} of $i$ rounds is $\tau = (\mathtt{x}, \mathsf{m}_1, \mathsf{c}_1, \dots, \mathsf{m}_i, \mathsf{c}_i)$; for $i = 0$ it is the \emph{empty transcript} $(\mathtt{x})$, and for $i = \nu$ it is a \emph{full transcript}.
Since the verifier's decision is a deterministic function of the full transcript (with oracles), we say that $\mathsf{V}$ \emph{accepts} or \emph{rejects} a full transcript.
When the verifier makes checks ``during'' a round, as in the protocols below, these checks are predicates on the transcript that are evaluated at the end; the interaction always runs all $\nu$ rounds.

\paragraph{Provers and the probability space.}
A \emph{deterministic prover} $\mathsf{P}^*$ is a function that maps each partial transcript of $i-1$ rounds to a message $\mathsf{m}_i$.
Its interaction with $\mathsf{V}$ on an instance $\mathtt{x}$ is determined by the challenges.
All probabilities are over the \emph{challenge space} $\Omega = \mathsf{Ch}_1 \times \dots \times \mathsf{Ch}_\nu$ with the uniform distribution, that is, with independent uniform challenges.
We write $\tau_i$ for the partial transcript after $i$ rounds, which is a function of $(\mathsf{c}_1,\dots,\mathsf{c}_i)$, and $\mathsf{m}_i$ for the $i$-th message, which is a function of $(\mathsf{c}_1,\dots,\mathsf{c}_{i-1})$.
A \emph{randomised prover} is a finitely supported probability distribution over deterministic provers, independent of the challenges.

\begin{lemma}[Sequential challenges]\label{lem:sequential}
For $i \in \iv{1}{\nu}$, let $E_i \subseteq \mathsf{Ch}_1 \times \dots \times \mathsf{Ch}_i$, and suppose that for every $(\mathsf{c}_1,\dots,\mathsf{c}_{i-1})$ there are at most $k_i$ values $\mathsf{c}_i \in \mathsf{Ch}_i$ with $(\mathsf{c}_1,\dots,\mathsf{c}_i) \in E_i$. Then
\[
  \Pr_{\Omega}\bigl[\, \exists\, i \in \iv{1}{\nu} : (\mathsf{c}_1,\dots,\mathsf{c}_i) \in E_i \,\bigr] \;\le\; \sum_{i=1}^{\nu} \frac{k_i}{|\mathsf{Ch}_i|}.
\]
\end{lemma}

\begin{proof}
Under the uniform distribution on $\Omega$, the prefix $(\mathsf{c}_1,\dots,\mathsf{c}_{i})$ is uniform on $\mathsf{Ch}_1 \times \dots \times \mathsf{Ch}_i$.
Counting the elements of $E_i$ prefix by prefix gives $|E_i| \le k_i \prod_{i' < i} |\mathsf{Ch}_{i'}|$, so the probability of the $i$-th event is at most $k_i/|\mathsf{Ch}_i|$. Conclude by the union bound.
\end{proof}

\begin{lemma}[Independent queries]\label{lem:queries}
Let $\Omega'$ be a finite nonempty set with the uniform distribution, $L$ a finite nonempty set, $\kappa \ge 1$, and let $\varpi \mapsto \mathcal{S}(\varpi) \subseteq L$ and $\mathcal{A} \subseteq \Omega'$ be given.
On $\Omega' \times L^{\times\kappa}$ with the uniform distribution, where $L^{\times\kappa}$ denotes the Cartesian power (the set of $\kappa$-tuples of elements of $L$),
\[
  \Pr\bigl[\, \varpi \in \mathcal{A} \text{ and } \xi^{(1)},\dots,\xi^{(\kappa)} \in \mathcal{S}(\varpi) \,\bigr] \;=\; \mathbb{E}_{\varpi}\Bigl[\mathbf{1}_{\mathcal{A}}(\varpi) \cdot \Bigl(\frac{|\mathcal{S}(\varpi)|}{|L|}\Bigr)^{\kappa}\Bigr],
\]
where $\mathbf{1}_{\mathcal{A}}$ is the indicator function of $\mathcal{A}$. In particular, if $\pi \ge 0$ and $|\mathcal{S}(\varpi)| \le \pi |L|$ for all $\varpi \in \mathcal{A}$, the probability is at most $\Pr[\mathcal{A}]\,\pi^\kappa \le \pi^{\kappa}$.
\end{lemma}

\begin{proof}
For fixed $\varpi \in \mathcal{A}$, the number of $\kappa$-tuples of elements of $\mathcal{S}(\varpi)$ is $|\mathcal{S}(\varpi)|^\kappa$; divide by $|\Omega'|\,|L|^\kappa$ and sum over $\varpi$.
\end{proof}

\begin{lemma}[Randomised provers]\label{lem:randomised}
Fix an instance $\mathtt{x}$, and let $E$ be a set of full transcripts of $\mathtt{x}$ (for instance, the accepted ones). If $\Pr_\Omega[\tau_\nu \in E] \le \varepsilon$ for every deterministic prover, then the same bound holds for every randomised prover. Equivalently, if a randomised prover achieves $\Pr[\tau_\nu \in E] > \varepsilon$, so does some deterministic prover in its support.
\end{lemma}

\begin{proof}
Since the prover's randomness is independent of the challenges, the probability for a randomised prover is the average, over its distribution, of the probabilities for deterministic provers.
\end{proof}

\begin{definition}[Completeness and soundness]\label{def:iop-sound}
The IOP is \emph{complete} for $\mathcal{R}$ if, for every $(\mathtt{x},\mathtt{w}) \in \mathcal{R}$, the honest prover makes $\mathsf{V}$ accept with probability $1$.
It has \emph{soundness error} $\varepsilon$ if, for every $\mathtt{x} \notin L_{\mathcal{R}}$ and every deterministic prover $\mathsf{P}^*$, $\mathsf{V}$ accepts with probability at most $\varepsilon$.
\end{definition}

By \cref{lem:randomised}, soundness also holds against randomised provers.

\subsection{Round-by-round soundness}\label{sec:proofs:rbr}

For the Fiat--Shamir compilation, the relevant notions are round by round: every partial transcript is classified as \emph{doomed} or not, and no round can turn a doomed transcript into a non-doomed one except with small probability.
The following definitions are those of~\cite[Definitions~3.11 and~3.12]{BGTZ24}; the soundness notion was introduced in~\cite{CCHLRRW19}, where it is phrased with a \emph{state function}, which amounts to specifying a set $\Dd$ of doomed transcripts, and its knowledge analogue in~\cite{CMS19}.
We place the verifier's query positions in the last challenge $\mathsf{c}_\nu$.

\begin{definition}[Round-by-round soundness]\label{def:rbr}
An IOP with $\nu$ rounds has \emph{round-by-round soundness error} $(\varepsilon_1,\dots,\varepsilon_\nu) \in [0,1]^\nu$ for $\mathcal{R}$ if there is a set $\Dd$ of partial transcripts, the \emph{doomed} ones, such that:
\begin{enumerate}
  \item if $\mathtt{x} \notin L_{\mathcal{R}}$, then the empty transcript $(\mathtt{x})$ is doomed;
  \item for every $i \in \iv{1}{\nu}$, every doomed partial transcript $\tau$ of $i-1$ rounds and every message $\mathsf{m}_i$,
  \[
    \Pr_{\mathsf{c}_i \leftarrow \mathsf{Ch}_i}\bigl[(\tau, \mathsf{m}_i, \mathsf{c}_i) \notin \Dd\bigr] \;\le\; \varepsilon_i ;
  \]
  \item $\mathsf{V}$ rejects every doomed full transcript.
\end{enumerate}
We write $\varepsilon_{\mathrm{rbr}} = \max_i \varepsilon_i$.
\end{definition}

\begin{definition}[Round-by-round knowledge soundness]\label{def:rbr-knowledge}
An IOP with $\nu$ rounds is \emph{round-by-round knowledge sound} for $\mathcal{R}$ with error $(\varepsilon_1,\dots,\varepsilon_\nu) \in [0,1]^\nu$ if there are a set $\Dd$ of partial transcripts and an algorithm $\mathsf{Ext}$, the \emph{extractor}, running in time polynomial in the size of the instance (including its oracle part), defined on all pairs of a partial transcript and a next message, such that:
\begin{enumerate}
  \item the empty transcript $(\mathtt{x})$ is doomed for every instance $\mathtt{x}$;
  \item for every $i \in \iv{1}{\nu}$, every doomed partial transcript $\tau$ of $i-1$ rounds and every message $\mathsf{m}_i$, if
  \[
    \Pr_{\mathsf{c}_i \leftarrow \mathsf{Ch}_i}\bigl[(\tau, \mathsf{m}_i, \mathsf{c}_i) \notin \Dd\bigr] \;>\; \varepsilon_i ,
  \]
  then $(\mathtt{x}, \mathsf{Ext}(\tau, \mathsf{m}_i)) \in \mathcal{R}$, where $\mathtt{x}$ is the instance of $\tau$;
  \item $\mathsf{V}$ rejects every doomed full transcript.
\end{enumerate}
We write $\varepsilon_{\mathrm{rbr}} = \max_i \varepsilon_i$.
\end{definition}

The polynomial-time requirement on $\mathsf{Ext}$ matters only for the compiled scheme of \cref{sec:proofs:bcs}; the results of \cref{sec:proofs:rbr,sec:proofs:pcs} use conditions 1--3 but not the running time of $\mathsf{Ext}$.

\begin{lemma}[Round-by-round implies overall]\label{lem:rbr-overall}
Let $\mathsf{P}^*$ be a deterministic prover and $\mathtt{x}$ an instance, and let $\tau_i$ and $\mathsf{m}_i$ be as in \cref{sec:proofs:iop}.
\begin{enumerate}
  \item If the IOP has round-by-round soundness error $(\varepsilon_i)_i$ and $\mathtt{x} \notin L_{\mathcal{R}}$, then $\mathsf{V}$ accepts with probability at most $\sum_{i=1}^{\nu} \varepsilon_i$.
  \item If the IOP is round-by-round knowledge sound with error $(\varepsilon_i)_i$, then
  \[
    \Pr[\mathsf{V} \text{ accepts}] \;\le\; \sum_{i=1}^{\nu} \varepsilon_i \;+\; \Pr\bigl[\, \exists\, i \in \iv{1}{\nu} : \tau_{i-1} \in \Dd \text{ and } (\mathtt{x}, \mathsf{Ext}(\tau_{i-1}, \mathsf{m}_i)) \in \mathcal{R} \,\bigr].
  \]
  In particular, if $\mathsf{V}$ accepts with probability greater than $\sum_i \varepsilon_i$, then there exist $i \in \iv{1}{\nu}$ and challenges $\mathsf{c}_1,\dots,\mathsf{c}_{i-1}$ such that $\tau_{i-1} \in \Dd$ and $(\mathtt{x}, \mathsf{Ext}(\tau_{i-1}, \mathsf{m}_i)) \in \mathcal{R}$.
\end{enumerate}
\end{lemma}

\begin{proof}
In both cases $\tau_0$ is doomed, and an accepted full transcript $\tau_\nu$ is not doomed (condition 3).
So if $\mathsf{V}$ accepts, there is a first round $i$ with $\tau_{i-1} \in \Dd$ and $\tau_i \notin \Dd$.
For (2), let $E_i$ be the set of $(\mathsf{c}_1,\dots,\mathsf{c}_i)$ such that $\tau_{i-1} \in \Dd$, $(\mathtt{x}, \mathsf{Ext}(\tau_{i-1}, \mathsf{m}_i)) \notin \mathcal{R}$ and $\tau_i \notin \Dd$.
For fixed $(\mathsf{c}_1,\dots,\mathsf{c}_{i-1})$, the transcript $\tau_{i-1}$ and the message $\mathsf{m}_i$ are fixed; if the first two conditions hold, condition~2 of \cref{def:rbr-knowledge} (in contrapositive form) shows that at most $\varepsilon_i |\mathsf{Ch}_i|$ values of $\mathsf{c}_i$ satisfy the third.
Acceptance implies that some $E_i$ occurs or that the event in the displayed probability occurs, and \cref{lem:sequential} bounds the probability of the former by $\sum_i \varepsilon_i$.
For (1), the same argument applies with $E_i$ defined without the condition on $\mathsf{Ext}$, using condition~2 of \cref{def:rbr}.
\end{proof}

\subsection{The sumcheck protocol}\label{sec:proofs:sumcheck}

The sumcheck protocol of Lund, Fortnow, Karloff and Nisan~\cite{LFKN92} reduces a claim about the sum of a polynomial over the Boolean hypercube to a claim about one evaluation.
Our commitment scheme runs it on a product of two multilinear polynomials, so we state it for polynomials of degree at most $2$ in each variable.

Let $\F \supseteq \Fq$ be a finite field, and let $F \in \F[x_1,\dots,x_n]$ have degree at most $2$ in each variable, with $n \ge 1$.
For $j \in \iv{1}{n}$ and $r_1,\dots,r_{j-1} \in \F$, the \emph{honest round polynomial} is
\[
  s^*_j(X) \;=\; \sum_{b = 0}^{2^{n-j}-1} F(r_1,\dots,r_{j-1}, X, b_1,\dots,b_{n-j}) \;\in\; \F[X]_{<3}.
\]
To prove the claim $\sum_{b=0}^{N-1} F(b) = v_0$, the prover sends, in round $j = 1,\dots,n$, a round polynomial $s_j \in \F[X]_{<3}$ (honestly $s_j = s^*_j$).
The verifier samples $r_j \in \F$ uniformly at random and sets $v_j = s_j(r_j)$; it accepts if $s_j(0) + s_j(1) = v_{j-1}$ for every $j \in \iv{1}{n}$ and $v_n = F(r_1,\dots,r_n)$.

\begin{lemma}[Sumcheck]\label{lem:sumcheck}
If $\sum_{b=0}^{N-1} F(b) = v_0$, the honest prover is accepted with probability $1$.
If $\sum_{b} F(b) \ne v_0$, then for every deterministic prover the verifier accepts with probability at most $2n/|\F|$.
More precisely, in that case the verifier accepts only if there is a round $j$ with $s_j \ne s^*_j$ and $s_j(r_j) = s^*_j(r_j)$.
\end{lemma}

\begin{proof}
Put $v^*_0 = \sum_b F(b)$ and $v^*_j = s^*_j(r_j)$ for $j \in \iv{1}{n}$.
By definition, $s^*_j(0) + s^*_j(1) = v^*_{j-1}$ for every $j$, and $v^*_n = F(r_1,\dots,r_n)$.
Completeness follows, since the honest prover has $v_j = v^*_j$ for all $j$.

Soundness: suppose that $v_0 \ne v^*_0$, that the verifier accepts, and that $s_j(r_j) \ne s^*_j(r_j)$ whenever $s_j \ne s^*_j$.
We show by induction that $v_j \ne v^*_j$ for all $j \in \iv{0}{n}$; for $j = n$ this contradicts the final check.
If $v_{j-1} \ne v^*_{j-1}$, the check of round $j$ gives $s_j(0) + s_j(1) = v_{j-1} \ne s^*_j(0) + s^*_j(1)$, so $s_j \ne s^*_j$, hence $v_j = s_j(r_j) \ne s^*_j(r_j) = v^*_j$ by assumption.

For the probability, let $E_j$ be the set of $(r_1,\dots,r_j)$ with $s_j \ne s^*_j$ and $s_j(r_j) = s^*_j(r_j)$.
For fixed $r_1,\dots,r_{j-1}$, both $s_j$ and $s^*_j$ are fixed, and their difference, when nonzero, has degree at most $2$, so at most $2$ values of $r_j$ lie in $E_j$ (\cref{lem:roots}).
\Cref{lem:sequential} gives the bound $2n/|\F|$.
\end{proof}

\subsection{Polynomial commitment schemes}\label{sec:proofs:pcs}

A \emph{multilinear polynomial commitment scheme} lets a prover commit to a table $f : \iv{0}{N-1} \to \F$, over a finite field $\F \supseteq \Fq$, and later prove claims of the form $\tilde f(z) = v$.
We describe it in the IOP model, where the commitment is an oracle.

\begin{definition}[Polynomial commitment in the IOP model]\label{def:pcs}
A multilinear polynomial commitment scheme in the IOP model consists of:
\begin{itemize}
  \item an \emph{encoding} $\Enc$ that maps each table $f$ to a word $\Enc(f)$, the \emph{commitment};
  \item an \emph{evaluation protocol}: a public-coin IOP whose instance is $(w; z, v)$, with a word $w$ (the commitment) as oracle part, a point $z \in \F^n$ and a value $v \in \F$.
\end{itemize}
It is \emph{complete} if the evaluation protocol accepts with probability $1$ when the prover is honest, $w = \Enc(f)$ and $\tilde f(z) = v$.
\end{definition}

Every claim $(z,v)$ is true for \emph{some} table (for instance the constant table with value $v$, since $\sum_b \eq(b,z) = \prod_{k=1}^n (z_k + 1 - z_k) = 1$), so the relation $\{((z,v), f) : \tilde f(z) = v\}$ says nothing about the commitment.
The meaningful requirement is that the commitment determines the table.
To express it, fix a function $\mathrm{dist}$ from pairs of words to real numbers, and a threshold $\delta$, such that
\begin{equation}\label{eq:dist-unique}
  \text{for every word } w, \text{ there is at most one table } f \text{ with } \mathrm{dist}(w, \Enc(f)) \le \delta,
\end{equation}
and consider the relation
\[
  \mathcal{R}^{\delta}_{\mathrm{eval}} \;=\; \bigl\{\, ((w; z,v), f) \;:\; \mathrm{dist}(w, \Enc(f)) \le \delta \ \text{ and } \ \tilde f(z) = v \,\bigr\}.
\]

\begin{definition}[Evaluation binding]\label{def:binding}
The evaluation protocol is \emph{$\varepsilon$-evaluation binding} if for every word $w$, point $z$ and values $v \ne v'$, it is not the case that some deterministic prover makes the verifier accept $(w; z, v)$ with probability greater than $\varepsilon$ and some deterministic prover makes it accept $(w; z, v')$ with probability greater than $\varepsilon$.
\end{definition}

By \cref{lem:randomised}, $\varepsilon$-evaluation binding also holds when randomised provers are allowed.

\begin{lemma}[Knowledge soundness implies binding]\label{lem:rbr-binding}
If the evaluation protocol is round-by-round knowledge sound for $\mathcal{R}^{\delta}_{\mathrm{eval}}$ with error $(\varepsilon_i)_i$, and \eqref{eq:dist-unique} holds, then it is $\varepsilon$-evaluation binding with $\varepsilon = \sum_i \varepsilon_i$.
\end{lemma}

\begin{proof}
Suppose that provers are accepted on $(w; z,v)$ and on $(w;z,v')$, each with probability greater than $\varepsilon$.
By \cref{lem:rbr-overall}(2), the extractor outputs, on some partial transcript of each interaction, tables $f$ and $f'$ with $((w;z,v),f) \in \mathcal{R}^{\delta}_{\mathrm{eval}}$ and $((w;z,v'),f') \in \mathcal{R}^{\delta}_{\mathrm{eval}}$.
Both tables satisfy $\mathrm{dist}(w, \Enc(\cdot)) \le \delta$, so $f = f'$ by \eqref{eq:dist-unique}, and $v = \tilde f(z) = \tilde f'(z) = v'$.
\end{proof}

\subsection{Non-interactive arguments}\label{sec:proofs:bcs}

This subsection concerns only the non-interactive scheme (\cref{cor:pq}).

\paragraph{Merkle trees.}
Let $\mathsf{H} : \{0,1\}^* \to \{0,1\}^{\lambda_{\mathsf{H}}}$ be a hash function with output length $\lambda_{\mathsf{H}}$, and let $\|$ denote concatenation of bit strings.
The \emph{Merkle tree}~\cite{Mer89} over a vector $\mathbf{d} = (\mathbf{d}_0,\dots,\mathbf{d}_{T-1})$ of bit strings, with $T$ a power of $2$, is the complete binary tree whose leaves are the hashes $\mathsf{H}(0 \,\|\, \mathbf{d}_i)$ and whose internal nodes are $\mathsf{H}(1 \,\|\, \text{left child} \,\|\, \text{right child})$; its \emph{root} $\mathrm{MT}(\mathbf{d})$ is a commitment to the vector.
An \emph{opening} of position $i$ consists of $\mathbf{d}_i$ and its \emph{authentication path}, the $\log_2 T$ siblings on the path from the leaf to the root; it is checked by recomputing the root.
The entries $\mathbf{d}_i$ may be encodings of several field elements; in \cref{sec:impl} each leaf holds the two values of a word on one fibre.

\begin{lemma}[Merkle binding]\label{lem:merkle}
\begin{enumerate}
  \item Given two accepting openings of the same position to different values under the same root, one can compute a collision of $\mathsf{H}$, that is, two distinct inputs with the same hash, with $O(\log_2 T)$ hash evaluations.
  \item Given two vectors $\mathbf{d} \ne \mathbf{d}'$ of length $T$ with $\mathrm{MT}(\mathbf{d}) = \mathrm{MT}(\mathbf{d}')$, one can compute a collision of $\mathsf{H}$ with $O(T)$ hash evaluations.
\end{enumerate}
\end{lemma}

\begin{proof}
(1) Number the levels of the tree from $0$ (leaf hashes) to $\log_2 T$ (root). Node values are strings of the fixed length $\lambda_{\mathsf{H}}$, and the position $i$ fixes, at every level, on which side the path node lies; so the input $1\|\text{left}\|\text{right}$ hashed at a level determines the path node below it and its sibling.
Recompute both paths from the leaves to the root, and let $t$ be the first level at which the two recomputed path nodes are equal; it exists because both paths end at the same root.
If $t = 0$, the hashed inputs $0\|\mathbf{d}_i$ and $0\|\mathbf{d}'_i$ are distinct. If $t > 0$, the inputs hashed at level $t$ are distinct, because they contain the path nodes of level $t-1$, which differ by the choice of $t$. In both cases these two inputs have the same hash.
(2) Compute both trees, and apply (1) at a position where $\mathbf{d}$ and $\mathbf{d}'$ differ, taking the authentication paths from the two trees.
\end{proof}

\paragraph{Random oracles.}
In the \emph{random oracle model}, $\mathsf{H}$ is modelled as a uniformly random function to which all parties have oracle access; in the \emph{quantum random oracle model} (QROM)~\cite{BDFLSZ11}, the adversary is a quantum algorithm that may query $\mathsf{H}$ in superposition.
The \emph{BCS transformation} of Ben-Sasson, Chiesa and Spooner~\cite{BCS16} compiles a public-coin IOP into a non-interactive argument:
each oracle message is replaced by a Merkle root, each challenge is derived by hashing the instance and the transcript so far (the Fiat--Shamir paradigm), and the proof contains the non-oracle messages, the roots, and an opening of every queried oracle entry.
The BCS theorem in the QROM of Chiesa, Manohar and Spooner~\cite[Theorem~8.6]{CMS19} covers IOPs whose instance has no oracle part.
An evaluation protocol has an oracle part, the commitment, which in the non-interactive scheme is replaced by its Merkle root and is chosen by the adversary.
This setting is covered by the extension of the BCS transformation to \emph{interactive oracle reductions} (IORs) of Chiesa, Di, Hu and Zheng~\cite{CDHZ25}, which we now recall in the special case that we need.
All definitions below are those of~\cite{CDHZ25}; we restate them to fix notation, and we indicate the specialisation.
We write $\varepsilon$ for the error terms that~\cite{CDHZ25} denotes by $\kappa$, since $\kappa$ is our number of queries.

\paragraph{Relations with an implicit instance~\cite[\S3.3]{CDHZ25}.}
Fix a finite alphabet $\Sigma$ and a length $l_y \ge 1$.
A \emph{partial string} is an element of $(\Sigma \sqcup \{\bot\})^{l_y}$, the symbol $\bot$ marking an undefined entry; $\mathrm{Dom}(y)$ is the set of positions where $y$ is defined, and $y$ is \emph{full} if $\mathrm{Dom}(y)$ is everything.
The fractional Hamming distance $\Delta(y,y')$ of two partial strings is the fraction of positions at which they differ, $\bot$ being different from every symbol of $\Sigma$.
A \emph{relation with implicit instance} is a set $\mathcal{R}$ of triples $(\mathtt{x}, y, \mathtt{w})$ of an \emph{explicit instance} $\mathtt{x}$, a full string $y \in \Sigma^{l_y}$, the \emph{implicit instance}, to which the verifier has oracle access only, and a witness $\mathtt{w}$.
For $\delta \in [0,1]$, the \emph{relaxed relation} is
\[
  \mathcal{R}_{\le\delta} \;=\; \bigl\{\, (\mathtt{x}, y, \mathtt{w}) \;:\; y \in (\Sigma \sqcup \{\bot\})^{l_y},\ \exists\, y^* \in \Sigma^{l_y} \text{ with } (\mathtt{x}, y^*, \mathtt{w}) \in \mathcal{R} \text{ and } \Delta(y, y^*) \le \delta \,\bigr\}.
\]

\paragraph{IOPs as interactive oracle reductions~\cite[\S3.4]{CDHZ25}.}
Consider an IOP (\cref{def:iop}) with $\nu$ rounds whose instance is $(\mathtt{x}; y)$ with oracle part $y$, in which every prover message is a string $\Pi_i$ over $\Sigma$, whose challenges are uniformly random bit strings $\mathsf{vr}_i \in \{0,1\}^{\beta_i}$, and whose verifier reads (by queries) positions of $y$ and of the $\Pi_i$ that are determined by $\mathtt{x}$ and the challenges only; a message part that the verifier reads entirely is a string all of whose positions are read.
Let $\mathcal{R}_{\mathrm{acc}} = \{(\top, (), \mathtt{w})\}$, where $\top$ is a fixed symbol, $()$ is the empty implicit instance, and $\mathtt{w}$ is arbitrary; then $(\mathcal{R}_{\mathrm{acc}})_{\le\delta} = \mathcal{R}_{\mathrm{acc}}$ for every $\delta$.
If the verifier outputs the instance $(\top, ())$ when it accepts and rejects otherwise, the IOP is an IOR from $\mathcal{R}$ to $\mathcal{R}_{\mathrm{acc}}$ in the sense of~\cite[\S3.4]{CDHZ25}, with $n_y = 1$ input implicit instance and $n'_y = 0$ output implicit instances.
$\mathcal{R}_{\mathrm{acc}}$ is \emph{monotone} in the sense of~\cite[Definition~11.1]{CDHZ25}, because it has no implicit instance.
A \emph{transcript} is a sequence $\mathrm{tr} = (\Pi_1, \mathsf{vr}_1, \Pi_2, \dots)$; the empty transcript is written $\mathrm{tr}_\emptyset$.

In the proof of the theorem below, the verifier and the knowledge state function are evaluated on strings produced by an extractor, whose entries may be undefined~\cite[Constructions~11.8 and~11.9]{CDHZ25}.
We therefore require the verifier to be defined when $y$ and the $\Pi_i$ are partial strings, and we take the statements $(\mathtt{x}, y)$ and the messages $\Pi_i$ in \cref{def:rrbr} to range over partial strings.
For the protocols of this paper, the verifier rejects as soon as one of the entries it reads is $\bot$.

\begin{definition}[Relaxed round-by-round knowledge soundness~\cite{CDHZ25}]\label{def:rrbr}
This is~\cite[Definitions~3.5 and~3.6]{CDHZ25} with output relation $\mathcal{R}_{\mathrm{acc}}$.
A \emph{knowledge state function} for an IOR from $\mathcal{R}$ to $\mathcal{R}_{\mathrm{acc}}$ is a deterministic function $\mathsf{KState}$ that, on input a proximity bound $\delta \in [0,1]$, a statement $(\mathtt{x}, y)$, a transcript $\mathrm{tr}$ and a witness $\mathtt{w}$, outputs a bit $\mathsf{KState}_\delta(\mathtt{x}, y, \mathrm{tr}, \mathtt{w})$, such that:
\begin{enumerate}
  \item (empty transcript) $\mathsf{KState}_\delta(\mathtt{x}, y, \mathrm{tr}_\emptyset, \mathtt{w}) = 1$ if and only if $(\mathtt{x}, y, \mathtt{w}) \in \mathcal{R}_{\le\delta}$;
  \item (prover moves) if the prover is about to move after $\mathrm{tr}$ and $\mathsf{KState}_\delta(\mathtt{x}, y, \mathrm{tr}, \mathtt{w}) = 0$, then $\mathsf{KState}_\delta(\mathtt{x}, y, \mathrm{tr} \| \Pi, \mathtt{w}) = 0$ for every message $\Pi$;
  \item (full transcript) for a full transcript, $\mathsf{KState}_\delta(\mathtt{x}, y, \mathrm{tr}, \mathtt{w}) = 1$ if and only if the verifier accepts $\mathrm{tr}$.
\end{enumerate}
The IOR has \emph{relaxed round-by-round knowledge soundness errors} $(\varepsilon^{\mathrm{rr}}_1, \dots, \varepsilon^{\mathrm{rr}}_\nu)$, functions of $(\mathtt{x}, y, \delta)$, if there are a knowledge state function $\mathsf{KState}$ and a deterministic polynomial-time algorithm $\mathsf{E}$, which does not depend on $\delta$, such that for every $\delta \in [0,1]$, every statement $(\mathtt{x}, y)$, every $i \in \iv{1}{\nu}$ and every transcript $\mathrm{tr} = (\Pi_1, \mathsf{vr}_1, \dots, \Pi_{i-1}, \mathsf{vr}_{i-1}, \Pi_i)$,
\begin{multline*}
  \Pr_{\mathsf{vr}_i \leftarrow \{0,1\}^{\beta_i}}\Bigl[\, \exists\, \mathtt{w} :\ \mathsf{KState}_\delta\bigl(\mathtt{x}, y, \mathrm{tr}, \mathsf{E}(\mathtt{x}, y, \mathrm{tr}\|\mathsf{vr}_i, \mathtt{w})\bigr) = 0 \\
  \text{ and } \ \mathsf{KState}_\delta(\mathtt{x}, y, \mathrm{tr}\|\mathsf{vr}_i, \mathtt{w}) = 1 \,\Bigr] \;\le\; \varepsilon^{\mathrm{rr}}_i(\mathtt{x}, y, \delta).
\end{multline*}
\end{definition}

For $\mathcal{R}_{\mathrm{acc}}$, condition 3 is the specialisation of~\cite[Definition~3.5]{CDHZ25}, which asks that the verifier output an instance $(\mathtt{x}', y')$ with $(\mathtt{x}', y', \mathtt{w}) \in (\mathcal{R}_{\mathrm{acc}})_{\le\delta}$, that is, that it accept.

\paragraph{Salted Merkle commitments~\cite[Definition~8.2]{CDHZ25}.}
The compiler of~\cite{CDHZ25} uses a salted variant of the Merkle tree: to commit to $\mathbf{a} = (a_1,\dots,a_l) \in \Sigma^{l}$ ($l$ a power of $2$, after padding), it samples uniformly random salts $\gamma_1,\dots,\gamma_l$ of a fixed length, sets the leaves to $\mathsf{H}(a_h, \gamma_h)$ and the internal nodes to $\mathsf{H}(\text{left child}, \text{right child})$, where $\mathsf{H}$ is a random oracle with output length $\lambda_{\mathsf{H}}$; the root is the commitment $\mathrm{cm}$, and the salts and the tree form the trapdoor $\mathrm{td}$.
An opening of a set $I$ of positions consists of $(a_h, \gamma_h)$ and the authentication path of every $h \in I$, and is checked by recomputing the root.
The proof of \cref{lem:merkle}(1) applies verbatim, with the input $(a_h, \gamma_h)$ in place of $0 \| \mathbf{d}_i$ and $\text{left}\|\text{right}$ in place of $1\|\text{left}\|\text{right}$: \emph{two accepting openings of the same position to different values under the same root and the same oracle yield a collision of that oracle.}

\begin{definition}[Committed relations~{\cite[Definition~11.2]{CDHZ25}}]\label{def:committed}
Let $\mathcal{R}$ be a relation with implicit instance and $\delta \in [0,1]$.
A \emph{committed instance} is $\mathtt{cx} = (\mathtt{x}, \mathrm{cm})$, where $\mathrm{cm}$ is a Merkle root, and a \emph{committed witness} is $\mathtt{cw} = (y, \mathtt{w}, \mathrm{td})$.
\begin{itemize}
  \item $(\mathtt{cx}, \mathtt{cw}) \in \mathsf{CR}[\mathcal{R}]$ if $(\mathtt{x}, y, \mathtt{w}) \in \mathcal{R}$ and $\mathrm{cm}$ and $\mathrm{td}$ are the root and the trapdoor of the salted Merkle tree of $y$ with the salts recorded in $\mathrm{td}$;
  \item $(\mathtt{cx}, \mathtt{cw}) \in \widehat{\mathsf{CR}}[\mathcal{R}, \delta]$ (the \emph{relaxed committed relation}) if $y$ is a partial string with $(\mathtt{x}, y, \mathtt{w}) \in \mathcal{R}_{\le\delta}$, and the opening of all positions of $\mathrm{Dom}(y)$ computed from $\mathrm{td}$ is accepted for $y$ under the root $\mathrm{cm}$.
\end{itemize}
\end{definition}

In~\cite{CDHZ25} the committed instance also carries an index that selects, by domain separation, the random oracle used by the commitment. We fix this index in the scheme (the verifier rejects any other value) and omit it from the notation. This extra check only removes accepted proofs and makes no oracle query, so the extraction event of~\cite[Definition~10.3]{CDHZ25} (acceptance and failed extraction) can only become smaller, and \cref{thm:cdhz} still applies.

\paragraph{The compiler.}
$\mathsf{BCS}[\Pi]$ denotes the non-interactive reduction of~\cite[Construction~11.7]{CDHZ25} applied to an IOR $\Pi$, with salted Merkle commitments of output length $\lambda_{\mathsf{H}}$.
Its prover receives $(\mathtt{cx}, \mathtt{cw}) \in \mathsf{CR}[\mathcal{R}]$; in round $i$ it commits to $\Pi_i$ by a salted Merkle tree, samples a fresh random salt, and derives $\mathsf{vr}_i$ by querying a random oracle on $\mathtt{x}$, the root $\mathrm{cm}$ of $y$, the roots of $\Pi_1,\dots,\Pi_i$ and the salts of rounds $1,\dots,i$.
The proof consists of the roots, the salts, and the openings of every position of $y$ and of the $\Pi_i$ that the IOR verifier reads; the verifier recomputes the challenges, checks the openings, and runs the IOR verifier on the opened values.
All random oracles are derived from one random oracle by domain separation.

\paragraph{Knowledge soundness.}
The security notion is post-quantum straightline knowledge soundness~\cite[Definition~10.3]{CDHZ25}.
It is stated for an \emph{emulated oracle}, which answers the adversary's queries and gives the extractor an additional extraction interface.
It has three error terms: a simulation error $\varepsilon_{\mathrm{sim}}$, bounding the distinguishing advantage between the emulated oracle and a true random oracle; a commutativity error $\varepsilon_{\mathrm{com}}$; and the extraction error $\varepsilon_{\mathrm{ext}}$, which bounds, for a $Q$-query quantum adversary $\tilde{\mathsf{P}}$ that outputs a committed instance $\mathtt{cx}$ of its choice, a proof, and a witness for the output relation, the probability that
\begin{center}
the verifier accepts, \quad and \quad the output $\mathtt{cw}$ of the extractor satisfies $(\mathtt{cx}, \mathtt{cw}) \notin \widehat{\mathsf{CR}}[\mathcal{R}, \delta]$.
\end{center}
(For $\mathcal{R}_{\mathrm{acc}}$ the output relation holds exactly when the verifier accepts, and its decider makes no oracle queries.)
The extractor is a polynomial-time quantum algorithm and is \emph{straightline}: it does not rewind the adversary.

\begin{theorem}[{\cite[Theorems~6.10, 8.3 and~11.3]{CDHZ25}}, specialised]\label{thm:cdhz}
Let $\Pi$ be an IOR from $\mathcal{R}$ to $\mathcal{R}_{\mathrm{acc}}$ as above, with $\nu$ rounds, challenge lengths $\beta_1,\dots,\beta_\nu$ and $\beta_{\min} = \min_i \beta_i$, relaxed round-by-round knowledge soundness errors $(\varepsilon^{\mathrm{rr}}_i)_i$ (\cref{def:rrbr}), and let $\varepsilon_{\mathrm{rr}}(\delta) = \max_{i, (\mathtt{x}, y)} \varepsilon^{\mathrm{rr}}_i(\mathtt{x}, y, \delta)$, the maximum being over the statements of bounded size under consideration.
Let $l_{\max}$ be the largest length of $y$ and of the strings $\Pi_i$, $q_{\mathsf{V}}$ the number of random-oracle queries of the verifier of $\mathsf{BCS}[\Pi]$, $q_{\widehat{\mathsf{CR}}}$ that of the decider of $\widehat{\mathsf{CR}}[\mathcal{R},\delta]$, $q_s$ the number of positions of $y$ and of the $\Pi_i$ that the IOR verifier reads, and $T_1 = Q + 1 + q_{\mathsf{V}} + \nu + q_{\widehat{\mathsf{CR}}}$.
Then there is a polynomial-time quantum extractor such that, for every $\delta \in [0,1]$ and every query bound $Q$, $\mathsf{BCS}[\Pi]$ is a non-interactive reduction from $(\mathsf{CR}[\mathcal{R}], \widehat{\mathsf{CR}}[\mathcal{R}, \delta])$ to $(\mathsf{CR}[\mathcal{R}_{\mathrm{acc}}], \widehat{\mathsf{CR}}[\mathcal{R}_{\mathrm{acc}}, \delta])$ with post-quantum straightline knowledge soundness errors $(0, \varepsilon_{\mathrm{com}}, \varepsilon_{\mathrm{ext}})$, where $\varepsilon_{\mathrm{com}}(q) \le 240\,(q+1)/2^{\lambda_{\mathsf{H}}}$ and
\begin{align*}
  \varepsilon_{\mathrm{ext}}(Q) \;\le\;\;& 320\,(Q+\nu+1)(Q+\nu)\,\varepsilon_{\mathrm{rr}}(\delta) \;+\; \frac{4\,(Q+\nu+1)\,\nu}{2^{\beta_{\min}}} \\
  &+\; \frac{32\, l_{\max} + 1280\, Q\,(2Q+1)^2 + 128\, q_s \lceil \log_2 l_{\max} \rceil}{2^{\lambda_{\mathsf{H}}}} \\
  &+\; \frac{960\,(Q+1)\,T_1^2 + 480\, q_{\mathsf{V}}^2\,(Q + \nu + q_{\mathsf{V}} + 2)}{2^{\lambda_{\mathsf{H}}}}.
\end{align*}
\end{theorem}

\begin{proof}[Derivation from~\cite{CDHZ25}]
By~\cite[Theorem~6.10]{CDHZ25}, $\Pi$ has post-quantum state-restoration straightline knowledge soundness error $\varepsilon^{\mathrm{sr}} \le 80\,(Q+\nu+1)(\mathfrak{m}_{\mathsf{FS}}+\nu)\,\varepsilon_{\mathrm{rr}}(\delta) + (Q+\nu+1)\,\nu/2^{\beta_{\min}}$ for a $Q$-query prover with total query mass $\mathfrak{m}_{\mathsf{FS}}$ to the challenge oracle.
By~\cite[Theorem~8.3]{CDHZ25}, the salted Merkle commitments are post-quantum multi-extractable with errors, at query bound $Q'$,
\begin{gather*}
  \varepsilon_{\mathrm{com}}(Q') \le \frac{240(Q'+1)}{2^{\lambda_{\mathsf{H}}}}, \qquad \varepsilon_{\mathrm{extract}} \le \frac{8\,n_y\,l_{\max}}{2^{\lambda_{\mathsf{H}}}}, \\
  \varepsilon_{\mathrm{offline}} \le \frac{160\, \mathfrak{m}\, Q'^2 + 16\, q_s \log_2 l_{\max}}{2^{\lambda_{\mathsf{H}}}}, \qquad \varepsilon_{\mathrm{online}} \le \frac{240\, \mathfrak{m}'\, Q'^2}{2^{\lambda_{\mathsf{H}}}},
\end{gather*}
where $\mathfrak{m}, \mathfrak{m}'$ are the total query masses in the third argument of these functions in~\cite{CDHZ25}.
Substitute these bounds into~\cite[Theorem~11.3]{CDHZ25} with $n_y = 1$, $n'_y = 0$ and $q_{\widehat{\mathsf{CR}}'} = 0$.
Every total query mass is at most the number of queries~\cite[Definition~4.3 and Remark~4.2]{CDHZ25}, so the masses $\mathfrak{m}_{\mathsf{FS}}, \mathfrak{m}_{\mathsf{VC}}, \mathfrak{m}_{\mathsf{MultiExtract}}$ over which the maximum in~\cite[Theorem~11.3]{CDHZ25} is taken are at most $Q$.
This gives the displayed terms in order: $4\varepsilon^{\mathrm{sr}}$; $4\varepsilon_{\mathrm{extract}}$ with $n_y = 1$; $8\varepsilon_{\mathrm{offline}}$ at query bound $2Q+1$ with mass $\mathfrak{m}_{\mathsf{VC}} \le Q$ (by~\cite[Remark~11.5]{CDHZ25} this term vanishes when $n'_y = 0$, and we keep it as an upper bound); $4\varepsilon_{\mathrm{online}}$ at query bound $T_1$ with mass $\mathfrak{m}_{\mathsf{MultiExtract}} + 1 \le Q+1$; and $2 q_{\mathsf{V}}^2 \varepsilon_{\mathrm{com}}$ at query bound $Q+1+\nu+q_{\mathsf{V}}$.
\end{proof}

We use the constants of~\cite[Theorem~8.3]{CDHZ25} as stated there (version of 2 March 2026); they follow from~\cite[Theorem~8.4, Lemmas~8.7 and~8.8]{CDHZ25} with the query counts of~\cite[\S8.3]{CDHZ25}.
A classical adversary is a special case of a quantum one, so \cref{thm:cdhz} also applies to classical adversaries in the random oracle model, with the same quantum extractor.
For classical adversaries, B\"unz, Mishra, Nguyen and Wang~\cite[Theorem~5.9]{BMNW24} give an asymptotic statement of the same shape for their notion of round-by-round knowledge~\cite[Definition~5.5]{BMNW24}; we do not use it.

\section{The Boolean-Kernel Basis}\label{sec:kernel}

This section introduces the kernel polynomials and shows that their Boolean members form a basis of $\F[X]_{<N}$, with a Kronecker-structured change of basis.
It then characterises low degree in the new coordinates and expresses univariate evaluation as a scaled multilinear evaluation.
Throughout, $\F \supseteq \Fq$ is a finite field (in the applications, $\F = \Fq$ or an extension field $\K$), and $n \ge 0$.
When the number of variables matters we write $K^{(n)}_y$, and otherwise $K_y$.

\subsection{Kernel polynomials}\label{sec:kernel:def}

\begin{definition}[Kernel polynomials]\label{def:kernel}
For $y = (y_1,\dots,y_n) \in \F^n$, the \emph{kernel polynomial} of $y$ is
\[
  K_y(X) \;=\; \prod_{k=1}^{n} \bigl( X^{2^{k-1}} + y_k \bigr) \;\in\; \F[X],
\]
with the empty product equal to $1$.
For $b \in \iv{0}{N-1}$ we write $K_b = K_{(b_1,\dots,b_n)}$, where $(b_1,\dots,b_n)$ is the bit vector of $b$.
The family $(K_b)_{b \in \iv{0}{N-1}}$ is the \emph{Boolean kernel family}; its members have coefficients in $\{0,1\}$, so they do not depend on $\F$.
For $n = 0$ the only kernel polynomial is $K^{(0)}_0 = 1$.
\end{definition}

Every $K_y$ is monic of degree $\sum_{k=1}^{n} 2^{k-1} = N-1$.

\begin{example}\label{ex:n2}
For $n = 2$ (so $N = 4$), $K_b(X) = (X + b_1)(X^2 + b_2)$, and
$K_0 = X^3$, \;
$K_1 = X^3 + X^2$, \;
$K_2 = X^3 + X$, \;
$K_3 = X^3 + X^2 + X + 1$.
\end{example}

\begin{lemma}[Coefficient rule]\label{lem:coeff-rule}
For every $y \in \F^n$ and $a \in \iv{0}{N-1}$,
\[
  \coef{a}{K_y} \;=\; \prod_{k \in \iv{1}{n} \,:\, a_k = 0} y_k ,
\]
with the empty product equal to $1$.
\end{lemma}

\begin{proof}
Expanding the product, each term chooses from factor $k$ either $X^{2^{k-1}}$ or $y_k$.
Encode the choice by $a_k = 1$ for $X^{2^{k-1}}$ and $a_k = 0$ for $y_k$.
The term is then $\bigl(\prod_{k : a_k = 0} y_k\bigr) X^{\sum_k a_k 2^{k-1}} = \bigl(\prod_{k : a_k = 0} y_k\bigr) X^{a}$.
Distinct choices give distinct exponents $a$ by uniqueness of the binary decomposition, so no two terms contribute to the same monomial.
\end{proof}

\begin{corollary}[Boolean coefficients]\label{cor:bool-coeff}
For $a, b \in \iv{0}{N-1}$,
\[
  \coef{a}{K_b} \;=\;
  \begin{cases}
    1 & \text{if } \cmp{b} \sub a,\\
    0 & \text{otherwise.}
  \end{cases}
\]
\end{corollary}

\begin{proof}
By \cref{lem:coeff-rule}, $\coef{a}{K_b} = \prod_{k : a_k = 0} b_k$, which is $1$ exactly when $a_k = 0 \Rightarrow b_k = 1$ for all $k$, that is, when $\cmp{b}_k = 1 \Rightarrow a_k = 1$ for all $k$.
\end{proof}

\begin{lemma}[Recursion]\label{lem:recursion}
Let $t \in \iv{0}{n}$, and write $b = c + 2^t h$ with $c \in \iv{0}{2^t-1}$ and $h \in \iv{0}{2^{n-t}-1}$ (\cref{sec:prelim:bits}). Then
\[
  K^{(n)}_{c + 2^t h}(X) \;=\; K^{(t)}_c(X)\; K^{(n-t)}_h\bigl(X^{2^t}\bigr).
\]
In particular, for $n \ge 1$ and $b = b_1 + 2b'$, $K^{(n)}_b(X) = (X + b_1)\, K^{(n-1)}_{b'}(X^2)$.
\end{lemma}

\begin{proof}
The bits of $b$ are $b_k = c_k$ for $k \le t$ and $b_k = h_{k-t}$ for $k > t$.
Split the product of \cref{def:kernel} into the factors $k \le t$, whose product is $K^{(t)}_c(X)$, and the factors $k > t$.
For $k > t$, $X^{2^{k-1}} + h_{k-t} = (X^{2^t})^{2^{(k-t)-1}} + h_{k-t}$, and the product of these factors over $k \in \iv{t+1}{n}$ is $K^{(n-t)}_h(X^{2^t})$.
The last statement is the case $t = 1$, since $K^{(1)}_{b_1} = X + b_1$.
\end{proof}

\subsection{Kronecker products}\label{sec:kernel:kron}

Matrices below have entries in $\F$, and their rows and columns are indexed by $\{0,1\}$ (for $2 \times 2$ matrices) or by $\iv{0}{N-1}$ (for $N \times N$ matrices).

\begin{definition}[Kronecker product]\label{def:kron}
For $2 \times 2$ matrices $S_1,\dots,S_n$, the $N \times N$ matrix $S_1 \otimes \dots \otimes S_n$ has entries
\[
  (S_1 \otimes \dots \otimes S_n)[a,b] \;=\; \prod_{k=1}^{n} S_k[a_k, b_k], \qquad a, b \in \iv{0}{N-1},
\]
where factor $k$ acts on bit $k$. For $n = 0$ it is the $1 \times 1$ identity matrix.
We write $S^{\otimes n}$ when all factors equal $S$, and $S_{(k)} = I \otimes \dots \otimes S \otimes \dots \otimes I$ (with $S$ in position $k$ and the $2 \times 2$ identity $I$ elsewhere).
\end{definition}

\begin{lemma}[Mixed-product rule]\label{lem:kron}
For $2 \times 2$ matrices $S_1,\dots,S_n$ and $T_1,\dots,T_n$,
\[
  (S_1 \otimes \dots \otimes S_n)(T_1 \otimes \dots \otimes T_n) \;=\; (S_1T_1) \otimes \dots \otimes (S_nT_n).
\]
In particular $I^{\otimes n}$ is the identity, $S^{\otimes n}T^{\otimes n} = (ST)^{\otimes n}$, and $S^{\otimes n} = S_{(1)} S_{(2)} \cdots S_{(n)}$, where the factors $S_{(k)}$ commute.
\end{lemma}

\begin{proof}
Entry $[a,b]$ of the left-hand side is
\[
  \sum_{c=0}^{N-1} \prod_{k=1}^{n} S_k[a_k,c_k]\,T_k[c_k,b_k] \;=\; \prod_{k=1}^{n} \sum_{\beta \in \{0,1\}} S_k[a_k,\beta]\,T_k[\beta,b_k],
\]
by distributivity, since $c \mapsto (c_1,\dots,c_n)$ is a bijection from $\iv{0}{N-1}$ onto $\{0,1\}^n$.
The right-hand side is $\prod_k (S_kT_k)[a_k,b_k]$, which is the same.
For the particular cases: $I^{\otimes n}[a,b] = \prod_k [a_k = b_k]$ equals $1$ if $a = b$ and $0$ otherwise, by uniqueness of the bit vector; $S_{(1)}\cdots S_{(n)}$ is, by induction on the number of factors, the Kronecker product whose factor in position $k$ is the product of $n$ matrices of which exactly one is $S$ and the others are $I$; and any two $S_{(k)}, S_{(k')}$ commute because both products equal the Kronecker product with $S$ in positions $k$ and $k'$.
\end{proof}

\subsection{The basis theorem}\label{sec:kernel:basis}

Let
\[
  C \;=\; \begin{pmatrix} 0 & 1 \\ 1 & 1 \end{pmatrix},
  \qquad
  C^{-1} \;=\; \begin{pmatrix} -1 & 1 \\ 1 & 0 \end{pmatrix},
\]
so that $C[0,\beta] = \beta$ and $C[1,\beta] = 1$ for $\beta \in \{0,1\}$; one checks $C^{-1}C = CC^{-1} = I$.
For $U \in \F[X]_{<N}$ we write $\mathrm{coef}(U) = (\coef{a}{U})_{a \in \iv{0}{N-1}}$ for its \emph{coefficient vector}; the map $\mathrm{coef}$ is an $\F$-linear bijection from $\F[X]_{<N}$ onto $\F^{\iv{0}{N-1}}$.

\begin{theorem}[Boolean-kernel basis]\label{thm:basis}
Let $n \ge 0$, and let $\mathbf{K}_n$ be the $N \times N$ matrix with entries $\mathbf{K}_n[a,b] = \coef{a}{K_b}$ for $a, b \in \iv{0}{N-1}$ (rows indexed by coefficient positions, columns by kernel indices).
\begin{enumerate}
  \item $\mathbf{K}_n = C^{\otimes n}$.
  \item $\mathbf{K}_n$ is invertible over every field, with inverse $(C^{-1})^{\otimes n}$.
  \item $(K_b)_{b \in \iv{0}{N-1}}$ is a basis of $\F[X]_{<N}$.
\end{enumerate}
\end{theorem}

\begin{proof}
(1) By \cref{lem:coeff-rule}, $\coef{a}{K_b} = \prod_{k : a_k = 0} b_k = \prod_{k} C[a_k,b_k]$. For $n = 0$ both sides are the $1\times 1$ identity.
(2) By \cref{lem:kron}, $(C^{-1})^{\otimes n} C^{\otimes n} = (C^{-1}C)^{\otimes n} = I^{\otimes n}$, and likewise in the other order.
(3) Each $K_b$ has degree $N-1$ and so lies in $\F[X]_{<N}$, and its coefficient vector is column $b$ of $\mathbf{K}_n$. These $N$ columns are linearly independent by (2), and $\dim \F[X]_{<N} = N$.
\end{proof}

The proof uses no property of $\F$; so $(K_b)_b$ is a basis of $\F[X]_{<N}$ over every field $\F$, of any characteristic.

\begin{remark}[Determinant]\label{rem:det}
For $n \ge 1$, $\det \mathbf{K}_n = (-1)^{n2^{n-1}}$, and $\det \mathbf{K}_0 = 1$.
Indeed, $C^{\otimes n} = C_{(1)} \cdots C_{(n)}$ by \cref{lem:kron}, and each $C_{(k)}$ is, after a simultaneous permutation of rows and columns grouping the indices into the $N/2$ pairs that differ only in bit $k$, block diagonal with $N/2$ blocks equal to $C$; so $\det C_{(k)} = (\det C)^{N/2} = (-1)^{2^{n-1}}$, since $\det C = -1$. Multiplying over $k$ gives the claim. This fact is not used below.
\end{remark}

\begin{definition}[Kernel coordinates]\label{def:kernel-coords}
For $U \in \F[X]_{<N}$, its \emph{kernel coordinates} are the unique table $\lambda : \iv{0}{N-1} \to \F$ with
$U = \sum_{b=0}^{N-1} \lambda(b)\, K_b$ (\cref{thm:basis}).
With the coefficient vector $\mathrm{coef}(U)$, we have $\mathrm{coef}(U) = C^{\otimes n}\lambda$ and $\lambda = (C^{-1})^{\otimes n}\, \mathrm{coef}(U)$.
\end{definition}

Since $(C^{-1})^{\otimes n}$ has integer entries, the kernel coordinates of a polynomial with coefficients in a subfield $\F_0 \subseteq \F$ lie in $\F_0$, and do not depend on whether it is regarded in $\F_0[X]$ or in $\F[X]$.

\begin{corollary}[Explicit transforms]\label{cor:transforms}
For $U = \sum_b \lambda(b) K_b \in \F[X]_{<N}$:
\[
  \coef{a}{U} \;=\; \sum_{\substack{b \in \iv{0}{N-1} \\ \cmp{b} \sub a}} \lambda(b),
  \qquad\qquad
  \lambda(b) \;=\; \sum_{\substack{a \in \iv{0}{N-1} \\ a \sub \cmp{b}}} (-1)^{\wt(\cmp{b}) - \wt(a)}\, \coef{a}{U} .
\]
Each of the two transforms is computed with exactly $\tfrac{n}{2}N$ field additions or subtractions and no multiplications.
\end{corollary}

\begin{proof}
The first formula is \cref{cor:bool-coeff}.
For the second, $C^{-1}[\beta,\beta']$ is nonzero only when $(\beta,\beta') \ne (1,1)$, and then equals $-1$ if $\beta = \beta' = 0$ and $1$ otherwise.
Hence $(C^{-1})^{\otimes n}[b,a] = \prod_k C^{-1}[b_k,a_k]$ is nonzero only if $a \sub \cmp{b}$, and then equals $(-1)^{|\{k : a_k = b_k = 0\}|} = (-1)^{\wt(\cmp{b}) - \wt(a)}$.
For the cost, $C^{\otimes n} = C_{(1)} \cdots C_{(n)}$ by \cref{lem:kron}.
The matrix $C_{(k)}$ acts on the $N/2$ pairs of entries whose indices differ only in bit $k$: a pair with values $(x, x')$, where $x$ sits at the index with bit $k$ equal to $0$, becomes $(x', x + x')$, which costs one addition.
Likewise $(C^{-1})_{(k)}$ maps $(x,x')$ to $(x' - x, x)$, and $(C^{-1})^{\otimes n} = (C^{-1})_{(1)}\cdots (C^{-1})_{(n)}$.
\end{proof}

\subsection{Low degree in kernel coordinates}\label{sec:kernel:lowdeg}

\begin{lemma}[The constant polynomial]\label{lem:one}
For every $n \ge 0$,
\[
  1 \;=\; \sum_{h=0}^{2^n-1} (-1)^{n - \wt(h)}\, K^{(n)}_h(X).
\]
\end{lemma}

\begin{proof}
Expand $1 = \prod_{k=1}^{n}\bigl( (X^{2^{k-1}} + 1) - X^{2^{k-1}} \bigr)$ by choosing, in factor $k$, either $X^{2^{k-1}}+1$ (bit $h_k = 1$) or $-X^{2^{k-1}}$ (bit $h_k = 0$). The term indexed by $h$ is $(-1)^{n-\wt(h)} K^{(n)}_h(X)$.
\end{proof}

\begin{theorem}[Low-degree criterion]\label{thm:lowdeg}
Let $t \in \iv{0}{n}$ and $U = \sum_{b=0}^{N-1} \lambda(b)\, K^{(n)}_b \in \F[X]_{<N}$.
Then $\deg U < 2^t$ if and only if there is a table $\varphi : \iv{0}{2^t-1} \to \F$ such that
\begin{equation}\label{eq:character}
  \lambda(c + 2^t h) \;=\; (-1)^{(n-t) - \wt(h)}\, \varphi(c)
  \qquad \text{for all } c \in \iv{0}{2^t-1},\ h \in \iv{0}{2^{n-t}-1}.
\end{equation}
In that case $\varphi$ is unique, $\varphi(c) = \lambda(c + N - 2^t)$, and $\varphi$ is the table of kernel coordinates of $U$ with $t$ variables: $U = \sum_{c=0}^{2^t-1} \varphi(c)\, K^{(t)}_c$.
\end{theorem}

\begin{proof}
By \cref{lem:recursion},
\[
  U(X) \;=\; \sum_{c=0}^{2^t-1} K^{(t)}_c(X)\, Y_c\bigl(X^{2^t}\bigr),
  \qquad
  Y_c(Z) \;=\; \sum_{h=0}^{2^{n-t}-1} \lambda(c+2^t h)\, K^{(n-t)}_h(Z) \;\in\; \F[Z]_{<2^{n-t}} .
\]
Put $R_i = \sum_c \coef{i}{Y_c}\, K^{(t)}_c \in \F[X]_{<2^t}$ for $i \in \iv{0}{2^{n-t}-1}$, so that $U(X) = \sum_{i} X^{2^t i} R_i(X)$.
By \cref{lem:splits}(2) (with $m = 2^{n-t}$), this is the block split of $U$, so $\deg U < 2^t$ if and only if $R_i = 0$ for all $i \ge 1$.
By \cref{thm:basis} with $t$ variables, this holds if and only if $\coef{i}{Y_c} = 0$ for all $c$ and all $i \ge 1$, that is, if and only if every $Y_c$ is a constant $\varphi(c) \in \F$.

By \cref{lem:one} applied with $n-t$ variables, the constant $\varphi(c)$ has kernel coordinates $h \mapsto (-1)^{(n-t)-\wt(h)} \varphi(c)$.
By uniqueness of kernel coordinates (\cref{thm:basis} with $n-t$ variables), $Y_c = \varphi(c)$ holds if and only if $\lambda(c+2^t h) = (-1)^{(n-t)-\wt(h)}\varphi(c)$ for all $h$, which is \eqref{eq:character}.
Taking $h = 2^{n-t}-1$ (so that $\wt(h) = n-t$) gives $\varphi(c) = \lambda(c + N-2^t)$, and then $U = \sum_c \varphi(c) K^{(t)}_c$.
\end{proof}

\begin{corollary}[Reed--Solomon codes in kernel coordinates]\label{cor:rs-kernel}
Let $t \in \iv{0}{n}$, and let $L$ be a smooth domain of order $M \ge 2^t$. Then
\[
  \RS_\F[L, 2^t] \;=\; \Bigl\{ \ev_L\Bigl(\sum_{b} \lambda(b) K_b\Bigr) \;:\; \lambda : \iv{0}{N-1} \to \F \text{ satisfies \eqref{eq:character} for some } \varphi \Bigr\}.
\]
\end{corollary}

\begin{proof}
Immediate from \cref{thm:lowdeg}, since every polynomial of degree $< 2^t$ lies in $\F[X]_{<N}$ and so has kernel coordinates.
\end{proof}

\subsection{Evaluation as a scaled multilinear evaluation}\label{sec:kernel:eval}

The next result shows that evaluating $U$ at a point $\zeta \in \F$ amounts, up to an explicit scalar, to evaluating the multilinear extension $\tilde\lambda$ of its kernel coordinates (\cref{prop:mle}) at a point $\psi(\zeta)$ of an explicit curve.

\begin{proposition}[Evaluation identity]\label{prop:eval}
Let $U = \sum_{b} \lambda(b) K_b \in \F[X]_{<N}$, and let $\zeta \in \F$ satisfy $\gamma_k(\zeta) := 2\zeta^{2^{k-1}} + 1 \ne 0$ for all $k \in \iv{1}{n}$. Define $\psi(\zeta) = (\psi_1(\zeta),\dots,\psi_n(\zeta)) \in \F^n$ by
\[
  \psi_k(\zeta) \;=\; \frac{\zeta^{2^{k-1}} + 1}{2\zeta^{2^{k-1}} + 1}, \qquad k \in \iv{1}{n}.
\]
Then
\[
  U(\zeta) \;=\; \Bigl( \prod_{k=1}^{n} \gamma_k(\zeta) \Bigr) \cdot \tilde\lambda\bigl(\psi(\zeta)\bigr).
\]
\end{proposition}

\begin{proof}
Fix $k$ and $\beta \in \{0,1\}$, and write $\gamma = \gamma_k(\zeta)$ and $\psi = \psi_k(\zeta)$, so that $\gamma \psi = \zeta^{2^{k-1}} + 1$.
Then $\gamma\, \eq_1(\beta, \psi) = \gamma\bigl(\beta \psi + (1-\beta)(1-\psi)\bigr)$ equals $\gamma \psi = \zeta^{2^{k-1}} + 1$ for $\beta = 1$, and $\gamma - \gamma \psi = \zeta^{2^{k-1}}$ for $\beta = 0$; in both cases it equals $\zeta^{2^{k-1}} + \beta$.
Taking the product over $k$ gives $K_b(\zeta) = \bigl(\prod_k \gamma_k(\zeta)\bigr)\, \eq(b, \psi(\zeta))$ for every $b \in \iv{0}{N-1}$.
Multiplying by $\lambda(b)$ and summing over $b$ gives the claim by \cref{prop:mle}.
\end{proof}

\begin{remark}\label{rem:eval-degenerate}
The condition $\gamma_k(\zeta) \ne 0$ for all $k$ excludes at most $N-1$ points of $\F$, namely the roots of $\prod_k (2X^{2^{k-1}} + 1)$, a polynomial of degree $N-1$ with leading coefficient $2^n \ne 0$ (\cref{lem:roots}).
At an excluded point, $U(\zeta) = \sum_b \lambda(b) K_b(\zeta)$ still holds, but the scalar vanishes and the identity carries no information about $\tilde\lambda$.
\end{remark}

\section{The Fold Dictionary}\label{sec:fold}

This section proves the central structural fact of the paper.
In kernel coordinates, a suitably normalised FRI fold is \emph{exactly} the restriction of the first variable of the multilinear extension of the coordinate table (\cref{def:restriction,lem:restriction}).
Consequently, folding a Reed--Solomon codeword $n$ times evaluates the multilinear extension of its kernel coordinates, with no change of representation.
As in \cref{sec:kernel}, $\F \supseteq \Fq$ is a finite field.
We define two folds: $\mathrm{pfold}_r$ acts on polynomials and $\fold_r$ acts on words; \cref{prop:fold-consistency} relates them.

\subsection{Even--odd split in kernel coordinates}\label{sec:fold:split}

\begin{proposition}[Even--odd split]\label{prop:split}
Let $n \ge 1$ and $U = \sum_{b=0}^{N-1} \lambda(b)\, K^{(n)}_b \in \F[X]_{<N}$, and let $U(X) = U_e(X^2) + X\,U_o(X^2)$ be its even--odd split, with $U_e, U_o \in \F[X]_{<N/2}$ (\cref{lem:splits}). Then
\[
  U_e \;=\; \sum_{b'=0}^{N/2-1} \lambda(2b'+1)\, K^{(n-1)}_{b'},
  \qquad
  U_o \;=\; \sum_{b'=0}^{N/2-1} \bigl(\lambda(2b') + \lambda(2b'+1)\bigr)\, K^{(n-1)}_{b'} .
\]
\end{proposition}

\begin{proof}
By \cref{lem:recursion}, $K^{(n)}_{2b'}(X) = X\, K^{(n-1)}_{b'}(X^2)$ and $K^{(n)}_{2b'+1}(X) = K^{(n-1)}_{b'}(X^2) + X\, K^{(n-1)}_{b'}(X^2)$. Hence
\[
  U(X) \;=\; \sum_{b'} \lambda(2b'+1)\, K^{(n-1)}_{b'}(X^2) \;+\; X \sum_{b'} \bigl(\lambda(2b') + \lambda(2b'+1)\bigr)\, K^{(n-1)}_{b'}(X^2).
\]
Both sums over $b'$ are polynomials of degree $< N/2$ evaluated at $X^2$, so the claim follows from the uniqueness in \cref{lem:splits}(1).
\end{proof}

\subsection{The kernel fold of a polynomial}\label{sec:fold:fold}

\begin{definition}[Kernel fold of a polynomial]\label{def:fold-poly}
Let $d \ge 1$. For $r \in \F$ and $U \in \F[X]_{<2d}$ with even--odd split $U = U_e(X^2) + X U_o(X^2)$, define
\[
  \mathrm{pfold}_r(U) \;=\; (2r-1)\,U_e \;+\; (1-r)\,U_o \;\in\; \F[X]_{<d}.
\]
\end{definition}

The map $U \mapsto \mathrm{pfold}_r(U)$ is $\F$-linear, by \cref{lem:splits}.
By the coefficient formulas of \cref{lem:splits}(1), $U_e$ and $U_o$, and hence $\mathrm{pfold}_r(U)$, do not depend on the choice of $d$ with $U \in \F[X]_{<2d}$: equivalently, $\coef{i}{U_e} = \coef{2i}{U}$ and $\coef{i}{U_o} = \coef{2i+1}{U}$ for all $U \in \F[X]$ and $i \in \N$.

\begin{theorem}[Fold dictionary]\label{thm:fold-dictionary}
Let $n \ge 1$, $U = \sum_{b=0}^{N-1} \lambda(b)\, K^{(n)}_b \in \F[X]_{<N}$ and $r \in \F$. Then the kernel coordinates of $\mathrm{pfold}_r(U)$, with $n-1$ variables, are the restriction $\res_r(\lambda)$ (\cref{def:restriction}):
\[
  \mathrm{pfold}_r(U) \;=\; \sum_{b'=0}^{N/2-1} \res_r(\lambda)(b')\, K^{(n-1)}_{b'},
  \qquad
  \res_r(\lambda)(b') \;=\; (1-r)\,\lambda(2b') + r\,\lambda(2b'+1).
\]
Consequently, $\widetilde{\res_r(\lambda)}(y) = \tilde\lambda(r, y)$ for $y = (y_1,\dots,y_{n-1})$ (\cref{lem:restriction}).
\end{theorem}

\begin{proof}
By \cref{prop:split}, the coordinate of $K^{(n-1)}_{b'}$ in $\mathrm{pfold}_r(U)$ is
$(2r-1)\lambda(2b'+1) + (1-r)\bigl(\lambda(2b') + \lambda(2b'+1)\bigr) = (1-r)\lambda(2b') + r\lambda(2b'+1)$.
\end{proof}

\begin{corollary}[Iterated folding evaluates the multilinear extension]\label{cor:iterated}
Let $U = \sum_{b=0}^{N-1} \lambda(b)\, K^{(n)}_b \in \F[X]_{<N}$ and $r = (r_1,\dots,r_n) \in \F^n$.
Define $U_0 = U$ and $U_j = \mathrm{pfold}_{r_j}(U_{j-1})$ for $j \in \iv{1}{n}$.
Then, for $j \in \iv{0}{n}$, the kernel coordinates of $U_j$ (with $n-j$ variables) are the iterated restriction $\res_{r_{\le j}}(\lambda)$, that is,
\[
  U_j \;=\; \sum_{b=0}^{2^{n-j}-1} \tilde\lambda(r_{\le j}, b)\, K^{(n-j)}_b .
\]
In particular, $U_n$ is the constant polynomial $\tilde\lambda(r_1,\dots,r_n)$.
\end{corollary}

\begin{proof}
Induction on $j$, using \cref{thm:fold-dictionary} for the step and \cref{def:restriction} for the iterated restriction; the displayed formula is \cref{lem:restriction}(2).
For $j = n$, $K^{(0)}_0 = 1$.
\end{proof}

\subsection{The kernel fold of a word}\label{sec:fold:words}

Let $L$ be any smooth domain of order $M \ge 2$; the definitions below are used for all such domains, in particular for the domains $L^{2^j}$.
For a word $w \in \F^L$, let $w_e, w_o \in \F^{L^2}$ be its even and odd parts \eqref{eq:even-odd}.

\begin{definition}[Kernel fold of a word]\label{def:fold-word}
For $r \in \F$ and $w \in \F^L$, define $\fold_r(w) \in \F^{L^2}$ by
\[
  \fold_r(w) \;=\; (2r-1)\,w_e + (1-r)\,w_o .
\]
\end{definition}

\begin{lemma}[Properties of the word fold]\label{lem:fold-word}
Let $w \in \F^L$ and $r \in \F$.
\begin{enumerate}
  \item (Explicit formula.) For every $\xi \in L$,
  \begin{equation}\label{eq:fold-explicit}
    \fold_r(w)(\xi^2) \;=\; \frac{\bigl((2r-1)\xi + (1-r)\bigr)\, w(\xi) \;+\; \bigl((2r-1)\xi - (1-r)\bigr)\, w(-\xi)}{2\xi}.
  \end{equation}
  \item (Linearity.) $w \mapsto \fold_r(w)$ is $\F$-linear.
  \item (Locality.) For $\eta \in L^2$, $\fold_r(w)(\eta)$ depends only on the values of $w$ on the fibre over $\eta$ (\cref{def:fibre}); a verifier computes it from two queries.
\end{enumerate}
\end{lemma}

\begin{proof}
(1) Substitute \eqref{eq:even-odd} into the definition and put everything over the denominator $2\xi$:
$(2r-1)\,\frac{w(\xi)+w(-\xi)}{2} = \frac{(2r-1)\xi\,(w(\xi)+w(-\xi))}{2\xi}$ and $(1-r)\,\frac{w(\xi)-w(-\xi)}{2\xi}$; collect the coefficients of $w(\xi)$ and $w(-\xi)$.
(2) and (3) follow from \cref{lem:split-words}(1,2).
\end{proof}

\begin{lemma}[Line form]\label{lem:line}
For $w \in \F^L$, put $u^0 = w_o - w_e$ and $u^1 = 2w_e - w_o$, both in $\F^{L^2}$. Then $\fold_r(w) = u^0 + r\,u^1$ for every $r \in \F$.
Conversely, $w_e = u^0 + u^1$ and $w_o = 2u^0 + u^1$.
\end{lemma}

\begin{proof}
$(2r-1)w_e + (1-r)w_o = (w_o - w_e) + r(2w_e - w_o)$; and $u^0 + u^1 = w_e$, $2u^0 + u^1 = w_o$.
\end{proof}

\begin{proposition}[Consistency with the polynomial fold]\label{prop:fold-consistency}
Let $1 \le d \le M/2$, $r \in \F$ and $U \in \F[X]_{<2d}$. Then
\[
  \fold_r(\ev_L(U)) \;=\; \ev_{L^2}(\mathrm{pfold}_r(U)).
\]
In particular, $\fold_r$ maps $\RS_\F[L,2d]$ into $\RS_\F[L^2,d]$, and $P_{\fold_r(c)} = \mathrm{pfold}_r(P_c)$ for every $c \in \RS_\F[L,2d]$.
\end{proposition}

\begin{proof}
By \cref{lem:split-words}(4), $\ev_L(U)_e = \ev_{L^2}(U_e)$ and $\ev_L(U)_o = \ev_{L^2}(U_o)$; since $\ev_{L^2}$ is linear, $\fold_r(\ev_L(U)) = \ev_{L^2}\bigl((2r-1)U_e + (1-r)U_o\bigr)$.
Moreover $\mathrm{pfold}_r(U) \in \F[X]_{<d}$ and $d \le |L^2|$, so the last statement follows from \cref{lem:rs-params}(1).
\end{proof}

\begin{corollary}[Folding a committed codeword]\label{cor:codeword-fold}
Let $L$ be a smooth domain of order $M = 2^{\mu}$ with $\mu \ge n$, let $\lambda : \iv{0}{N-1} \to \F$, and let $w_0 = \ev_L(U)$ with $U = \sum_b \lambda(b) K_b$.
For $r \in \F^n$, define $w_j = \fold_{r_j}(w_{j-1}) \in \F^{L^{2^j}}$ for $j \in \iv{1}{n}$.
Then $w_j = \ev_{L^{2^j}}(U_j)$, with $U_j$ as in \cref{cor:iterated}; in particular $w_n$ is the constant word on $L^{2^n}$ with value $\tilde\lambda(r_1,\dots,r_n)$.
\end{corollary}

\begin{proof}
Induction on $j$. For the step, apply \cref{prop:fold-consistency} on the domain $L^{2^{j-1}}$, of order $M/2^{j-1} \ge 2$ (\cref{lem:power}), with $2d = N/2^{j-1} \le M/2^{j-1}$ and $(L^{2^{j-1}})^2 = L^{2^j}$. Conclude with \cref{cor:iterated}.
\end{proof}

\subsection{Relation to classical FRI folding}\label{sec:fold:fri}

\begin{definition}[Classical FRI fold]\label{def:cfold}
The classical FRI fold~\cite{BBHR18} with challenge $\theta \in \F$ is $\mathrm{cfold}_\theta(U) = U_e + \theta\, U_o$ for $U \in \F[X]$ (which, as for $\mathrm{pfold}$, does not depend on a degree bound), and $\mathrm{cfold}_\theta(w) = w_e + \theta\, w_o$ for a word $w$ on a smooth domain of order at least $2$.
\end{definition}

\begin{proposition}[Classical versus kernel fold]\label{prop:classical}
\begin{enumerate}
  \item Let $n \ge 1$, and let $U \in \F[X]_{<N}$ have kernel coordinates $\lambda$. Then the kernel coordinates of $\mathrm{cfold}_\theta(U)$, with $n-1$ variables, are
  \[
    b' \;\longmapsto\; \theta\,\lambda(2b') + (1+\theta)\,\lambda(2b'+1).
  \]
  \item As maps $\F[X]_{<2d} \to \F[X]_{<d}$, for every $d \ge 1$: if $\theta \ne -1/2$, then $\mathrm{cfold}_\theta = (1+2\theta)\,\mathrm{pfold}_{r}$ with $r = (1+\theta)/(1+2\theta)$. Conversely, if $r \ne 1/2$, then $\mathrm{pfold}_r = (2r-1)\,\mathrm{cfold}_{\theta}$ with $\theta = (1-r)/(2r-1)$.
  \item Let $U = \sum_b \lambda(b) K_b \in \F[X]_{<N}$ and $\theta_1,\dots,\theta_n \in \F \setminus \{-1/2\}$. Applying $\mathrm{cfold}_{\theta_1},\dots,\mathrm{cfold}_{\theta_n}$ successively to $U$ gives the constant
  \[
    \prod_{j=1}^{n} (1+2\theta_j) \cdot \tilde\lambda\Bigl( \tfrac{1+\theta_1}{1+2\theta_1}, \dots, \tfrac{1+\theta_n}{1+2\theta_n} \Bigr).
  \]
  \item For $\zeta \in \F$, applying $\mathrm{cfold}_{\theta_1},\dots,\mathrm{cfold}_{\theta_n}$ with $\theta_j = \zeta^{2^{j-1}}$ to $U$ gives the constant $U(\zeta)$.
\end{enumerate}
\end{proposition}

\begin{proof}
(1) By \cref{prop:split}, the coordinate of $K^{(n-1)}_{b'}$ in $U_e + \theta U_o$ is $\lambda(2b'+1) + \theta(\lambda(2b') + \lambda(2b'+1))$.
(2) With $r = (1+\theta)/(1+2\theta)$ we have $2r - 1 = 1/(1+2\theta)$ and $1 - r = \theta/(1+2\theta)$, so $(1+2\theta)\bigl((2r-1)U_e + (1-r)U_o\bigr) = U_e + \theta U_o$. The converse is the same computation with $\theta = (1-r)/(2r-1)$, for which $1+2\theta = 1/(2r-1)$ and $(1+\theta)/(1+2\theta) = r$.
(3) By (2) and the linearity of $\mathrm{cfold}$ and $\mathrm{pfold}$, the composition equals $\prod_j (1+2\theta_j)$ times $\mathrm{pfold}_{r_n} \circ \dots \circ \mathrm{pfold}_{r_1}(U)$ with $r_j = (1+\theta_j)/(1+2\theta_j)$; conclude with \cref{cor:iterated}.
(4) Let $V_0 = U$ and $V_j = \mathrm{cfold}_{\theta_j}(V_{j-1})$. We show $V_{j-1}(\zeta^{2^{j-1}}) = U(\zeta)$ for $j \in \iv{1}{n+1}$ by induction: the case $j = 1$ is trivial, and $V_{j-1}(X) = (V_{j-1})_e(X^2) + X(V_{j-1})_o(X^2)$ gives $V_{j-1}(\zeta^{2^{j-1}}) = (V_{j-1})_e(\zeta^{2^j}) + \zeta^{2^{j-1}}(V_{j-1})_o(\zeta^{2^j}) = V_j(\zeta^{2^j})$. For $j = n+1$, $V_n$ is a constant, equal to $U(\zeta)$.
\end{proof}

For $\zeta$ with $\gamma_j(\zeta) \ne 0$ for all $j$, items 3 and 4 with $\theta_j = \zeta^{2^{j-1}}$ give a second proof of \cref{prop:eval}: $1 + 2\theta_j = \gamma_j(\zeta)$ and $(1+\theta_j)/(1+2\theta_j) = \psi_j(\zeta)$.

\begin{remark}[A consistency pitfall]\label{rem:pitfall}
By \cref{prop:classical}(2), the two folds differ by a scalar and by a reparametrisation of the challenge.
An implementation must therefore use \emph{the same} fold on the committed word and on the table that the verifier tracks.
Folding the word with $\mathrm{cfold}_\theta$ while restricting the table with $r = \theta$ (or folding with $\mathrm{pfold}_\theta$ while tracking $\mathrm{cfold}_\theta$) yields in general a final constant different from $\tilde\lambda(r_1,\dots,r_n)$, even for an honest prover; see \cref{ex:pitfall}.
In an early version of our implementation this mismatch caused honest proofs to fail the final consistency check; \cref{def:fold-word} together with \cref{thm:fold-dictionary} fixes the convention.
\end{remark}

\begin{example}\label{ex:pitfall}
Over $\mathbb{F}_{17}$ with $n = 1$, take $\lambda = (\lambda(0), \lambda(1)) = (1, 0)$, so $U = K_0 = X$, $U_e = 0$ and $U_o = 1$.
With $\theta = 2$, the classical fold gives $\mathrm{cfold}_2(U) = 2$, while the restriction at $r = 2$ gives $\res_2(\lambda)(0) = (1-2)\cdot 1 + 2 \cdot 0 = -1 = 16$.
The kernel fold agrees with the restriction: $\mathrm{pfold}_2(U) = 3 \cdot 0 + (1-2)\cdot 1 = 16$. So the classical final constant $2$ differs from $\tilde\lambda(2) = 16$.
\end{example}

\begin{remark}[Comparison with coefficient-form folding]\label{rem:coeff-form}
In the monomial basis, $\mathrm{cfold}_\theta$ acts on coefficients by $\coef{a'}{\mathrm{cfold}_\theta(U)} = \coef{2a'}{U} + \theta\,\coef{2a'+1}{U}$ (\cref{lem:splits}(1)).
By induction on $n$, applying $\mathrm{cfold}_{\theta_1},\dots,\mathrm{cfold}_{\theta_n}$ gives $\sum_a \coef{a}{U} \prod_{k : a_k = 1} \theta_k = \sum_a \coef{a}{U}\, m_a(\theta_1,\dots,\theta_n)$, the evaluation at $(\theta_1,\dots,\theta_n)$ of the multilinear polynomial whose \emph{coefficient form} (\cref{def:coeff-form}) is $\mathrm{coef}(U)$.
This is the mechanism behind Reed--Solomon instantiations of BaseFold~\cite{ZCF24} and behind DeepFold~\cite{GLHQTZ25}.
\Cref{cor:codeword-fold} is its counterpart for the \emph{table form}: if the prover encodes a table $f$ as $U = \sum_b f(b) K_b$, folding computes $\tilde f(r_1,\dots,r_n)$ directly, without converting $f$ to coefficient form.
\end{remark}

\section{The Polynomial Commitment Scheme}\label{sec:pcs}

We now use the fold dictionary to build a multilinear polynomial commitment scheme for tables, in the sense of \cref{def:pcs}.
The prover encodes a table $f$ as the Reed--Solomon codeword of $U_f = \sum_b f(b) K_b$.
An evaluation claim $\tilde f(z) = v$ is proved by a sumcheck (\cref{sec:proofs:sumcheck}) whose round challenges are also the fold challenges, exactly as in BaseFold~\cite{ZCF24}.
By \cref{cor:codeword-fold}, no change of representation is needed: the folds of the committed word track the partially restricted table that the sumcheck prover already holds.
We present the scheme as a public-coin IOP (\cref{def:iop}); its compilation into a non-interactive scheme is discussed in \cref{sec:proofs:bcs,sec:sound:pq}.

\subsection{Parameters and encoding}\label{sec:pcs:params}

\paragraph{Fields.}
The smooth domain lies in the base field $\Fq$, while tables, words, the evaluation point and the challenges live in a finite field $\F \supseteq \Fq$, the \emph{challenge field}.
Taking $\F = \Fq$ gives the plain scheme; taking $\F = \K$, an extension of $\Fq$, gives the scheme with extension-field challenges used in \cref{sec:impl}, in which honest tables have entries in $\Fq$.
All results of \cref{sec:pcs,sec:sound} hold for every such $\F$, with $|\F|$ in the error bounds.

\paragraph{Parameters.}
\begin{itemize}
  \item the number of variables $n \ge 1$, and $N = 2^n$;
  \item a rate $\rho = 2^{-R}$ with $R \ge 1$, and a smooth domain $L \le \Fq^\times$ of order $M = N/\rho = 2^{n+R}$ (\cref{def:domain}; it exists if and only if $2^{n+R} \mid q-1$, by \cref{lem:power});
  \item a number of folding rounds $\ell \in \iv{0}{n}$;
  \item a number of queries $\kappa \ge 1$.
\end{itemize}
For $j \in \iv{0}{\ell}$ let $L_j = L^{2^j}$, the smooth domain of order $M_j = M/2^j \ge 2^R \ge 2$ (\cref{lem:power}), let $N_j = N/2^j$, and let $\Cc_j = \RS_\F[L_j, N_j]$.
Every code $\Cc_j$ has the same rate $N_j/M_j = \rho$, and $L_{j+1} = (L_j)^2$.

\begin{definition}[Encoding]\label{def:encoding}
The \emph{encoding} of a table $f : \iv{0}{N-1} \to \F$ is the word
\[
  \Enc(f) \;=\; \ev_L(U_f) \;\in\; \Cc_0,
  \qquad
  U_f \;=\; \sum_{b=0}^{N-1} f(b)\, K_b \;\in\; \F[X]_{<N} .
\]
\end{definition}

\begin{lemma}[Encoding]\label{lem:encoding}
$\Enc$ is an $\F$-linear bijection from the tables over $\F$ onto $\Cc_0$, and for $c \in \Cc_0$, the table $f$ with $\Enc(f) = c$ is the table of kernel coordinates of $P_c$.
\end{lemma}

\begin{proof}
$f \mapsto U_f$ is a linear bijection onto $\F[X]_{<N}$ by \cref{thm:basis}, and $\ev_L$ is a linear bijection from $\F[X]_{<N}$ onto $\Cc_0$ by \cref{lem:rs-params}(1), with inverse $c \mapsto P_c$.
\end{proof}

$\Enc$ is computed by the transform of \cref{cor:transforms} ($\tfrac{n}{2}N$ additions), followed by one number-theoretic transform (NTT, the fast evaluation of a polynomial on a smooth domain) of size $M$ on the zero-padded coefficient vector.

\subsection{The evaluation protocol}\label{sec:pcs:protocol}

The commitment to $f$ is the word $w_0 = \Enc(f)$.
The evaluation protocol $\Pieval$ proves the claim $\tilde f(z) = v$ for a public point $z = (z_1,\dots,z_n) \in \F^n$ and a public value $v \in \F$.
As in \cref{def:pcs}, its instance is $(w_0; z, v)$, and the verifier has oracle access to $w_0$.

\paragraph{Round structure.}
$\Pieval$ is a public-coin IOP with $\ell+1$ rounds.
\begin{itemize}
  \item In round $j \in \iv{1}{\ell}$, the prover sends a \emph{round polynomial} $s_j \in \F[X]_{<3}$, together with, if $j \ge 2$, an oracle $w_{j-1} \in \F^{L_{j-1}}$; the verifier replies with a uniformly random challenge $r_j \in \F$.
  \item In round $\ell+1$, the prover sends a \emph{final table} $g : \iv{0}{N_\ell-1} \to \F$ in the clear; the verifier replies with $\kappa$ independent uniformly random \emph{query points} $\xi^{(1)}_0,\dots,\xi^{(\kappa)}_0 \in L$, and then decides.
\end{itemize}
Thus the challenge sets are $\mathsf{Ch}_j = \F$ for $j \in \iv{1}{\ell}$ and $\mathsf{Ch}_{\ell+1} = L^{\times\kappa}$, the set of $\kappa$-tuples of elements of $L$ (a Cartesian power, not the image of a power map).
A round polynomial is transmitted as its values at $0$, $1$ and $2$, which are distinct because the characteristic is odd. The map $\F[X]_{<3} \to \F^3$, $s \mapsto (s(0), s(1), s(2))$, is injective by \cref{lem:roots}, hence bijective since both spaces have dimension $3$; so every triple sent by a prover identifies a unique $s_j \in \F[X]_{<3}$.
As in \cref{sec:proofs:iop}, the checks below are predicates on the full transcript.

\paragraph{Honest prover.}
The honest prover keeps two tables over $\F$, initialised to
\[
  f_0 = f \qquad\text{and}\qquad \chi_0 = \bigl(b \mapsto \eq(b,z)\bigr), \qquad b \in \iv{0}{N-1},
\]
and after round $j$ replaces them by their restrictions (\cref{def:restriction}) at $r_j$: $f_j = \res_{r_j}(f_{j-1})$ and $\chi_j = \res_{r_j}(\chi_{j-1})$, tables on $\iv{0}{N_j-1}$.
In round $j$ it sends
\begin{equation}\label{eq:honest-sj}
  s_j(X) \;=\; \sum_{b=0}^{N_j - 1}
      \bigl((1-X)\, f_{j-1}(2b) + X f_{j-1}(2b+1)\bigr)\,
      \bigl((1-X)\, \chi_{j-1}(2b) + X \chi_{j-1}(2b+1)\bigr),
\end{equation}
and, if $j \ge 2$, the oracle $w_{j-1} = \fold_{r_{j-1}}(w_{j-2})$ (\cref{def:fold-word}).
In round $\ell+1$ it sends $g = f_\ell$.

\paragraph{Verifier.}
Set $v_0 = v$ and $v_j = s_j(r_j)$ for $j \in \iv{1}{\ell}$.
Let $G = \sum_{b=0}^{N_\ell - 1} g(b)\, K^{(n-\ell)}_b \in \F[X]_{<N_\ell}$ and $w_\ell = \ev_{L_\ell}(G)$; the verifier computes values of $w_\ell$ itself.
The verifier accepts if and only if all of the following checks pass.
\begin{enumerate}
  \item \emph{Sumcheck.} $s_j(0) + s_j(1) = v_{j-1}$ for every $j \in \iv{1}{\ell}$.
  \item \emph{Closure.}
  \[
    v_\ell \;=\; \Bigl(\prod_{j=1}^{\ell} \eq_1(r_j, z_j)\Bigr)\cdot \tilde g\bigl(z_{\ell+1},\dots,z_{n}\bigr).
  \]
  \item \emph{Queries.} For every $i \in \iv{1}{\kappa}$, with $\xi_0 = \xi^{(i)}_0$ and $\xi_j = \xi_0^{2^j}$ for $j \in \iv{1}{\ell}$:
  \begin{itemize}
    \item if $\ell \ge 1$: for every $j \in \iv{1}{\ell}$, the value $\fold_{r_j}(w_{j-1})(\xi_j)$, computed by \eqref{eq:fold-explicit} from the two queried values $w_{j-1}(\pm\xi_{j-1})$, equals $w_j(\xi_j)$ (a query to the oracle $w_j$ if $j < \ell$, and $G(\xi_\ell)$ if $j = \ell$);
    \item if $\ell = 0$: $w_0(\xi_0) = G(\xi_0)$.
  \end{itemize}
\end{enumerate}

\begin{remark}[Stopping parameter]\label{rem:stop}
With $\ell = n$, the final table is the single value $\tilde f(r_1,\dots,r_n)$ and $G$ is a constant; with $\ell = 0$, the prover sends $f$ in the clear.
Intermediate values trade the $N_\ell = 2^{n-\ell}$ field elements of the final message against $\ell$ folding rounds.
\end{remark}

\subsection{Completeness}\label{sec:pcs:complete}

Recall from \cref{sec:prelim:bits} the notation $(r_{\le j}, b) = (r_1,\dots,r_j,b_1,\dots,b_{n-j})$ for an index $b$ with $n-j$ bits.

\begin{lemma}[Honest prover]\label{lem:honest}
Let $f$ be a table over $\F$, $z \in \F^n$, and $r_1,\dots,r_\ell \in \F$, and let $f_j$, $\chi_j$ and $s_j$ be computed by the honest prover. Then, for $j \in \iv{0}{\ell}$ and $b \in \iv{0}{N_j-1}$:
\begin{enumerate}
  \item $f_j(b) = \tilde f(r_{\le j}, b)$ and $\chi_j(b) = \eq\bigl((r_{\le j}, b), z\bigr)$;
  \item for $j \ge 1$, $s_j$ is the honest round polynomial $s^*_j$ of \cref{sec:proofs:sumcheck} for $F(x) = \tilde f(x)\,\eq(x,z)$, and $s_j \in \F[X]_{<3}$;
  \item for $j \ge 1$, $s_j(r_j) = \sum_b f_j(b)\,\chi_j(b)$, and $s_j(0) + s_j(1) = \sum_b f_{j-1}(b)\,\chi_{j-1}(b)$;
  \item $\sum_b f_0(b)\chi_0(b) = \tilde f(z)$ and, for every $j$,
  \[
    \sum_{b=0}^{N_j-1} f_j(b)\,\chi_j(b) \;=\; \Bigl(\prod_{i=1}^{j}\eq_1(r_i,z_i)\Bigr)\cdot \widetilde{f_j}(z_{j+1},\dots,z_n).
  \]
\end{enumerate}
\end{lemma}

\begin{proof}
(1) \Cref{lem:restriction}(2), and \cref{cor:eq-mle} for $\chi_0$.
(2) Regard $f_{j-1}$ and $\chi_{j-1}$ as tables over the ring $\F[X] \supseteq \F$ (\cref{sec:prelim:ml}). The two factors in \eqref{eq:honest-sj} are then $\res_X(f_{j-1})(b)$ and $\res_X(\chi_{j-1})(b)$, and \cref{lem:restriction}(2) over $\F[X]$, with (1), gives
\[
  s_j(X) \;=\; \sum_{b=0}^{N_j-1} \tilde f(r_1,\dots,r_{j-1}, X, b_1,\dots,b_{n-j})\;\eq\bigl((r_1,\dots,r_{j-1},X,b_1,\dots,b_{n-j}),z\bigr) \;=\; s^*_j(X).
\]
Each summand of \eqref{eq:honest-sj} is a product of two polynomials of degree $\le 1$ in $X$, so $s_j \in \F[X]_{<3}$.
(3) By \cref{lem:natural} applied to the evaluation $X \mapsto r_j$, the factors become $f_j(b)$ and $\chi_j(b)$. Setting $X = 0$ and $X = 1$ in \eqref{eq:honest-sj} gives $\sum_b \bigl(f_{j-1}(2b)\chi_{j-1}(2b) + f_{j-1}(2b+1)\chi_{j-1}(2b+1)\bigr)$.
(4) The first identity is \cref{prop:mle}. For the second, (1) and \cref{lem:eq}(3) give
\[
  \chi_j(b) \;=\; \Bigl(\prod_{i=1}^{j}\eq_1(r_i,z_i)\Bigr)\,\eq\bigl(b,(z_{j+1},\dots,z_n)\bigr);
\]
conclude with \cref{prop:mle}.
\end{proof}

\begin{theorem}[Completeness]\label{thm:complete}
If $w_0 = \Enc(f)$, $v = \tilde f(z)$ and the prover is honest, the verifier accepts with probability $1$.
\end{theorem}

\begin{proof}
\emph{Sumcheck.} By \cref{lem:honest}(3,4), $v_0 = \tilde f(z) = \sum_b f_0(b)\chi_0(b)$ and $v_j = s_j(r_j) = \sum_b f_j(b)\chi_j(b)$, so $s_j(0) + s_j(1) = \sum_b f_{j-1}(b)\chi_{j-1}(b) = v_{j-1}$ for every $j \in \iv{1}{\ell}$.

\emph{Closure.} By \cref{lem:honest}(4) with $j = \ell$ and $g = f_\ell$, $v_\ell = \prod_{j}\eq_1(r_j,z_j) \cdot \tilde g(z_{\ell+1},\dots,z_n)$.

\emph{Queries.} If $\ell = 0$, then $g = f$, $G = U_f$ and $w_0 = \Enc(f) = \ev_L(G)$.
If $\ell \ge 1$, extend $(r_1,\dots,r_\ell)$ arbitrarily to $r \in \F^n$, and let $U_j$ be as in \cref{cor:iterated} for $U = U_f$; its kernel coordinates are $f_j$.
Since $L$ has order $2^{n+R}$ with $n+R \ge n$, \cref{cor:codeword-fold} shows that the successive folds of $w_0$ satisfy $\fold_{r_j}(\cdots\fold_{r_1}(w_0)) = \ev_{L_j}(U_j)$ for $j \in \iv{1}{\ell}$.
The honest oracles are these folds for $j < \ell$, and $U_\ell = \sum_b f_\ell(b) K^{(n-\ell)}_b = G$ because $g = f_\ell$.
Hence $\fold_{r_j}(w_{j-1}) = w_j$ on all of $L_j$ for every $j \in \iv{1}{\ell}$, and every check of every query passes.
\end{proof}

\subsection{Costs}\label{sec:pcs:costs}

Let $t_{\mathsf{H}}$ denote the cost of one hash evaluation. We count field operations up to constant factors.

\begin{itemize}
  \item \emph{Commit.} $\tfrac{n}{2}N$ additions (\cref{cor:transforms}), one NTT of size $M$, that is, $O(M\log M)$ operations, and a Merkle tree over the $M/2$ fibres of $w_0$ (one leaf per fibre), $O(M)\cdot t_{\mathsf{H}}$.
  \item \emph{Prover, evaluation.} Sumcheck: $O(N)$ in total, since each round is linear in the current table size and the sizes halve.
  Folds: $\fold_{r_j}(w_{j-1})$ costs $O(M_j)$ using precomputed inverses of $2\xi$ for $\xi \in L$; $O(M)$ in total.
  Merkle trees for $w_1,\dots,w_{\ell-1}$: $O(M)\cdot t_{\mathsf{H}}$ in total.
  \item \emph{Verifier.} Sumcheck: $O(\ell)$. Closure: $O(N_\ell)$. Queries: $O(\kappa \ell)$ field operations plus $O(\kappa \ell \log M)\cdot t_{\mathsf{H}}$ for authentication paths, plus $\kappa$ evaluations of $G$, $O(\kappa N_\ell)$.
  \item \emph{Proof size.} For $\ell \ge 1$: $3\ell$ field elements for the sumcheck, $N_\ell$ for the final table, $\ell-1$ Merkle roots besides the commitment, and $\kappa\ell$ opened leaves ($2\kappa \ell$ values) with their authentication paths. For $\ell = 0$: the table $g = f$ and $\kappa$ opened values of $w_0$.
\end{itemize}

\begin{remark}[Relation to BaseFold and DeepFold]\label{rem:basefold}
The interleaving of sumcheck rounds with fold rounds under shared challenges is the BaseFold paradigm~\cite{ZCF24}, and the asymptotic costs above match Reed--Solomon instantiations of BaseFold and DeepFold~\cite{GLHQTZ25}.
In particular, converting a table to coefficient form by \cref{lem:mobius} also costs $\tfrac{n}{2}N$ additions, so the kernel encoding offers no asymptotic advantage.
What differs is structural: by \cref{thm:fold-dictionary}, the committed word, the sumcheck table and the folds are three views of \emph{one} object, the restricted table $f_j$, and the consistency between them is an identity rather than a bookkeeping convention (compare \cref{rem:pitfall}).
\end{remark}

\begin{remark}[Batching and zero knowledge]\label{rem:batch}
Because $\Enc$ and $\fold_r$ are linear, several tables can be opened at the same point by running $\Pieval$ on a random linear combination of them, the verifier answering queries to the combined word from queries to each committed word.
We do not analyse batching in this paper, and zero knowledge is out of scope.
\end{remark}

\section{Soundness}\label{sec:sound}

We prove that $\Pieval$ is sound, evaluation binding, and round-by-round knowledge sound, in the unique-decoding regime.
The only external result about codes used is the correlated-agreement theorem (\cref{thm:bciks}) of~\cite{BCIKS23}, which enters through the line form of the kernel fold (\cref{lem:line}); the efficiency of the extractor also uses \cref{fact:decode}.
Throughout this section we use the notation of \cref{sec:pcs}: the challenge field $\F \supseteq \Fq$, the domains $L_j$ of order $M_j$, and the codes $\Cc_j = \RS_\F[L_j, N_j]$, all of rate $\rho$.
We fix a proximity parameter $0 < \delta \le (1-\rho)/2$.
We assume $1 \le \ell \le n$, except in \cref{sec:sound:l0}, which treats $\ell = 0$.
Probabilities are over the challenge space $\F^\ell \times L^{\times\kappa}$ of $\Pieval$ (\cref{sec:proofs:iop,sec:pcs:protocol}).

\subsection{Fibre distance and decoding}\label{sec:sound:fibre}

For $j \in \iv{0}{\ell-1}$, the domain $L_j$ has order $M_j \ge 2$, and its fibres (\cref{def:fibre}) are indexed by the $M_{j+1}$ points of $L_{j+1} = (L_j)^2$.

\begin{definition}[Fibre distance]\label{def:fibre-dist}
Let $j \in \iv{0}{\ell-1}$. For $w, c \in \F^{L_j}$,
\[
  \Delta^{\mathrm{fib}}_j(w,c) \;=\; \frac{\bigl|\{\, \eta \in L_{j+1} \;:\; w \text{ and } c \text{ differ somewhere on the fibre over } \eta \,\}\bigr|}{M_{j+1}},
\]
and $\Delta^{\mathrm{fib}}_j(w, \Cc_j) = \min_{c \in \Cc_j} \Delta^{\mathrm{fib}}_j(w,c)$.
\end{definition}

\begin{lemma}[Fibre distance]\label{lem:fibre-dist}
Let $j \in \iv{0}{\ell-1}$ and $w, c \in \F^{L_j}$.
\begin{enumerate}
  \item $\Delta(w,c) \le \Delta^{\mathrm{fib}}_j(w,c)$.
  \item $\Delta^{\mathrm{fib}}_j(w,c) \le \delta$ if and only if $w$ and $c$ agree on every fibre over some set of at least $(1-\delta)M_{j+1}$ points of $L_{j+1}$.
\end{enumerate}
\end{lemma}

\begin{proof}
(1) Each point of $L_j$ at which $w$ and $c$ differ lies in a fibre over which they differ, and each of these $\Delta^{\mathrm{fib}}_j(w,c)\,M_{j+1}$ fibres has two points, so $\Delta(w,c)\,M_j \le 2\,\Delta^{\mathrm{fib}}_j(w,c)\,M_{j+1} = \Delta^{\mathrm{fib}}_j(w,c)\,M_j$.
(2) The set of points over whose fibre $w$ and $c$ agree has $(1 - \Delta^{\mathrm{fib}}_j(w,c))\,M_{j+1}$ elements.
\end{proof}

\begin{lemma}[Unique fibre decoding]\label{lem:unique}
Let $j \in \iv{0}{\ell-1}$. For every $w \in \F^{L_j}$ there is at most one $c \in \Cc_j$ with $\Delta^{\mathrm{fib}}_j(w,c) \le \delta$.
\end{lemma}

\begin{proof}
Such a $c$ satisfies $\Delta(w,c) \le \delta \le (1-\rho)/2$ by \cref{lem:fibre-dist}(1); apply \cref{lem:rs-params}(3).
\end{proof}

\begin{definition}[Decoded codeword and decoded table]\label{def:decoded}
Let $j \in \iv{0}{\ell-1}$ and $w \in \F^{L_j}$ with $\Delta^{\mathrm{fib}}_j(w, \Cc_j) \le \delta$.
The \emph{decoded codeword} $\mathrm{dec}_j(w)$ is the unique $c \in \Cc_j$ of \cref{lem:unique}.
For $j = 0$, the \emph{decoded table} of $w$ is the unique table $f^*$ over $\F$ with $\Enc(f^*) = \mathrm{dec}_0(w)$, that is, the table of kernel coordinates of $P_{\mathrm{dec}_0(w)}$ (\cref{lem:encoding}).
\end{definition}

\begin{remark}[Computing the decoded table]\label{rem:ext}
Since $\Delta(w, \mathrm{dec}_0(w)) \le \delta < \tfrac12(1 - (N-1)/M)$ (\cref{lem:rs-params}), the algorithm of \cref{fact:decode} finds $\mathrm{dec}_0(w)$; an inverse NTT gives $P_{\mathrm{dec}_0(w)}$, and \cref{cor:transforms} gives $f^*$.
Checking whether $\Delta^{\mathrm{fib}}_0(w, \mathrm{dec}_0(w)) \le \delta$ costs $O(M)$ comparisons.
Conversely, for an arbitrary $w$, if the decoder fails or its output fails this check, then no codeword is within fibre distance $\delta$ of $w$; so decoding followed by the check decides whether $\Delta^{\mathrm{fib}}_0(w, \Cc_0) \le \delta$.
If $w \in \Fq^L$ (as for a commitment computed over the base field), then $f^*$ has entries in $\Fq$, by \cref{lem:descent} and the integrality of \cref{cor:transforms}.
\end{remark}

\subsection{Two folding lemmas}\label{sec:sound:lemmas}

\begin{lemma}[Folding far words]\label{lem:far}
Let $j \in \iv{0}{\ell-1}$ and $w \in \F^{L_j}$ with $\Delta^{\mathrm{fib}}_j(w, \Cc_j) > \delta$. Then
\[
  \bigl|\{\, r \in \F \;:\; \Delta\bigl(\fold_r(w),\, \Cc_{j+1}\bigr) \le \delta \,\}\bigr| \;\le\; M_{j+1}.
\]
\end{lemma}

\begin{proof}
Suppose the set has more than $M_{j+1}$ elements. By \cref{lem:line}, $\fold_r(w) = u^0 + r u^1$ with $u^0 = w_o - w_e$ and $u^1 = 2w_e - w_o$.
Apply \cref{thm:bciks} to the code $\Cc_{j+1}$ (length $M_{j+1}$, rate $\rho$): there are $L' \subseteq L_{j+1}$ with $|L'| \ge (1-\delta)M_{j+1}$ and codewords $\hat c^0, \hat c^1 \in \Cc_{j+1}$ with $u^0 = \hat c^0$ and $u^1 = \hat c^1$ on $L'$.
Let $T_e = P_{\hat c^0} + P_{\hat c^1}$ and $T_o = 2P_{\hat c^0} + P_{\hat c^1}$ in $\F[X]_{<N_{j+1}}$, and $T(X) = T_e(X^2) + X T_o(X^2) \in \F[X]_{<N_j}$.
By \cref{lem:line}, $w_e = T_e$ and $w_o = T_o$ on $L'$, so for every $\eta \in L'$ and $\xi \in L_j$ with $\xi^2 = \eta$, \eqref{eq:even-odd-inverse} gives
$w(\pm \xi) = T_e(\eta) \pm \xi T_o(\eta) = T(\pm \xi)$.
Thus $w$ agrees with the codeword $\ev_{L_j}(T) \in \Cc_j$ on every fibre over $L'$, so $\Delta^{\mathrm{fib}}_j(w, \Cc_j) \le \delta$, a contradiction.
\end{proof}

\begin{lemma}[Fibre errors survive folding]\label{lem:survive}
Let $j \in \iv{0}{\ell-1}$, $w, c \in \F^{L_j}$, and $\eta \in L_{j+1}$ such that $w$ and $c$ differ somewhere on the fibre over $\eta$. Then there is at most one $r \in \F$ with $\fold_r(w)(\eta) = \fold_r(c)(\eta)$.
\end{lemma}

\begin{proof}
Let $y = w - c$. By \cref{lem:fold-word}(2) and \cref{lem:line}, $\fold_r(w)(\eta) - \fold_r(c)(\eta) = \fold_r(y)(\eta) = \bigl(y_o(\eta) - y_e(\eta)\bigr) + r\bigl(2y_e(\eta) - y_o(\eta)\bigr)$.
This is a polynomial of degree $\le 1$ in $r$. If it were identically zero, the inverse formulas of \cref{lem:line} would give $y_e(\eta) = y_o(\eta) = 0$, so $y$ would vanish on the fibre over $\eta$ (\cref{lem:split-words}(3)). Hence it has at most one root.
\end{proof}

\subsection{Good challenges}\label{sec:sound:good}

\begin{definition}[Good challenge]\label{def:good}
Let $j \in \iv{1}{\ell}$, $w \in \F^{L_{j-1}}$ and $r \in \F$. We say that $r$ is \emph{good for $w$ at level $j$}, and write $\mathsf{Good}_j(w,r)$, if:
\begin{enumerate}
  \item in case $\Delta^{\mathrm{fib}}_{j-1}(w, \Cc_{j-1}) > \delta$: $\Delta(\fold_r(w), \Cc_{j}) > \delta$;
  \item in case $\Delta^{\mathrm{fib}}_{j-1}(w, \Cc_{j-1}) \le \delta$, with $c = \mathrm{dec}_{j-1}(w)$: for every $\eta \in L_j$ such that $w$ and $c$ differ somewhere on the fibre over $\eta$, $\fold_r(w)(\eta) \ne \fold_r(c)(\eta)$.
\end{enumerate}
\end{definition}

\begin{lemma}[Few bad challenges]\label{lem:good}
For every $j \in \iv{1}{\ell}$ and $w \in \F^{L_{j-1}}$, at most $M_j$ values $r \in \F$ are not good for $w$ at level $j$.
\end{lemma}

\begin{proof}
In case 1, this is \cref{lem:far}. In case 2, by \cref{lem:survive}, each of the at most $M_j$ points $\eta$ over which $w$ and $c$ differ excludes at most one $r$.
\end{proof}

\begin{lemma}[One folding step]\label{lem:step}
Let $j \in \iv{1}{\ell}$, $w \in \F^{L_{j-1}}$, and $r \in \F$ good for $w$ at level $j$.
Let $c' \in \Cc_j$ and $Y = \{\eta \in L_j : \fold_r(w)(\eta) = c'(\eta)\}$, and suppose $|Y| \ge (1-\delta) M_j$.
Then $\Delta^{\mathrm{fib}}_{j-1}(w, \Cc_{j-1}) \le \delta$; the codeword $c = \mathrm{dec}_{j-1}(w)$ satisfies $\fold_r(c) = c'$; and $w = c$ on the fibre over every $\eta \in Y$.
\end{lemma}

\begin{proof}
We have $\Delta(\fold_r(w), c') \le \delta$, so case 1 of \cref{def:good} is excluded: $\Delta^{\mathrm{fib}}_{j-1}(w, \Cc_{j-1}) \le \delta$.
By \cref{prop:fold-consistency} (with $2N_j = N_{j-1} \le M_{j-1}$), $\fold_r(c) \in \Cc_j$.
By locality (\cref{lem:fold-word}(3)), $\fold_r(c)$ agrees with $\fold_r(w)$ at every $\eta$ over whose fibre $w = c$, hence on at least $(1-\delta)M_j$ points; and $\fold_r(w)$ agrees with $c'$ on $Y$.
By \cref{lem:metric}, $\fold_r(c)$ and $c'$ agree on at least $(1-2\delta)M_j \ge \rho M_j = N_j$ points, so they are equal (\cref{lem:rs-params}(2)).
Finally, for $\eta \in Y$, $\fold_r(w)(\eta) = c'(\eta) = \fold_r(c)(\eta)$, so by case 2 of \cref{def:good}, $w = c$ on the fibre over $\eta$.
\end{proof}

\subsection{The soundness theorem}\label{sec:sound:thm}

Fix a deterministic prover $\mathsf{P}^*$ (\cref{lem:randomised} extends the results to randomised provers) and an instance $(w_0; z, v)$ with $w_0 \in \F^L$.
The prover's messages are functions of the earlier challenges: $s_j$ of $r_{\le j-1}$, the oracle $w_j$ ($j \in \iv{1}{\ell-1}$) of $r_{\le j}$, and $g$ of $r_{\le \ell}$.
As in \cref{sec:pcs:protocol}, $G = \sum_b g(b) K^{(n-\ell)}_b$ and $w_\ell = \ev_{L_\ell}(G) \in \Cc_\ell$.

\paragraph{Passing sets.}
For fixed challenges $r \in \F^\ell$, let $\Pp_\ell = L_\ell$ and, for $j = \ell-1, \dots, 0$,
\[
  \Pp_j \;=\; \bigl\{\, \xi \in L_j \;:\; \xi^2 \in \Pp_{j+1} \text{ and } \fold_{r_{j+1}}(w_j)(\xi^2) = w_{j+1}(\xi^2) \,\bigr\}.
\]
Let $\pi = \pi(r) = |\Pp_0|/M$. We record these facts.

\begin{lemma}[Passing sets]\label{lem:passing}
For fixed $r \in \F^\ell$:
\begin{enumerate}
  \item for $j \in \iv{0}{\ell}$, a point $\xi_0 \in L$ has $\xi_0^{2^j} \in \Pp_j$ if and only if the query checks at levels $j+1,\dots,\ell$ pass for $\xi_0$; in particular $\xi_0$ passes all the query checks if and only if $\xi_0 \in \Pp_0$;
  \item for $j < \ell$, $\Pp_j$ is a union of fibres, and the squaring map sends it two-to-one onto a set $\Pp_j^2 \subseteq \Pp_{j+1}$;
  \item $|\Pp_j| \ge \pi M_j$ for every $j \in \iv{0}{\ell}$.
\end{enumerate}
\end{lemma}

\begin{proof}
(1) Downward induction on $j$, from the recursive definition. (2) The defining condition of $\Pp_j$ depends only on $\xi^2$, and squaring is two-to-one on $L_j$ (\cref{lem:power}). (3) By (2), $|\Pp_j|/M_j = |\Pp_j^2|/M_{j+1} \le |\Pp_{j+1}|/M_{j+1}$; iterate from $j = 0$.
\end{proof}

\begin{lemma}[The codeword chain]\label{lem:chain}
Fix $r \in \F^\ell$ such that $r_j$ is good for $w_{j-1}$ at level $j$ for every $j \in \iv{1}{\ell}$, and $\pi > 1-\delta$.
Then $\Delta^{\mathrm{fib}}_j(w_j, \Cc_j) \le \delta$ for every $j \in \iv{0}{\ell-1}$, and the codewords $c_j = \mathrm{dec}_j(w_j)$ for $j < \ell$ and $c_\ell = w_\ell$ satisfy $\fold_{r_{j+1}}(c_j) = c_{j+1}$ and $w_j = c_j$ on $\Pp_j$, for every $j \in \iv{0}{\ell-1}$.
Moreover, $g = \res_r(f^*)$, where $f^*$ is the decoded table of $w_0$.
\end{lemma}

\begin{proof}
We show by downward induction on $j \in \iv{0}{\ell}$ that $w_j = c_j$ on $\Pp_j$, together with the other claims for $j < \ell$. For $j = \ell$ this is trivial.
Let $j < \ell$. For $\eta \in \Pp_j^2 \subseteq \Pp_{j+1}$, $\fold_{r_{j+1}}(w_j)(\eta) = w_{j+1}(\eta) = c_{j+1}(\eta)$ by the definition of $\Pp_j$ and the induction hypothesis.
So the set $Y$ of \cref{lem:step} (with $w = w_j$, $r = r_{j+1}$, $c' = c_{j+1}$, at level $j+1$) contains $\Pp_j^2$, of size $|\Pp_j|/2 \ge \pi M_{j+1} > (1-\delta)M_{j+1}$ (\cref{lem:passing}).
\Cref{lem:step} gives $\Delta^{\mathrm{fib}}_j(w_j,\Cc_j) \le \delta$, $\fold_{r_{j+1}}(c_j) = c_{j+1}$, and $w_j = c_j$ on every fibre over $\Pp_j^2$, that is, on $\Pp_j$.

For the last claim, extend $r$ arbitrarily to a point of $\F^n$. Then $c_0 = \Enc(f^*) = \ev_L(U_{f^*})$, and by \cref{cor:codeword-fold} (with $\mu = n + R \ge n$), $c_\ell$, obtained from $c_0$ by $\ell$ folds, is $\ev_{L_\ell}(U_\ell)$, where $U_\ell$ has kernel coordinates $\res_{r}(f^*)$ (\cref{cor:iterated}).
Since $c_\ell = w_\ell = \ev_{L_\ell}(G)$ and both $G$ and $U_\ell$ lie in $\F[X]_{<N_\ell}$ with $N_\ell \le M_\ell$, \cref{lem:rs-params}(1) gives $G = U_\ell$, and \cref{thm:basis} gives $g = \res_r(f^*)$.
\end{proof}

\begin{theorem}[Soundness]\label{thm:sound}
Let $1 \le \ell \le n$, and put
\[
  \varepsilon_{\mathrm{fold}} \;=\; \sum_{j=1}^{\ell} \frac{M_j}{|\F|} \;=\; \frac{M\,(1 - 2^{-\ell})}{|\F|} \;<\; \frac{M}{|\F|}.
\]
\begin{enumerate}
  \item If $\Delta^{\mathrm{fib}}_0(w_0, \Cc_0) > \delta$, then $\mathsf{V}$ accepts with probability at most $\varepsilon_{\mathrm{fold}} + (1-\delta)^{\kappa}$.
  \item If $\Delta^{\mathrm{fib}}_0(w_0, \Cc_0) \le \delta$ and the decoded table $f^*$ of $w_0$ satisfies $\tilde{f^*}(z) \ne v$, then $\mathsf{V}$ accepts with probability at most
  \[
    \varepsilon_{\mathrm{fold}} \;+\; \frac{2\ell}{|\F|} \;+\; (1-\delta)^{\kappa}.
  \]
\end{enumerate}
\end{theorem}

\begin{proof}
\emph{Honest values.} In case 2, let $f^*_j = \res_{r_{\le j}}(f^*)$ and $\chi^*_j = \res_{r_{\le j}}(\chi_0)$ be the tables of the honest prover for $f^*$ (\cref{sec:pcs:protocol}), let $s^*_j$ be its round polynomials, which are those of \cref{sec:proofs:sumcheck} for $F^*(x) = \tilde{f^*}(x)\,\eq(x,z)$ (\cref{lem:honest}), and let $\sigma^*_j = \sum_b f^*_j(b)\,\chi^*_j(b)$ for $j \in \iv{0}{\ell}$.
By \cref{lem:honest}, $\sigma^*_0 = \tilde{f^*}(z)$, $s^*_j(0) + s^*_j(1) = \sigma^*_{j-1}$, $s^*_j(r_j) = \sigma^*_j$, and $\sigma^*_\ell = \bigl(\prod_{j}\eq_1(r_j,z_j)\bigr)\cdot\widetilde{f^*_\ell}(z_{\ell+1},\dots,z_n)$.
Note that $s^*_j$ depends on $r_{\le j-1}$ only.

\emph{Bad events.} For $j \in \iv{1}{\ell}$, let $E_j$ be the set of $r_{\le j}$ such that $r_j$ is not good for $w_{j-1}$ at level $j$, or (in case 2) $s_j \ne s^*_j$ and $s_j(r_j) = s^*_j(r_j)$.
For fixed $r_{\le j-1}$, the word $w_{j-1}$ and the polynomials $s_j, s^*_j$ are fixed, so at most $M_j$ values of $r_j$ (\cref{lem:good}) plus, in case 2, at most $2$ values (\cref{lem:roots}) complete $r_{\le j-1}$ into $E_j$.
By \cref{lem:sequential}, $\Pr[\exists j : r_{\le j} \in E_j] \le \varepsilon_{\mathrm{fold}}$ in case 1 and $\le \varepsilon_{\mathrm{fold}} + 2\ell/|\F|$ in case 2.

\emph{No bad event.} Suppose that no $E_j$ occurs and that $\mathsf{V}$ accepts. We show that $\pi \le 1-\delta$.
Assume instead $\pi > 1-\delta$. Then \cref{lem:chain} applies. In case 1 it gives $\Delta^{\mathrm{fib}}_0(w_0,\Cc_0) \le \delta$, a contradiction.
In case 2 it gives $g = f^*_\ell$, so the closure check reads $v_\ell = \sigma^*_\ell$.
Since $v_0 = v \ne \tilde{f^*}(z) = \sigma^*_0$, there is a largest $j \in \iv{0}{\ell-1}$ with $v_j \ne \sigma^*_j$; then $v_{j+1} = \sigma^*_{j+1}$.
The sumcheck check for index $j+1$ gives $s_{j+1}(0) + s_{j+1}(1) = v_j \ne \sigma^*_j = s^*_{j+1}(0) + s^*_{j+1}(1)$, so $s_{j+1} \ne s^*_{j+1}$, while $s_{j+1}(r_{j+1}) = v_{j+1} = \sigma^*_{j+1} = s^*_{j+1}(r_{j+1})$; that is, $r_{\le j+1} \in E_{j+1}$, a contradiction.

\emph{Conclusion.} Acceptance implies that some $E_j$ occurs, or that $\pi(r) \le 1-\delta$ and all $\kappa$ query points lie in $\Pp_0$ (\cref{lem:passing}(1)).
By \cref{lem:queries} with $\Omega' = \F^\ell$, $\mathcal{A} = \{r : \pi(r) \le 1-\delta\}$ and $\mathcal{S}(r) = \Pp_0$, the second event has probability at most $(1-\delta)^\kappa$.
\end{proof}

\begin{corollary}[Evaluation binding]\label{cor:binding}
Let $1 \le \ell \le n$ and $\varepsilon = \varepsilon_{\mathrm{fold}} + 2\ell/|\F| + (1-\delta)^{\kappa}$.
For every $w_0 \in \F^L$ and $z \in \F^n$, if some prover makes $\mathsf{V}$ accept the instance $(w_0; z,v)$ with probability greater than $\varepsilon$, then $\Delta^{\mathrm{fib}}_0(w_0,\Cc_0) \le \delta$ and $v = \tilde{f^*}(z)$ for the decoded table $f^*$ of $w_0$.
In particular, $\Pieval$ is $\varepsilon$-evaluation binding (\cref{def:binding}).
\end{corollary}

\begin{proof}
Immediate from \cref{thm:sound} and \cref{lem:randomised}; $f^*$ depends on $w_0$ only.
\end{proof}

\subsection{The case \texorpdfstring{$\ell = 0$}{l = 0}}\label{sec:sound:l0}

For $\ell = 0$ there is no fold, and the fibre distance is replaced by the Hamming distance.

\begin{proposition}[No folding]\label{prop:l0}
Let $\ell = 0$. If a deterministic prover makes $\mathsf{V}$ accept $(w_0; z,v)$ with probability greater than $(1-\delta)^{\kappa}$, then $\Delta(w_0, \Cc_0) < \delta$, and $v = \tilde{f^H}(z)$, where $f^H$ is the table such that $\Enc(f^H)$ is the unique codeword within Hamming distance $\delta$ of $w_0$ (\cref{lem:rs-params}(3)).
\end{proposition}

\begin{proof}
For $\ell = 0$ the prover sends $g$ before any challenge, so $g$ and $G$ are fixed, and the verifier accepts if and only if $v = \tilde g(z)$ and all query points lie in $S = \{\xi \in L : w_0(\xi) = G(\xi)\}$.
This has probability $0$ or $(|S|/M)^\kappa$. If it exceeds $(1-\delta)^\kappa$, then $v = \tilde g(z)$ and $|S| > (1-\delta)M$, so $\Delta(w_0, \ev_L(G)) < \delta$; hence $\ev_L(G) = \Enc(g)$ is the unique codeword within distance $\delta$ of $w_0$, and $g = f^H$.
\end{proof}

\subsection{Round-by-round knowledge soundness}\label{sec:sound:rbr}

\Cref{thm:sound} bounds the soundness of the interactive protocol.
We now prove the stronger property of round-by-round knowledge soundness (\cref{def:rbr-knowledge}); \cref{sec:sound:pq} derives from it the relaxed notion of \cref{def:rrbr}, which is the hypothesis of \cref{thm:cdhz}.

\paragraph{Relation.}
Let $1 \le \ell \le n$.
We view $\Pieval$ as an IOP for the relation $\mathcal{R}^{\delta}_{\mathrm{eval}}$ of \cref{sec:proofs:pcs}, with $\mathrm{dist}(w, \Enc(f)) = \Delta^{\mathrm{fib}}_0(w, \Enc(f))$; condition~\eqref{eq:dist-unique} holds by \cref{lem:unique} and the injectivity of $\Enc$.
A table $f$ is a witness for $(w_0;z,v)$ if and only if $\Delta^{\mathrm{fib}}_0(w_0, \Cc_0) \le \delta$, $f$ is the decoded table of $w_0$, and $\tilde f(z) = v$.

\paragraph{Notation.}
The rounds are those of \cref{sec:pcs:protocol}, and $\tau_j$ is the partial transcript after $j$ rounds.
For $j \in \iv{1}{\ell}$ let $\hat w_j = \fold_{r_j}(w_{j-1}) \in \F^{L_j}$, which is determined by $\tau_j$, and for $j \in \iv{1}{\ell-1}$ let $\mathcal{B}_j = \{\eta \in L_j : \hat w_j(\eta) \ne w_j(\eta)\}$, which is determined by $\tau_{j+1}$.
The \emph{sumcheck check of index $j$} is the equation $s_j(0) + s_j(1) = v_{j-1}$, which is determined by $\tau_j$.
For $c \in \Cc_j$, let $\lambda[c]$ be the table of kernel coordinates, with $n-j$ variables, of $P_c$ (\cref{lem:rs-params,def:kernel-coords}).

\paragraph{Witness sets and claimed values.}
For $c \in \Cc_0$ let $\mathcal{W}_0(c) = \{\xi \in L : w_0(\xi) = c(\xi) \text{ and } w_0(-\xi) = c(-\xi)\}$, and for $j \in \iv{1}{\ell}$ and $c \in \Cc_j$ let
\[
  \mathcal{W}_j(c) \;=\; \bigl\{\, \xi \in L \;:\; \xi^{2^i} \notin \mathcal{B}_i \text{ for all } i \in \iv{1}{j-1}, \text{ and } \hat w_j(\xi^{2^j}) = c(\xi^{2^j}) \,\bigr\},
\]
which is determined by $\tau_j$. Let $\mathrm{dens}_j(c) = |\mathcal{W}_j(c)|/M$, and for $j \in \iv{0}{\ell}$ and $c \in \Cc_j$ let
\[
  \sigma_j(c) \;=\; \Bigl(\prod_{i=1}^{j} \eq_1(r_i, z_i)\Bigr) \cdot \widetilde{\lambda[c]}(z_{j+1}, \dots, z_{n}),
\]
the value that the running claim $v_j$ takes if the prover is honest with respect to $c$ (compare \cref{lem:honest}(4)).

\begin{definition}[Doomed set and extractor]\label{def:state}
\begin{itemize}
  \item The empty transcript $\tau_0$ is doomed.
  \item For $j \in \iv{1}{\ell}$, the partial transcript $\tau_j$ is \emph{not doomed} if the sumcheck checks of indices $1,\dots,j$ hold and there is $c \in \Cc_j$ with
  \[
    \mathrm{dens}_j(c) \;\ge\; 1 - \delta
    \qquad\text{and}\qquad
    v_j \;=\; \sigma_j(c);
  \]
  otherwise it is \emph{doomed}.
  \item A full transcript $\tau_{\ell+1}$ is doomed if and only if the verifier rejects it.
\end{itemize}
Let $\Dd$ be the set of doomed transcripts.
The \emph{extractor} $\mathsf{Ext}$, on every input $(\tau, \mathsf{m})$, reads the oracle part $w_0$ of the instance of $\tau$, and outputs the decoded table of $w_0$ if $\Delta^{\mathrm{fib}}_0(w_0,\Cc_0) \le \delta$, and the zero table otherwise.
\end{definition}

By \cref{rem:ext}, $\mathsf{Ext}$ runs in time polynomial in $M$.
Its output is a witness for $(w_0;z,v)$ if and only if $\Delta^{\mathrm{fib}}_0(w_0,\Cc_0) \le \delta$ and the decoded table $f^*$ satisfies $\tilde{f^*}(z) = v$: if $\Delta^{\mathrm{fib}}_0(w_0,\Cc_0) > \delta$, no table is a witness.

\begin{lemma}[One round of the chain]\label{lem:rbr-step}
Let $j \in \iv{1}{\ell}$, fix $\tau_{j-1}$ and the round-$j$ message (hence $w_{j-1}$), and let $r_j$ be good for $w_{j-1}$ at level $j$.
If $c' \in \Cc_j$ satisfies $\mathrm{dens}_{j}(c') \ge 1-\delta$, then $\Delta^{\mathrm{fib}}_{j-1}(w_{j-1}, \Cc_{j-1}) \le \delta$, the codeword $c = \mathrm{dec}_{j-1}(w_{j-1})$ satisfies $\fold_{r_j}(c) = c'$, and $\mathcal{W}_{j}(c') \subseteq \mathcal{W}_{j-1}(c)$; in particular $\mathrm{dens}_{j-1}(c) \ge 1-\delta$.
\end{lemma}

\begin{proof}
Let $Y = \{\eta \in L_{j} : \hat w_j(\eta) = c'(\eta)\}$. Every $\xi \in \mathcal{W}_{j}(c')$ has $\xi^{2^j} \in Y$, and the map $\xi \mapsto \xi^{2^j}$ is $2^j$-to-$1$ from $L$ onto $L_j$ (\cref{lem:power}); so $|Y| \ge |\mathcal{W}_j(c')|/2^j \ge (1-\delta)M_j$.
\Cref{lem:step} gives the first two claims, and $w_{j-1} = c$ on the fibre over every $\eta \in Y$.
Let $\xi \in \mathcal{W}_{j}(c')$, and $\eta = \xi^{2^j} \in Y$; the fibre over $\eta$ is $\{\pm \xi^{2^{j-1}}\}$, on which $w_{j-1} = c$.
If $j = 1$, this says $\xi \in \mathcal{W}_0(c)$.
If $j \ge 2$, the condition $\xi^{2^{j-1}} \notin \mathcal{B}_{j-1}$ (part of the definition of $\mathcal{W}_{j}(c')$) gives $\hat w_{j-1}(\xi^{2^{j-1}}) = w_{j-1}(\xi^{2^{j-1}}) = c(\xi^{2^{j-1}})$, and the conditions $\xi^{2^i} \notin \mathcal{B}_i$ for $i \le j-2$ are common to both sets; hence $\xi \in \mathcal{W}_{j-1}(c)$.
\end{proof}

\begin{lemma}[Honest round polynomial for a codeword]\label{lem:rbr-sumcheck}
Let $j \in \iv{1}{\ell}$, $c \in \Cc_{j-1}$ and $r_1,\dots,r_{j-1} \in \F$, and define
\[
  s^{c}_j(X) \;=\; \Bigl(\prod_{i=1}^{j-1}\eq_1(r_i,z_i)\Bigr) \sum_{b=0}^{N_{j}-1} \widetilde{\lambda[c]}(X, b_1,\dots,b_{n-j})\;\eq\bigl((X,b_1,\dots,b_{n-j}),\, (z_j,\dots,z_{n})\bigr).
\]
Then $s^{c}_j \in \F[X]_{<3}$, $s^{c}_j(0) + s^{c}_j(1) = \sigma_{j-1}(c)$, and $s^{c}_j(r) = \sigma_{j}(\fold_r(c))$ for every $r \in \F$, where $\sigma_j$ is computed with $r_j = r$.
\end{lemma}

\begin{proof}
Each summand is a product of two polynomials of degree $\le 1$ in $X$.
Summing over $X \in \{0,1\}$ and $b$ amounts to summing over $a = X + 2b \in \iv{0}{N_{j-1}-1}$, whose bit vector is $(X, b_1,\dots,b_{n-j})$; this gives $\sum_{a} \lambda[c](a)\, \eq(a, (z_j,\dots,z_n)) = \widetilde{\lambda[c]}(z_j,\dots,z_n)$ by \cref{prop:mle}, times the prefix, that is, $\sigma_{j-1}(c)$.
By \cref{prop:fold-consistency} and \cref{thm:fold-dictionary}, $\lambda[\fold_r(c)] = \res_r(\lambda[c])$, whose multilinear extension is $\widetilde{\lambda[c]}(r,\cdot)$ (\cref{lem:restriction}).
Since $\eq\bigl((r,b), (z_j,\dots,z_n)\bigr) = \eq_1(r,z_j)\,\eq\bigl(b,(z_{j+1},\dots,z_n)\bigr)$ (\cref{lem:eq}(3)), evaluating at $X = r$ (\cref{lem:natural}) gives
\[
  s^c_j(r) \;=\; \Bigl(\prod_{i=1}^{j-1}\eq_1(r_i,z_i)\Bigr)\,\eq_1(r,z_j)\;\widetilde{\lambda[\fold_r(c)]}(z_{j+1},\dots,z_n) \;=\; \sigma_j(\fold_r(c)). \qedhere
\]
\end{proof}

\begin{theorem}[Round-by-round knowledge soundness]\label{thm:rbr}
Let $1 \le \ell \le n$, and let $\Dd$ and $\mathsf{Ext}$ be as in \cref{def:state}.
\begin{enumerate}
  \item (Round $1$.) If the output of $\mathsf{Ext}$ is not a witness for $(w_0; z,v)$, then for every round-$1$ message $s_1$, $\tau_1$ is not doomed with probability at most $(M_1 + 2)/|\F|$ over $r_1$.
  \item (Round $j$, $2 \le j \le \ell$.) If $\tau_{j-1}$ is doomed, then for every round-$j$ message $(w_{j-1}, s_j)$, $\tau_j$ is not doomed with probability at most $(M_j+2)/|\F|$ over $r_j$.
  \item (Round $\ell+1$.) If $\tau_\ell$ is doomed, then for every final message $g$, the verifier accepts with probability at most $(1-\delta)^{\kappa}$ over the query points.
\end{enumerate}
Consequently, $\Pieval$ is round-by-round knowledge sound for $\mathcal{R}^{\delta}_{\mathrm{eval}}$ (\cref{def:rbr-knowledge}) with errors $\varepsilon_j = (M_j+2)/|\F|$ for $j \in \iv{1}{\ell}$ and $\varepsilon_{\ell+1} = (1-\delta)^\kappa$, and in particular with
\[
  \varepsilon_{\mathrm{rbr}} \;=\; \max\Bigl\{ \frac{M_1 + 2}{|\F|},\; (1-\delta)^{\kappa} \Bigr\}.
\]
\end{theorem}

\begin{proof}
\emph{Items 1 and 2.} Fix $j \in \iv{1}{\ell}$, the instance, $\tau_{j-1}$ and the round-$j$ message.
Then $w_{j-1}$, $s_j$ and, if $\Delta^{\mathrm{fib}}_{j-1}(w_{j-1},\Cc_{j-1}) \le \delta$, the codeword $c = \mathrm{dec}_{j-1}(w_{j-1})$ and the polynomial $s^{c}_j$ of \cref{lem:rbr-sumcheck} are fixed.
Call $r_j$ \emph{bad} if it is not good for $w_{j-1}$ at level $j$, or if $\Delta^{\mathrm{fib}}_{j-1}(w_{j-1},\Cc_{j-1}) \le \delta$, $s_j \ne s^{c}_j$ and $s_j(r_j) = s^{c}_j(r_j)$.
By \cref{lem:good,lem:roots}, at most $M_j + 2$ values of $r_j$ are bad.
We show that if $r_j$ is not bad and $\tau_j$ is not doomed, then the hypothesis of the item fails.

Let $c' \in \Cc_j$ witness that $\tau_j$ is not doomed: $\mathrm{dens}_j(c') \ge 1-\delta$, $v_j = \sigma_j(c')$, and the sumcheck checks of indices $1,\dots,j$ hold.
By \cref{lem:rbr-step}, $c$ exists, $\fold_{r_j}(c) = c'$ and $\mathrm{dens}_{j-1}(c) \ge 1-\delta$.
By \cref{lem:rbr-sumcheck}, $s_j(r_j) = v_j = \sigma_j(\fold_{r_j}(c)) = s^{c}_j(r_j)$, so $s_j = s^{c}_j$ because $r_j$ is not bad.
The sumcheck check of index $j$ then gives $v_{j-1} = s_j(0) + s_j(1) = s^{c}_j(0) + s^{c}_j(1) = \sigma_{j-1}(c)$.

For $j \ge 2$, the codeword $c$ therefore witnesses that $\tau_{j-1}$ is not doomed, contradicting the hypothesis of item 2.
For $j = 1$, $\Delta^{\mathrm{fib}}_0(w_0,\Cc_0) \le \delta$ and $c = \mathrm{dec}_0(w_0) = \Enc(f^*)$ for the decoded table $f^*$, so $\lambda[c] = f^*$ and $v = v_0 = \sigma_0(c) = \tilde{f^*}(z)$; thus the output $f^*$ of $\mathsf{Ext}$ is a witness, contradicting the hypothesis of item 1.

\emph{Item 3.} If a sumcheck check or the closure check fails, the verifier rejects. Otherwise, since $\lambda[w_\ell] = g$, the closure check reads $v_\ell = \sigma_\ell(w_\ell)$; as $\tau_\ell$ is doomed, $\mathrm{dens}_\ell(w_\ell) < 1-\delta$.
A query point $\xi$ passes all its checks exactly when $\hat w_i(\xi^{2^i}) = w_i(\xi^{2^i})$ for all $i \in \iv{1}{\ell}$, that is, when $\xi \in \mathcal{W}_\ell(w_\ell)$.
The $\kappa$ query points are independent and uniform, so the verifier accepts with probability at most $\mathrm{dens}_\ell(w_\ell)^{\kappa} < (1-\delta)^{\kappa}$ (\cref{lem:queries}).

\emph{Round-by-round knowledge soundness.} We check the three conditions of \cref{def:rbr-knowledge}.
Condition 1 holds because $\tau_0$ is doomed, and condition 3 because a full transcript is doomed exactly when it is rejected.
For condition 2 in round $1$, item 1 says that an escape probability above $\varepsilon_1$ forces $\mathsf{Ext}$ to output a witness.
In rounds $2,\dots,\ell+1$, items 2 and 3 bound the escape probability from a doomed transcript by $\varepsilon_j$ outright.
Finally, $M_j \le M_1$ for all $j \ge 1$.
\end{proof}

\begin{remark}[Relation to the definition of~\cite{BGTZ24}]\label{rem:bgtz}
\Cref{def:rbr-knowledge} is Definition~3.12 of Block, Garreta, Tiwari and Zaj\k{a}c~\cite{BGTZ24}, with the query positions placed in the challenge of the last round.
Because the relation $\mathcal{R}^{\delta}_{\mathrm{eval}}$ requires the witness to be the table decoded from the commitment, the extracted witness is bound to the commitment; by \cref{lem:rbr-binding}, \cref{thm:rbr} implies $\varepsilon$-evaluation binding with $\varepsilon = \sum_{j=1}^{\ell}(M_j+2)/|\F| + (1-\delta)^\kappa$, which is also the bound of \cref{cor:binding}.
\end{remark}

\begin{proposition}[Round-by-round knowledge soundness for $\ell = 0$]\label{prop:rbr-l0}
Let $\ell = 0$, and define $\mathcal{R}^{\delta}_{\mathrm{eval}}$ with $\mathrm{dist} = \Delta$; condition~\eqref{eq:dist-unique} holds by \cref{lem:rs-params}(3).
Let $\mathsf{Ext}$ output the table $f^H$ of \cref{prop:l0} if $\Delta(w_0, \Cc_0) \le \delta$ (computed with \cref{fact:decode}), and the zero table otherwise, and let $\Dd$ consist of $\tau_0$ and the rejected full transcripts.
Then $\Pieval$ is round-by-round knowledge sound for $\mathcal{R}^{\delta}_{\mathrm{eval}}$ with error $\varepsilon_1 = (1-\delta)^\kappa$.
\end{proposition}

\begin{proof}
Conditions 1 and 3 of \cref{def:rbr-knowledge} hold by construction. For condition 2, fix the message $g$: if the acceptance probability exceeds $(1-\delta)^\kappa$, the proof of \cref{prop:l0} gives $\Delta(w_0, \Cc_0) < \delta$ and $v = \tilde{f^H}(z)$, so the output of $\mathsf{Ext}$ is a witness.
\end{proof}

\subsection{The non-interactive scheme}\label{sec:sound:pq}

We now show that $\Pieval$, written as an interactive oracle reduction, satisfies \cref{def:rrbr}, and apply \cref{thm:cdhz}.
Let $1 \le \ell \le n$ and $\delta^* = (1-\rho)/2$.
Everything proved so far in this section holds for every $\delta \in (0, \delta^*]$; below, $\delta$ varies in $[0,1]$.

\paragraph{Fibre strings.}
For each $j \in \iv{0}{\ell-1}$ and $\eta \in L_{j+1}$, fix one of the two square roots $\xi_\eta \in L_j$ of $\eta$ (\cref{lem:power}); the other is $-\xi_\eta$.
Let $\Sigma = \F^2$, and for $w \in \F^{L_j}$ let $\mathrm{fib}_j(w) \in \Sigma^{L_{j+1}}$ be the string $\eta \mapsto (w(\xi_\eta), w(-\xi_\eta))$, indexed by $L_{j+1}$ in a fixed order.
The map $\mathrm{fib}_j$ is a bijection from $\F^{L_j}$ onto $\Sigma^{M_{j+1}}$, and by \cref{def:fibre-dist},
\begin{equation}\label{eq:fib-hamming}
  \Delta\bigl(\mathrm{fib}_j(w), \mathrm{fib}_j(c)\bigr) \;=\; \Delta^{\mathrm{fib}}_j(w, c) \qquad (w, c \in \F^{L_j}).
\end{equation}
For a partial string $u$ over $\Sigma$, the \emph{filling} $\bar u$ replaces every $\bot$ by $(0,0)$, and $\mathcal{U}(u)$ is the set of positions where $u$ is $\bot$.
Since $\bar u$ and $u$ agree on $\mathrm{Dom}(u)$, $\Delta(\bar u, u') \le \Delta(u, u')$ for every full string $u'$.

\paragraph{The relation.}
Let
\[
  \mathcal{R}_{\mathrm{eval}} \;=\; \bigl\{\, ((z,v),\, y,\, f) \;:\; y = \mathrm{fib}_0(\Enc(f)) \ \text{ and } \ \tilde f(z) = v \,\bigr\},
\]
with explicit instance $(z, v) \in \F^n \times \F$, implicit instance $y \in \Sigma^{M/2}$ and witness a table $f$ over $\F$.
For a partial string $y$ and $\delta \in [0,1]$, $((z,v), y, f) \in (\mathcal{R}_{\mathrm{eval}})_{\le\delta}$ if and only if $\Delta(y, \mathrm{fib}_0(\Enc(f))) \le \delta$ and $\tilde f(z) = v$.
For a full string $y = \mathrm{fib}_0(w_0)$, by \eqref{eq:fib-hamming}, this says $((w_0;z,v),f) \in \mathcal{R}^{\delta}_{\mathrm{eval}}$ in the notation of \cref{sec:sound:rbr}.

\paragraph{The protocol as an IOR.}
Let $\Pieval^{\Sigma}$ be $\Pieval$ with instance $((z,v); y)$, and with messages and challenges encoded as follows.
\begin{itemize}
  \item For $j \in \iv{1}{\ell}$, the round-$j$ message is the string over $\Sigma$ formed by $(s_j(0), s_j(1))$ and $(s_j(2), 0)$, followed, if $j \ge 2$, by the oracle part $\mathrm{fib}_{j-1}(w_{j-1})$. The final message is $g$, written as $\lceil N_\ell/2 \rceil$ symbols, padded with a zero if $N_\ell = 1$. Padding entries are ignored.
  \item Fix an integer $\beta \ge \log_2 |\F|$ and a bijection $\mathrm{num} : \iv{0}{|\F|-1} \to \F$. The challenge of round $j \in \iv{1}{\ell}$ is $\mathsf{vr}_j \in \{0,1\}^\beta$, and $r_j = \phi(\mathsf{vr}_j)$, where $\phi(\mathsf{vr}) = \mathrm{num}(\mathrm{int}(\mathsf{vr}) \bmod |\F|)$ and $\mathrm{int}$ reads a bit string as a binary integer.
  \item Fix a bijection $\mathrm{pos} : \{0,1\}^{n+R} \to L$ (recall $M = 2^{n+R}$). The challenge of round $\ell+1$ is $\mathsf{vr}_{\ell+1} \in \{0,1\}^{\kappa(n+R)}$, cut into $\kappa$ blocks of $n+R$ bits, whose images under $\mathrm{pos}$ are the query points.
  \item On partial strings, the verifier rejects if an entry that it reads is $\bot$. Otherwise it runs the verifier of $\Pieval$ on the words $w_0 = \mathrm{fib}_0^{-1}(\bar y)$ and $w_{j-1} = \mathrm{fib}_{j-1}^{-1}(\overline{\text{oracle part of round } j})$, the round polynomials and the final table read from the filled messages; it outputs $(\top, ())$ if that verifier accepts, and rejects otherwise.
\end{itemize}
Every string of the prescribed length over $\Sigma$ encodes a message of $\Pieval$, and filled strings are such strings.
The query point $\xi_0 = \xi^{(i)}_0$ reads, for each level $i' \in \iv{1}{\ell}$, the entry at position $\xi_0^{2^{i'}}$ of the string $y$ (if $i' = 1$) or of the oracle part of round $i'$ (if $i' \ge 2$): this entry contains the two values $w_{i'-1}(\pm\xi_0^{2^{i'-1}})$ of the fold check of level $i'$, and, for $i' \ge 2$, one of them is the value $w_{i'-1}(\xi_0^{2^{i'-1}})$ compared at level $i'-1$.
The positions read are determined by the challenges.
So $\Pieval^{\Sigma}$ is an IOR from $\mathcal{R}_{\mathrm{eval}}$ to $\mathcal{R}_{\mathrm{acc}}$ with $\nu = \ell + 1$ rounds, $\beta_{\min} = \min(\beta, \kappa(n+R))$ and $l_{\max} = M/2$.

\begin{lemma}[Sampling]\label{lem:sampling}
For every $S \subseteq \F$, $\Pr_{\mathsf{vr} \leftarrow \{0,1\}^\beta}[\phi(\mathsf{vr}) \in S] \le |S|\,\bigl(1/|\F| + 2^{-\beta}\bigr)$.
The query points of round $\ell+1$ are independent and uniform in $L$.
\end{lemma}

\begin{proof}
Each $a \in \F$ has at most $\lceil 2^\beta/|\F| \rceil < 2^\beta/|\F| + 1$ preimages under $\phi$; divide by $2^\beta$ and sum over $S$.
The map $\{0,1\}^{\kappa(n+R)} \to L^{\times\kappa}$ induced by $\mathrm{pos}$ is a bijection.
\end{proof}

\paragraph{Erasures.}
For a statement and messages given as partial strings, let $\mathcal{U}_0 = \mathcal{U}(y) \subseteq L_1$ and, for $i \in \iv{1}{\ell-1}$, let $\mathcal{U}_i \subseteq L_{i+1}$ be the set of erased positions of the oracle part of round $i+1$.
We apply the notation of \cref{sec:sound:rbr} to the filled words $w_0, \dots, w_{\ell-1}$, round polynomials and final table defined above, and, for $j \in \iv{0}{\ell}$ and $c \in \Cc_j$, we put
\[
  \mathcal{W}^{\bot}_j(c) \;=\; \bigl\{\, \xi \in \mathcal{W}_j(c) \;:\; \xi^{2^{i}} \notin \mathcal{U}_{i-1} \text{ for all } i \in \iv{1}{\max(j,1)} \,\bigr\}, \qquad \mathrm{dens}^{\bot}_j(c) = |\mathcal{W}^{\bot}_j(c)|/M.
\]
For $\delta \in (0,\delta^*]$ and the statement at hand, let $\Dd^{\bot}_\delta$ be the set of partial transcripts obtained from \cref{def:state} with $\mathrm{dens}^{\bot}_j$ in place of $\mathrm{dens}_j$; it depends on the erasures of $y$ and of the messages, not only on the filled words.
Without erasures, $\mathcal{W}^{\bot}_j = \mathcal{W}_j$ and $\Dd^{\bot}_\delta$ is the doomed set of \cref{def:state}.

\begin{definition}[Knowledge state function of $\Pieval^{\Sigma}$]\label{def:kstate}
The algorithm $\mathsf{E}$, on every input, reads $y$, lets $w_0 = \mathrm{fib}_0^{-1}(\bar y)$, and outputs the decoded table of $w_0$ at radius $\delta^*$ (\cref{def:decoded} with $\delta = \delta^*$) if $\Delta^{\mathrm{fib}}_0(w_0, \Cc_0) \le \delta^*$, and the zero table otherwise.
For $\delta \in [0,1]$, a statement $((z,v), y)$, a transcript $\mathrm{tr}$ and a table $f$, let $\mathsf{KState}_\delta((z,v), y, \mathrm{tr}, f)$ be:
\begin{enumerate}
  \item $1$ if and only if $((z,v), y, f) \in (\mathcal{R}_{\mathrm{eval}})_{\le\delta}$, if $\mathrm{tr} = \mathrm{tr}_\emptyset$;
  \item $1$ if and only if the verifier of $\Pieval^{\Sigma}$ accepts, if $\mathrm{tr}$ is a full transcript;
  \item for $\mathrm{tr} = \tau_j$ the partial transcript after $j \in \iv{1}{\ell}$ rounds: $1$ if and only if $\tau_j \notin \Dd^{\bot}_\delta$, when $\delta \in (0, \delta^*]$; and $1$ when $\delta \notin (0, \delta^*]$;
  \item $\mathsf{KState}_\delta((z,v), y, \mathrm{tr}', f)$, if $\mathrm{tr} = \mathrm{tr}' \| \Pi$ ends with a prover message.
\end{enumerate}
\end{definition}

$\mathsf{E}$ runs in time polynomial in $M$ by \cref{rem:ext}, and it does not depend on $\delta$.
Conditions 1 and 3 of \cref{def:rrbr} hold by items 1 and 2, and condition 2 by item 4.

\begin{lemma}[Round-by-round steps with erasures]\label{lem:rbr-erasures}
Let $\delta \in (0, \delta^*]$, fix a statement given by a partial string $y$, and let $j \in \iv{1}{\ell+1}$.
Fix the partial transcript $\tau_{j-1}$ and the round-$j$ message (partial strings).
\begin{enumerate}
  \item If $j = 1$ and $((z,v), y, \mathsf{E}(y)) \notin (\mathcal{R}_{\mathrm{eval}})_{\le\delta}$, then at most $M_1 + 2$ values of $r_1$ make $\tau_1 \notin \Dd^{\bot}_\delta$.
  \item If $2 \le j \le \ell$ and $\tau_{j-1} \in \Dd^{\bot}_\delta$, then at most $M_j + 2$ values of $r_j$ make $\tau_j \notin \Dd^{\bot}_\delta$.
  \item If $j = \ell+1$ and $\tau_\ell \in \Dd^{\bot}_\delta$, then the verifier of $\Pieval^{\Sigma}$ accepts with probability at most $(1-\delta)^{\kappa}$ over the query points.
\end{enumerate}
\end{lemma}

\begin{proof}
We follow the proof of \cref{thm:rbr}, applied to the filled words, which are words of $\Pieval$.
First, \cref{lem:rbr-step} holds with $\mathcal{W}^{\bot}$: if $\xi \in \mathcal{W}^{\bot}_j(c')$, then $\xi \in \mathcal{W}_j(c')$, so $\xi \in \mathcal{W}_{j-1}(c)$ by \cref{lem:rbr-step}; the erasure conditions defining $\mathcal{W}^{\bot}_{j-1}(c)$, for $i \le \max(j-1,1) \le j$, are among those defining $\mathcal{W}^{\bot}_{j}(c')$. So $\mathcal{W}^{\bot}_j(c') \subseteq \mathcal{W}^{\bot}_{j-1}(c)$, and the size bound on $Y$ in that proof still holds because $\mathcal{W}^{\bot}_j(c') \subseteq \mathcal{W}_j(c')$.
\Cref{lem:rbr-sumcheck} does not involve witness sets.

(1) and (2). Call $r_j$ \emph{bad} as in the proof of \cref{thm:rbr}; at most $M_j + 2$ values are bad. As there, if $r_j$ is not bad and $\tau_j \notin \Dd^{\bot}_\delta$, then $c = \mathrm{dec}_{j-1}(w_{j-1})$ exists, $\mathrm{dens}^{\bot}_{j-1}(c) \ge 1-\delta$ and $v_{j-1} = \sigma_{j-1}(c)$, and the sumcheck checks of indices $1, \dots, j-1$ hold.
For $j \ge 2$ this says $\tau_{j-1} \notin \Dd^{\bot}_\delta$, contradicting the hypothesis of (2).
For $j = 1$: $\mathcal{W}^{\bot}_0(c)$ is the set of $\xi \in L$ with $\xi^2 \notin \mathcal{U}_0$ and $w_0 = c$ on $\{\pm\xi\}$, so its size is twice the number of fibres $\eta \in L_1$ with $y(\eta) = \mathrm{fib}_0(c)(\eta)$ (in particular $y(\eta) \ne \bot$). So $\mathrm{dens}^{\bot}_0(c) \ge 1-\delta$ gives $\Delta(y, \mathrm{fib}_0(c)) \le \delta$.
Hence $\Delta^{\mathrm{fib}}_0(w_0, c) = \Delta(\bar y, \mathrm{fib}_0(c)) \le \delta \le \delta^*$, so $\mathsf{E}(y) = \lambda[c]$ by \cref{lem:unique} at radius $\delta^*$ and \cref{lem:encoding}; and $v = v_0 = \sigma_0(c) = \widetilde{\lambda[c]}(z)$. Thus $((z,v), y, \mathsf{E}(y)) \in (\mathcal{R}_{\mathrm{eval}})_{\le\delta}$, contradicting the hypothesis of (1).

(3). If the verifier accepts, no entry it reads is $\bot$, its round polynomials and final table are those of the filled messages, and the checks of $\Pieval$ pass on the filled words. As in item 3 of \cref{thm:rbr}, the closure check gives $v_\ell = \sigma_\ell(w_\ell)$, so $\mathrm{dens}^{\bot}_\ell(w_\ell) < 1-\delta$ because $\tau_\ell \in \Dd^{\bot}_\delta$.
A query point $\xi$ that passes its checks lies in $\mathcal{W}_\ell(w_\ell)$, and it reads the positions $\xi^{2^{i}}$ of the strings with erasure sets $\mathcal{U}_{i-1}$, $i \in \iv{1}{\ell}$, which are therefore not erased; so $\xi \in \mathcal{W}^{\bot}_\ell(w_\ell)$.
By \cref{lem:sampling} and \cref{lem:queries}, the verifier accepts with probability at most $\mathrm{dens}^{\bot}_\ell(w_\ell)^\kappa < (1-\delta)^\kappa$.
\end{proof}

\begin{theorem}[Relaxed round-by-round knowledge soundness]\label{thm:rrbr}
$\Pieval^{\Sigma}$ has relaxed round-by-round knowledge soundness errors (\cref{def:rrbr}), with the knowledge state function and the algorithm $\mathsf{E}$ of \cref{def:kstate},
\[
  \varepsilon^{\mathrm{rr}}_j(\delta) \;=\; (M_j + 2)\Bigl(\frac{1}{|\F|} + \frac{1}{2^\beta}\Bigr) \quad (j \in \iv{1}{\ell}), \qquad \varepsilon^{\mathrm{rr}}_{\ell+1}(\delta) \;=\; (1-\delta)^{\kappa},
\]
for $\delta \in (0, \delta^*]$, and $\varepsilon^{\mathrm{rr}}_j(\delta) = 1$ for $\delta \notin (0, \delta^*]$.
\end{theorem}

\begin{proof}
For $\delta \notin (0,\delta^*]$ there is nothing to prove. Fix $\delta \in (0,\delta^*]$, a statement, a round $j \in \iv{1}{\ell+1}$ and $\mathrm{tr} = \tau_{j-1} \| \Pi_j$, where $\tau_0 = \mathrm{tr}_\emptyset$.
By item 4 of \cref{def:kstate}, $\mathsf{KState}_\delta(\cdot, \mathrm{tr}, \cdot) = \mathsf{KState}_\delta(\cdot, \tau_{j-1}, \cdot)$, and by items 2 and 3, $\mathsf{KState}_\delta(\cdot, \mathrm{tr}\|\mathsf{vr}_j, f)$ does not depend on $f$.
So the event of \cref{def:rrbr} is contained in the conjunction of $\mathsf{KState}_\delta(\cdot, \tau_{j-1}, \mathsf{E}(\cdot)) = 0$, which does not depend on $\mathsf{vr}_j$, and $\mathsf{KState}_\delta(\cdot, \tau_j, \cdot) = 1$.
If the first condition fails, the event is empty. Otherwise:
for $j = 1$, $((z,v), y, \mathsf{E}(y)) \notin (\mathcal{R}_{\mathrm{eval}})_{\le\delta}$; for $2 \le j \le \ell+1$, $\tau_{j-1} \in \Dd^{\bot}_\delta$.
For $j \le \ell$, \cref{lem:rbr-erasures}(1,2) and \cref{lem:sampling} with $|S| \le M_j + 2$ bound the probability that $\tau_j \notin \Dd^{\bot}_\delta$ by $\varepsilon^{\mathrm{rr}}_j(\delta)$.
For $j = \ell+1$, $\mathsf{KState}_\delta(\cdot, \tau_{\ell+1}, \cdot) = 1$ means acceptance, whose probability is at most $(1-\delta)^\kappa$ by \cref{lem:rbr-erasures}(3).
\end{proof}

\begin{corollary}[Post-quantum knowledge soundness]\label{cor:pq}
Let $\delta \in (0,\delta^*]$ and $\varepsilon_{\mathrm{rr}}(\delta) = \max\{(M_1+2)(1/|\F| + 2^{-\beta}),\ (1-\delta)^{\kappa}\}$.
The non-interactive scheme $\mathsf{BCS}[\Pieval^{\Sigma}]$ (\cref{sec:proofs:bcs}) satisfies the conclusion of \cref{thm:cdhz} with $\nu = \ell + 1$, $\beta_{\min} = \min(\beta, \kappa(n+R))$, $l_{\max} = M/2$ and this $\varepsilon_{\mathrm{rr}}(\delta)$.
In particular, for every quantum adversary that makes at most $Q$ queries to the random oracle and outputs a Merkle root $\mathrm{cm}$, a claim $(z, v)$ and a proof, except with probability $\varepsilon_{\mathrm{ext}}(Q)$, if the verifier accepts then the extractor outputs a partial string $y \in (\Sigma \sqcup \{\bot\})^{M/2}$, a table $f$ and a trapdoor such that:
the openings of all positions of $\mathrm{Dom}(y)$ are accepted under $\mathrm{cm}$; $y$ agrees with $\mathrm{fib}_0(\Enc(f))$ on at least $(1-\delta)M/2$ fibres; and $\tilde f(z) = v$.
\end{corollary}

\begin{proof}
\Cref{thm:rrbr} and \cref{thm:cdhz}; the last sentence unfolds \cref{def:committed} for $\mathcal{R}_{\mathrm{eval}}$ ($\bot$ entries count as disagreements).
\end{proof}

The extracted table is bound to the root: two extracted witnesses for the same root with different tables expose a hash collision.

\begin{lemma}[Uniqueness of relaxed committed witnesses]\label{lem:cr-unique}
Let $\delta \in (0,\delta^*]$ and let $\mathrm{cm}$ be a Merkle root. Let $\mathtt{cx} = ((z,v), \mathrm{cm})$ and $\mathtt{cx}' = ((z',v'), \mathrm{cm})$, and suppose that $(\mathtt{cx}, (y, f, \mathrm{td}))$ and $(\mathtt{cx}', (y', f', \mathrm{td}'))$ belong to $\widehat{\mathsf{CR}}[\mathcal{R}_{\mathrm{eval}}, \delta]$.
If $f \ne f'$, then $\mathrm{td}$ and $\mathrm{td}'$ yield two accepting openings, under $\mathrm{cm}$, of one position to different values, hence a collision of the random oracle.
In particular, if $z = z'$ and $v \ne v'$, such a collision is obtained.
\end{lemma}

\begin{proof}
Let $A = \{\eta \in L_1 : y(\eta) = \mathrm{fib}_0(\Enc(f))(\eta)\}$ and $A' = \{\eta \in L_1 : y'(\eta) = \mathrm{fib}_0(\Enc(f'))(\eta)\}$.
Then $A \subseteq \mathrm{Dom}(y)$ and $A' \subseteq \mathrm{Dom}(y')$, and $|A|, |A'| \ge (1-\delta)M/2$, so $|A \cap A'| \ge (1-2\delta)M/2 \ge \rho M/2 = N/2$.
If $y = y'$ on $A \cap A'$, then $\Enc(f)$ and $\Enc(f')$ agree on the at least $N$ points of the fibres over $A \cap A'$, hence are equal (\cref{lem:rs-params}(2)), and $f = f'$ (\cref{lem:encoding}).
So if $f \ne f'$, some $\eta \in A \cap A'$ has $y(\eta) \ne y'(\eta)$, and the openings of $\eta$ computed from $\mathrm{td}$ and $\mathrm{td}'$ are both accepted under $\mathrm{cm}$ (\cref{def:committed}), with the same random oracle since the commitment index is fixed by the scheme; apply the salted form of \cref{lem:merkle}(1) (\cref{sec:proofs:bcs}).
If $z = z'$ and $v \ne v'$, then $\tilde f(z) = v \ne v' = \tilde f'(z)$, so $f \ne f'$.
\end{proof}

\Cref{lem:cr-unique} is deterministic. We do not state a probability bound for an adversary that produces two accepting proofs for one root, since that would require running the extractor of \cref{thm:cdhz} on two proofs of one adversary, which the theorem does not address.

\begin{remark}[Post-quantum parameters]\label{rem:pq-params}
Take $\rho = 1/4$, $\delta = \delta^* = 3/8$, $M \le 2^{32}$ (so $\ell \le n \le 30$), the quartic extension $\F = \mathbb{F}_p[\iota]/(\iota^4 - 7)$ of the Goldilocks field, of size $p^4 > 2^{255.99}$ ($\iota^4 - 7$ is irreducible over $\mathbb{F}_p$ because $7$ is a non-square in $\mathbb{F}_p$ and $p \equiv 1 \pmod 4$~\cite[Theorem~3.75]{LN97}), $\beta = 320$, hash output length $\lambda_{\mathsf{H}} = 256$, and $\kappa = 248$ queries, so that $(5/8)^{\kappa} < 2^{-168}$ and $(M_1+2)(1/|\F| + 2^{-\beta}) < 2^{-224}$.
The quantities of \cref{thm:cdhz} satisfy $q_s \le \kappa(\ell+1) + 2\ell + \lceil N_\ell/2 \rceil \le 2^{30}$ (positions read: one per query and level, the sumcheck symbols, and the final table), $q_{\mathsf{V}} \le \nu + q_s(\lceil\log_2 l_{\max}\rceil + 1) \le 2^{36}$ (one challenge query per round and one authentication path, of at most $\lceil\log_2 l_{\max}\rceil + 1$ hashes, per position read), and $q_{\widehat{\mathsf{CR}}} \le (M/2)\log_2 M \le 2^{36}$ (one leaf hash and $\log_2(M/2)$ internal hashes per position of $y$).
For an adversary making $Q = 2^{64}$ queries, \cref{thm:cdhz} then gives $\varepsilon_{\mathrm{ext}}(2^{64}) < 2^{-31.8}$: the first term is $2^{-31.84}$ and the other terms together are below $2^{-51}$ (computed exactly in rational arithmetic by \texttt{rust/examples/pq\_params.rs} and \texttt{checks/pq\_params.py} in the companion repository).
The dominant term is $320\, Q^2 \varepsilon_{\mathrm{rr}}(\delta)$, so each additional factor $2$ in $Q$ costs about $2/\log_2(8/5) \approx 3$ further queries.
The parameters of \cref{sec:exp} (quadratic extension, $\kappa = 148$) give $\varepsilon_{\mathrm{rr}}(\delta) \approx 2^{-100}$ and are sized for classical adversaries only.
\end{remark}

\begin{remark}[Parameters]\label{rem:params}
For $\rho = 1/4$ and $\delta = (1-\rho)/2 = 3/8$, one query passes with probability at most $5/8$; $\kappa = 148$ queries give $(5/8)^{\kappa} < 2^{-100}$.
With $M = 2^{24}$ and $|\F| = p^2 > 2^{127.9}$ (the quadratic extension of \cref{sec:impl}), $\varepsilon_{\mathrm{fold}} + 2\ell/|\F| < 2^{-103}$.
These unique-decoding parameters are conservative; list-decoding analyses of FRI-type folding~\cite{BCIKS23} allow fewer queries, and we leave that extension to future work.
\end{remark}

\begin{remark}[Scope]\label{rem:fs}
\begin{itemize}
  \item The analysis concerns one commitment opened at one point $z$ per proof. Reusing a commitment across several openings, batching, or choosing $z$ after the commitment inside a larger Fiat--Shamir transcript is not analysed here.
  \item \Cref{cor:pq} applies to the compiler $\mathsf{BCS}$ of~\cite[Construction~11.7]{CDHZ25}, with salted Merkle trees and the challenge derivation described in \cref{sec:proofs:bcs}. The implementation of \cref{sec:impl} uses unsalted, domain-separated Merkle trees and a hash-chain transcript, and is therefore not covered by \cref{cor:pq} as it stands.
  \item For $\ell = 0$ the same argument applies with $\Sigma = \F$, the implicit instance $y = w_0$ and the Hamming distance, giving the single error $(1-\delta)^{\kappa}$: the verifier rejects on $\bot$, $\mathsf{E}$ decodes the filled word at radius $\delta^*$, and, for a fixed final table $g$, the accepted query points lie in $\{\xi : y(\xi) \ne \bot,\ y(\xi) = G(\xi)\}$, whose density is less than $1-\delta$ unless $\Delta(y, \ev_L(G)) \le \delta$ and $v = \tilde g(z)$, in which case $\mathsf{E}(y) = g$ is a witness (as in \cref{prop:l0}).
  \item The bounds also hold against classical adversaries in the random oracle model, with the same (quantum) extractor (\cref{sec:proofs:bcs}).
\end{itemize}
\end{remark}

\section{Implementation}\label{sec:impl}

We implemented $\Pieval$ of \cref{sec:pcs}, compiled with Merkle trees and Fiat--Shamir (\cref{sec:proofs:bcs}), as a self-contained Rust library, \texttt{kbfold} (about $1{,}200$ lines including tests and benchmarks; its only dependency is the BLAKE3 hash function).
The code follows the paper line by line; this section records the concrete choices, and \cref{sec:exp} reports measurements.
The library also has a coefficient-form mode, used only as the baseline of \cref{sec:exp}.

\paragraph{Index conventions.}
The paper numbers bits, variables and rounds from $1$; the code, like most Rust code, numbers them from $0$.
Bit $b_k$ of the paper is bit $k-1$ of the code (\texttt{(b >> (k-1)) \& 1}), the variable $x_k$ and the coordinate $z_k$ are entries $k-1$ of the corresponding arrays, and the round polynomial $s_j$ and the challenge $r_j$ are stored at index $j-1$, while the oracle $w_j$ (sent in round $j+1$) is stored at index $j$.
Tables $f : \iv{0}{N-1} \to \Fq$ are arrays of length $N$ indexed identically in both.

\subsection{Fields}\label{sec:impl:fields}

\paragraph{Base field.}
Tables and the initial codeword $w_0$ live over the Goldilocks field $\mathbb{F}_p$, $p = 2^{64} - 2^{32} + 1$.
The largest power of $2$ dividing $p - 1$ is $2^{32}$, so smooth domains of order up to $2^{32}$ exist (\cref{lem:power}).
We take $L = \langle \omega \rangle$ with $\omega = 7^{(p-1)/M}$, where $7$ generates $\mathbb{F}_p^{\times}$: since $p - 1 = 2^{32}\cdot 3 \cdot 5 \cdot 17 \cdot 257 \cdot 65537$, this amounts to $7^{(p-1)/\ell'} \ne 1$ for each of the six prime factors $\ell'$, which the test suite checks.
Multiplication uses the standard reduction of a $128$-bit product modulo $p$ based on $2^{64} \equiv 2^{32}-1$ and $2^{96} \equiv -1$; elements are kept in canonical form.

\paragraph{Extension field.}
The challenge field of \cref{sec:pcs:params} is the quadratic extension $\F = \K = \mathbb{F}_p[\iota]/(\iota^2 - 7)$, of size $p^2 \approx 2^{128}$: the challenges $r_j$, the evaluation point $z$, the value $v$, the sumcheck tables and all folded words $w_j$ ($j \ge 1$) live in $\K$, while tables and $w_0$ have entries in $\mathbb{F}_p$.
The polynomial $\iota^2 - 7$ is irreducible because $7$, a generator of $\mathbb{F}_p^\times$, is a non-square.
The post-quantum parameters of \cref{rem:pq-params} use the quartic extension instead; the code does not implement it yet, and its Merkle trees and transcript differ from the compiler of~\cite{CDHZ25} (\cref{rem:fs}).
Multiplication in $\K$ uses three base-field multiplications (Karatsuba) and one multiplication by the constant $7$.
Mixing base-field words with extension-field challenges is only needed in the first fold, which the code handles generically.

\subsection{Encoding and folding}\label{sec:impl:enc}

\paragraph{Commit.}
$\Enc(f)$ (\cref{def:encoding}) is computed in three steps:
the kernel-to-monomial transform $\mathrm{coef}(U_f) = C^{\otimes n} f$ of \cref{cor:transforms}, applied in place one tensor factor at a time as $(x, x') \mapsto (x', x + x')$ (exactly $\tfrac{n}{2}N$ additions, no multiplications);
zero-padding to length $M$;
and an iterative radix-$2$ NTT producing $w_0(\omega^i)$ in natural order.

\paragraph{Domain layout.}
With natural ordering, $L_j = \langle \omega^{2^j} \rangle$ is stored as $(\omega^{2^j i})_{i \in \iv{0}{M_j-1}}$.
The fibre of index $i \in \iv{0}{M_{j+1}-1}$ is the pair of positions $(i, i + M_{j+1})$, since $\omega^{2^j M_{j+1}} = -1$, and its image under squaring is position $i$ of $L_{j+1}$.
Folding therefore reads the two halves of the array and writes one array of half size, with no permutation.

\paragraph{Fold.}
Both prover and verifier evaluate \eqref{eq:fold-explicit} in the line form of \cref{lem:line}.
With $a = w(\xi)$ and $a' = w(-\xi)$, they compute $a_e = w_e(\xi^2) = (a+a')/2$ and $a_o = w_o(\xi^2) = (a-a')/(2\xi)$, and then $\fold_r(w)(\xi^2) = (a_o - a_e) + r\,(2a_e - a_o)$ (\cref{lem:line}).
This costs one multiplication by $r$ in $\K$ and two multiplications by base-field constants.
The inverses $(2\xi)^{-1} = 2^{-1}\omega^{-2^j i}$ are read from a single table of $\omega^{-i}$, $i < M/2$, shared by all levels.
The same function is used by the prover on whole words and by the verifier on single fibres, which rules out the parametrisation mismatch of \cref{rem:pitfall} by construction.

\paragraph{Sumcheck.}
The prover keeps the tables $f_{j-1}$ and $\chi_{j-1}$ of \cref{sec:pcs:protocol} as arrays over $\K$.
It sends $s_j(0)$, $s_j(1)$ and $s_j(2)$, computed in one pass over the pairs of entries at positions $2c$ and $2c+1$ of both tables, using
\[
  s_j(2) \;=\; \sum_{c} \bigl(2f_{j-1}(2c+1) - f_{j-1}(2c)\bigr)\bigl(2\chi_{j-1}(2c+1) - \chi_{j-1}(2c)\bigr).
\]
Both tables are then restricted in place (\cref{def:restriction}).
The verifier recovers $s_j(r_j)$ by Lagrange interpolation on $\{0,1,2\}$.

\subsection{Commitments and transcript}\label{sec:impl:merkle}

\paragraph{Fibre leaves.}
Every query at level $j$ reads a whole fibre $\{w_j(\xi), w_j(-\xi)\}$ (\cref{sec:pcs:protocol}).
We therefore build each Merkle tree over the $M_{j+1}$ fibres of $w_j$, one leaf per fibre containing both values.
A query at level $j$ costs one authentication path of length $\log_2 M_{j+1}$.
The value $w_{j+1}(\xi_{j+1})$ needed for the fold check at level $j+1$ is one of the two entries of the fibre opened at level $j+1$, so each query opens exactly one leaf per level.
Leaves and internal nodes are hashed with BLAKE3 under distinct one-byte domain separators, as in \cref{sec:proofs:bcs}.
In the IOP view, this means that the oracle $w_j$ is presented as the word $\mathrm{fib}_j(w_j)$ over the alphabet $\F^2$ indexed by $L_{j+1}$, one pair per fibre (\cref{sec:sound:pq}); every query of $\Pieval$ reads whole fibres, so this changes neither the protocol nor the analysis of \cref{sec:sound}.

\paragraph{Fiat--Shamir.}
The transcript is a BLAKE3 hash chain.
It absorbs the parameters $(n, R, \ell, \kappa)$, the root of $w_0$, the point $z$ and the value $v$.
Then, for $j = 1,\dots,\ell$, it absorbs $s_j$ before deriving $r_j$, and (for $j < \ell$) the root of $w_j$ before deriving $r_{j+1}$.
It then absorbs the final table $g$, and finally derives the $\kappa$ query indices.
A challenge in $\K$ is obtained from $256$ output bits as two $128$-bit integers reduced modulo $p$, at statistical distance below $2^{-64}$ from uniform per coordinate.

\paragraph{Final message.}
The verifier converts $g$ to monomial form with the transform $C^{\otimes (n-\ell)}$ and evaluates $G$ at each $\xi_\ell$ by Horner's rule.
It evaluates $\tilde g(z_{\ell+1},\dots,z_{n})$ for the closure check by $n-\ell$ successive restrictions (\cref{lem:restriction}).

\subsection{Testing}\label{sec:impl:tests}

The test suite checks each component against a direct definition (field arithmetic against $128$-bit integer arithmetic, the transforms against the expansion of \cref{cor:bool-coeff}, the NTT against Horner evaluation) and the protocol end to end:
\begin{itemize}
  \item \emph{Completeness.} Honest proofs verify for all $1 \le n \le 10$, all $0 \le \ell \le n$ and rates $\rho \in \{1/2, 1/4, 1/8\}$, and the returned value equals $\tilde f(z)$ computed independently.
  \item \emph{Tampering.} Proofs are rejected after changing the claimed value, the point, the commitment, one sumcheck value, one entry of $g$, one opened value, or one intermediate root, and when a proof for a different table is checked against the original commitment.
  \item \emph{A consistent cheater.} Consider a prover that commits honestly and claims $v+1$. In every round it adds half of the current discrepancy to each of $s_j(0)$, $s_j(1)$ and $s_j(2)$, so that every sumcheck check passes; it folds honestly, and modifies $g(0)$ so that the closure check passes. This prover is rejected in every tested configuration, and always in the query phase, as predicted by case~2 of \cref{thm:sound}.
\end{itemize}
Independently, each algebraic statement of \cref{sec:kernel,sec:fold} and the two lemmas of \cref{sec:sound:lemmas} were checked numerically over small prime fields ($\mathbb{F}_{17}$ exhaustively where feasible, and $\mathbb{F}_{257}$); these scripts are part of the artefact.

\section{Experiments}\label{sec:exp}

This section measures the implementation of \cref{sec:impl}.
The goals are modest and specific: to check that the kernel encoding costs the same as coefficient-form folding at identical engineering (as predicted by \cref{rem:basefold}), to locate where the time goes, and to quantify the two protocol parameters $\ell$ and $\rho$.
We do not claim competitive absolute performance: the code is single-threaded and uses no SIMD, no Merkle multiproofs and no hash-specific batching.

\subsection{Setup}\label{sec:exp:setup}

\paragraph{Machine.}
All runs use one core of an Intel Xeon processor at $2.10$\,GHz (cloud virtual machine, $2$ vCPUs, $7$\,GB RAM), Rust $1.95.0$, release profile with fat LTO and one codegen unit, and no \texttt{target-cpu} flag.
Times are medians over $11$ runs for $n \le 18$, $5$ for $n = 20$ and $3$ for $n = 22$ and for \cref{tab:stop,tab:rate}; verification times are medians over $21$ runs.
The machine is a shared virtual machine, and timings of identical work vary by up to about $20\%$ between runs (see \cref{sec:exp:scaling}).

\paragraph{Parameters.}
Unless stated otherwise: rate $\rho = 1/4$ ($R = 2$), number of folding rounds $\ell = n - 4$ (final table of $16$ elements of $\K$), and $\kappa = 148$ queries.
By \cref{thm:sound} with $\delta = (1-\rho)/2 = 3/8$, the query term is $(5/8)^{148} < 2^{-100}$, and $\varepsilon_{\mathrm{fold}} + 2\ell/|\K| < 2^{-100}$ for all sizes tested ($M \le 2^{24}$, $|\K| \approx 2^{128}$).
Tables $f$ are uniformly random over $\mathbb{F}_p$ and the point $z$ is uniformly random in $\K^n$.

\paragraph{Baseline.}
To isolate the effect of the encoding, the library has a second mode that differs from $\Pieval$ in exactly two places: the table is converted to the monomial coefficients of $\tilde f$ by M\"obius inversion (\cref{lem:mobius}) instead of the transform of \cref{cor:transforms}, and words are folded with the classical FRI fold $\mathrm{cfold}_\theta(w) = w_e + \theta\,w_o$ (\cref{def:cfold}) instead of the kernel fold (\cref{rem:coeff-form}).
This is the coefficient-form folding used in Reed--Solomon instantiations of BaseFold and in DeepFold; everything else (field, NTT, Merkle trees, transcript, sumcheck, query phase) is shared.
The baseline passes the same completeness and tampering tests.
It is \emph{not} the BaseFold or DeepFold code base, and \cref{tab:scaling} should be read as a comparison of encodings, not of implementations.

\subsection{Kernel versus coefficient-form encoding}\label{sec:exp:scaling}

\begin{table}[t]
\centering
\small
\begin{tabular}{r rr rr rr r}
\toprule
 & \multicolumn{2}{c}{Commit (ms)} & \multicolumn{2}{c}{Open (ms)} & \multicolumn{2}{c}{Verify (ms)} & Proof \\
$n$ & kernel & coeff. & kernel & coeff. & kernel & coeff. & (KiB) \\
\midrule
12 & 2.5 & 2.3 & 2.6 & 2.6 & 2.3 & 2.3 & 387 \\
14 & 10.2 & 9.8 & 10.4 & 10.3 & 3.1 & 3.1 & 531 \\
16 & 44.7 & 36.5 & 42.1 & 34.8 & 3.1 & 3.5 & 693 \\
18 & 158.8 & 176.4 & 173.6 & 171.7 & 5.0 & 4.1 & 873 \\
20 & 821.5 & 779.0 & 734.8 & 751.8 & 6.1 & 6.1 & 1072 \\
22 & 4139.7 & 3998.1 & 3354.9 & 3401.3 & 6.1 & 7.2 & 1290 \\
\bottomrule
\end{tabular}
\caption{Kernel encoding ($\Pieval$) versus the coefficient-form baseline, single-threaded, $\rho = 1/4$, $\ell = n-4$, $\kappa = 148$. Proof sizes are identical for both encodings.}
\label{tab:scaling}
\end{table}

\Cref{tab:scaling} reports commit, open and verify times for $n = 12,\dots,22$ ($N = 2^{12},\dots,2^{22}$).
Proof sizes are identical because the proofs have the same shape.
The timing differences between the two encodings are of the same size as the differences in steps whose code is identical in both modes: in \cref{tab:breakdown}, the NTT and Merkle columns differ by up to about $20\%$ between modes although they perform exactly the same computation.
Moreover, in an earlier full run of the same benchmark, the sign of the difference in commit time reversed at $n = 18$ and $n = 20$.
We conclude that the two encodings are indistinguishable at the resolution of this machine.
This is the outcome predicted by \cref{rem:basefold}: the kernel transform and M\"obius inversion both cost $\tfrac{n}{2}N$ additions, and in line form (\cref{lem:line}) the kernel fold costs one multiplication by $r$ per output element, the same as the classical fold.
The kernel encoding therefore buys its structural property (\cref{thm:fold-dictionary}) at no measurable cost.

\subsection{Where the time goes}\label{sec:exp:breakdown}

\begin{table}[t]
\centering
\small
\begin{tabular}{r rr rr rr}
\toprule
 & \multicolumn{2}{c}{Transform (ms)} & \multicolumn{2}{c}{NTT (ms)} & \multicolumn{2}{c}{Merkle tree (ms)} \\
$n$ & kernel & coeff. & kernel & coeff. & kernel & coeff. \\
\midrule
16 & 0.48 & 0.32 & 9.4 & 7.8 & 32.4 & 28.4 \\
18 & 2.68 & 2.61 & 40.6 & 46.3 & 104.6 & 128.6 \\
20 & 15.2 & 13.6 & 210.2 & 227.5 & 521.2 & 509.1 \\
22 & 60.8 & 66.1 & 1476.1 & 1469.5 & 2391.2 & 2271.7 \\
\bottomrule
\end{tabular}
\caption{Components of the commitment ($M = 4N$), measured separately in the same runs as \cref{tab:scaling}. The NTT and Merkle columns execute identical code in both modes.}
\label{tab:breakdown}
\end{table}

\Cref{tab:breakdown} splits the commitment into its three steps.
The kernel transform accounts for $1$--$2\%$ of the commitment at every size, and takes the same time as M\"obius inversion (it performs $\tfrac{n}{2}N$ additions, against $\tfrac{M}{2}\log_2 M$ multiplications for the NTT on the $4\times$ larger domain).
The commitment is dominated by hashing: BLAKE3 over $M/2$ fibre leaves and $M/2 - 1$ internal nodes takes $1.6$--$3.5$ times as long as the NTT.
Opening costs about as much as committing: the folded words $w_1,\dots,w_{\ell-1}$ have total length below $M$ but live in $\K$, and each needs its own Merkle tree; the sumcheck ($O(N)$ multiplications in $\K$) is a minor term.

Verification takes $2$--$7$\,ms and grows slowly with $n$, as expected from its $O(\kappa \ell \log M)$ hashing cost.
Proof sizes ($0.4$--$1.3$\,MiB) are dominated by the $\kappa \ell$ authentication paths: we open one Merkle leaf per query and level and do not merge paths across queries.
Both the unique-decoding query count and the absence of path merging are deliberate simplifications; \cref{rem:params} discusses the former.

\subsection{The number of folding rounds}\label{sec:exp:stop}

\begin{table}[t]
\centering
\small
\begin{tabular}{r r rrr}
\toprule
$\ell$ & final table $N_\ell$ & open (ms) & verify (ms) & proof (KiB) \\
\midrule
8 & 4096 & 769.7 & 9.1 & 747 \\
10 & 1024 & 748.8 & 5.8 & 824 \\
12 & 256 & 746.9 & 4.5 & 918 \\
14 & 64 & 726.8 & 5.6 & 1003 \\
16 & 16 & 666.2 & 6.0 & 1072 \\
18 & 4 & 744.9 & 6.4 & 1123 \\
20 & 1 & 678.9 & 5.2 & 1156 \\
\bottomrule
\end{tabular}
\caption{Effect of the number of folding rounds $\ell$ (\cref{rem:stop}), kernel encoding, $n = 20$, $\rho = 1/4$, $\kappa = 148$.}
\label{tab:stop}
\end{table}

\Cref{tab:stop} varies $\ell$ at $n = 20$.
Each folding round adds $\kappa$ authentication paths to the proof, while stopping earlier adds $16 N_\ell$ bytes for the final table; in our range the paths dominate, so proofs shrink as $\ell$ decreases, from $1156$\,KiB at $\ell = n$ to $747$\,KiB at $\ell = n-12$.
Verification time is smallest around $\ell = 12$: for smaller $\ell$ the verifier's $\kappa$ evaluations of $G$, each of cost $O(N_\ell)$, take over.
Opening time is essentially independent of $\ell$, since the first folds dominate it.

\subsection{The rate}\label{sec:exp:rate}

\begin{table}[t]
\centering
\small
\begin{tabular}{r r r rrrr}
\toprule
$\rho$ & $\delta$ & $\kappa$ & commit (ms) & open (ms) & verify (ms) & proof (KiB) \\
\midrule
$1/2$ & $1/4$ & 241 & 422.2 & 405.3 & 7.9 & 1625 \\
$1/4$ & $3/8$ & 148 & 704.5 & 765.5 & 6.1 & 1072 \\
$1/8$ & $7/16$ & 121 & 1880.5 & 1521.2 & 5.3 & 937 \\
\bottomrule
\end{tabular}
\caption{Effect of the rate at $100$ bits of query security, $\delta = (1-\rho)/2$, kernel encoding, $n = 20$, $\ell = 16$.}
\label{tab:rate}
\end{table}

\Cref{tab:rate} fixes the query term of \cref{thm:sound} at $2^{-100}$ and varies the rate.
Halving the rate doubles the domain and increases the prover's work by a factor of $1.7$--$2.7$; it also reduces the number of queries, which shortens proofs and speeds up verification.
Within the unique-decoding analysis of this paper, $\rho = 1/4$ is a reasonable default.

\section{Conclusion and Open Problems}\label{sec:concl}

The Boolean-kernel basis $\{K_b\}$ gives an exact dictionary between two operations that are usually treated separately.
One is even--odd folding of univariate polynomials over smooth multiplicative domains; the other is restriction of a variable of a multilinear polynomial in table form.
In kernel coordinates the normalised fold $\mathrm{pfold}_r$ \emph{is} the restriction (\cref{thm:fold-dictionary}).
A Reed--Solomon codeword can therefore carry a table directly, and the BaseFold paradigm runs on it with no change of representation.
We gave a self-contained proof of soundness, evaluation binding and round-by-round knowledge soundness in the unique-decoding regime, whose only external input about codes is the correlated-agreement theorem for Reed--Solomon codes (with efficient unique decoding for the extractor).
With the BCS theorem for interactive oracle reductions of~\cite{CDHZ25}, this makes the non-interactive scheme knowledge sound, for a single opening, against quantum adversaries in the quantum random oracle model (\cref{cor:pq}).
We also gave an implementation, in which the kernel encoding and coefficient-form folding are indistinguishable in running time at the resolution of our machine.

Several questions remain open.
\begin{itemize}
  \item \emph{List decoding.} Extend \cref{thm:sound} to proximity parameters up to the Johnson bound, in the spirit of~\cite{Hab24,Zei26}. The line form of \cref{lem:line} reduces each fold to a Reed--Solomon line, so correlated agreement beyond unique decoding~\cite{BCIKS23} is the natural tool. The main work is to replace the unique codewords $c_j$ of the passing-set argument by lists.
  \item \emph{Several openings.} Extend \cref{cor:pq} to several openings of one commitment, and to batched openings (\cref{rem:batch}).
  \item \emph{Conformance of the implementation.} Implement the compiler of~\cite[Construction~11.7]{CDHZ25} (salted Merkle trees, their challenge derivation) and the quartic extension, so that \cref{cor:pq} applies to the code.
  \item \emph{Tighter post-quantum parameters.} The quadratic loss in the number of queries in \cref{thm:cdhz} makes the post-quantum parameters of \cref{rem:pq-params} conservative; a list-decoding version of \cref{thm:rbr} would reduce the number of queries.
  \item \emph{Other domains.} The dictionary uses a smooth multiplicative subgroup in odd characteristic. Is there a kernel-type basis whose natural fold is table-form restriction on circle domains~\cite{HLP24} or on additive subspaces~\cite{LCH14,DP24}?
  \item \emph{Engineering.} Merkle multiproofs, batched openings (\cref{rem:batch}), parallelism and SIMD field arithmetic would bring proof sizes and running times in line with production systems. None of them interacts with the encoding.
  \item \emph{Zero knowledge.} Adapt standard masking techniques for folding-based schemes to the kernel encoding.
\end{itemize}

\paragraph{Code availability.}
The Rust implementation of \cref{sec:impl}, the benchmark scripts of \cref{sec:exp} and the numerical checks of \cref{sec:impl:tests} are released as open source together with this paper.


\bibliographystyle{alpha}
\bibliography{refs}

\appendix
\section{Notation}\label{app:notation}

\begin{center}
\small
\begin{longtable}{@{}p{0.27\textwidth} p{0.49\textwidth} >{\raggedright\arraybackslash}p{0.17\textwidth}@{}}
\toprule
Symbol & Meaning & Where \\
\midrule
\endhead
\multicolumn{3}{l}{\emph{Integers, bits and points}} \\
$\N$, $\iv{a}{c}$ & non-negative integers; integers $i$ with $a \le i \le c$ (empty if $a > c$) & \S\ref{sec:prelim:bits} \\
$n$, $N = 2^n$ & number of variables; size of a table & \S\ref{sec:prelim:bits} \\
$b_k$, $(b_1,\dots,b_n)$ & bit $k$ of an index $b$ (least significant first); bit vector & \S\ref{sec:prelim:bits} \\
$B_n$ & Boolean hypercube $\{0,1\}^n$, identified with $\iv{0}{N-1}$ & \S\ref{sec:prelim:bits} \\
$\wt(b)$, $\cmp{b}$, $c \sub b$ & Hamming weight; complement; submask & \S\ref{sec:prelim:bits} \\
$b = c + 2^t h$, $b = b_1 + 2b'$ & splitting an index & \S\ref{sec:prelim:bits} \\
$(x,y)$, $(x,b)$, $r_{\le j}$ & concatenated points; prefix of a point & \S\ref{sec:prelim:bits} \\
\midrule
\multicolumn{3}{l}{\emph{Fields and polynomials}} \\
$\Fq$ & base field, finite, odd characteristic & \S\ref{sec:prelim:fields} \\
$\F$, $\K$ & finite field containing $\Fq$ (challenge field); extension field & \S\ref{sec:prelim:fields}, \S\ref{sec:pcs:params} \\
$\F[X]_{<d}$, $\coef{a}{P}$ & polynomials of degree $< d$; coefficient of $X^a$ & \S\ref{sec:prelim:fields} \\
$P_e$, $P_o$; $R_i$ & even and odd parts; blocks of a polynomial & \Cref{lem:splits} \\
\midrule
\multicolumn{3}{l}{\emph{Multilinear polynomials and tables}} \\
$A$ & commutative ring & \S\ref{sec:prelim:ml} \\
$\mathcal{L}_n(A)$, $\mathcal{L}_n$ & multilinear polynomials over $A$, over $\Fq$ & \S\ref{sec:prelim:ml} \\
$m_a$, $e_b$ & canonical monomial; Lagrange polynomial & \Cref{def:bases} \\
$\eq_1$, $\eq$ & equality polynomials; $e_b(x) = \eq(b,x)$ & \Cref{def:bases,lem:eq} \\
$f$, $f(b)$, $\tilde f$ & table; its entries; its multilinear extension & \Cref{prop:mle} \\
$\alpha_a$ & coefficient form of $\tilde f$ & \Cref{def:coeff-form} \\
$\res_r(f)$ & restriction of a table at $r \in A$ or $r \in A^j$ & \Cref{def:restriction} \\
$\vartheta$ & ring homomorphism, applied entrywise & \Cref{lem:natural} \\
\midrule
\multicolumn{3}{l}{\emph{Smooth domains and codes}} \\
$L$, $M = 2^{\mu}$ & smooth domain and its order & \Cref{def:domain} \\
$L^{2^j}$ & image of $L$ under $\xi \mapsto \xi^{2^j}$ (not a Cartesian power) & \Cref{def:domain,lem:power} \\
$\xi$, $\eta$, $\zeta$ & points of a domain; points of the squared domain; arbitrary field point & \S\ref{sec:prelim:rs}, \S\ref{sec:kernel:eval} \\
fibre $\{\xi, -\xi\}$ & the two square roots of $\eta = \xi^2$ in $L$ & \Cref{def:fibre} \\
$\ev_L$, $\RS_\F[L,d]$, $\rho$ & evaluation map; Reed--Solomon code; rate $d/M$ & \Cref{def:rs} \\
$P_c$ & polynomial of degree $< d$ whose evaluation is the codeword $c$ & \Cref{lem:rs-params} \\
$\mathrm{Frob}$ & Frobenius map $x \mapsto x^q$ & \Cref{lem:descent} \\
$\Delta$ & relative Hamming distance & \Cref{def:rs,lem:metric} \\
$w_e$, $w_o$ & even and odd parts of a word & \eqref{eq:even-odd}, \Cref{lem:split-words} \\
$\delta$ & proximity parameter, $0 < \delta \le (1-\rho)/2$ & \Cref{thm:bciks} \\
$u^0$, $u^1$; $\hat c^0$, $\hat c^1$ & words on a line; codewords of correlated agreement & \Cref{thm:bciks,lem:line} \\
\midrule
\multicolumn{3}{l}{\emph{Proof systems}} \\
$\mathcal{R}$, $\mathtt{x}$, $\mathtt{w}$, $L_{\mathcal{R}}$ & relation, instance, witness, language & \S\ref{sec:proofs:iop} \\
$\mathsf{P}$, $\mathsf{V}$, $\mathsf{P}^*$ & prover, verifier, arbitrary prover & \Cref{def:iop} \\
$\nu$, $\mathsf{m}_i$, $\mathsf{c}_i$, $\mathsf{Ch}_i$ & number of rounds; message; challenge; challenge set & \Cref{def:iop} \\
$\tau$, $\tau_i$, $\Omega$ & partial transcript; after $i$ rounds; challenge space & \S\ref{sec:proofs:iop} \\
$E_i$, $k_i$ & events and counts of the sequential-challenge bound & \Cref{lem:sequential} \\
$\Omega'$, $\varpi$, $\mathcal{S}(\varpi)$, $\mathcal{A}$, $\mathbf{1}_{\mathcal{A}}$ & finite probability space, its element, a set of points, an event, its indicator & \Cref{lem:queries} \\
$L^{\times\kappa}$ & set of $\kappa$-tuples of elements of $L$ (Cartesian power) & \Cref{lem:queries}, \S\ref{sec:pcs:protocol} \\
$\Dd$, $\mathsf{Ext}$, $\varepsilon_i$, $\varepsilon_{\mathrm{rbr}}$ & doomed set; extractor; round errors; maximum error & \Cref{def:rbr,def:rbr-knowledge} \\
$F$, $s_j$, $s^*_j$, $r_j$, $v_j$ & summed polynomial; round polynomial sent; honest round polynomial; challenge; running claim & \S\ref{sec:proofs:sumcheck} \\
$\Enc$, $(w; z, v)$ & encoding; instance of an evaluation protocol & \Cref{def:pcs} \\
$\mathrm{dist}$, $\mathcal{R}^{\delta}_{\mathrm{eval}}$ & distance to encodings; evaluation relation & \S\ref{sec:proofs:pcs} \\
$\mathsf{H}$, $\lambda_{\mathsf{H}}$, $Q$ & hash function (random oracle); output length; number of oracle queries of the adversary & \S\ref{sec:proofs:bcs} \\
$\mathbf{d}$, $T$, $\mathrm{MT}(\mathbf{d})$ & vector of leaves; its length; Merkle root & \S\ref{sec:proofs:bcs}, \Cref{lem:merkle} \\
$\Sigma$, $y$, $\bot$, $\mathrm{Dom}(y)$ & alphabet; implicit instance (partial string); undefined entry; defined positions & \S\ref{sec:proofs:bcs} \\
$\mathrm{tr}$, $\mathrm{tr}_\emptyset$, $\mathbf{a}$, $\gamma_h$ & transcript; empty transcript; committed vector; salts & \S\ref{sec:proofs:bcs} \\
$\mathcal{R}_{\le\delta}$, $\mathcal{R}_{\mathrm{acc}}$, $(\top, ())$ & relaxed relation; acceptance relation; its instance (explicit $\top$, empty implicit instance $()$) & \S\ref{sec:proofs:bcs} \\
$n_y$, $n'_y$ & numbers of input and output implicit instances ($1$ and $0$ here) & \S\ref{sec:proofs:bcs} \\
$I$, $h$ & set of opened positions; position index in a salted Merkle tree & \S\ref{sec:proofs:bcs} \\
$\varepsilon_{\mathrm{sim}}$, $\varepsilon^{\mathrm{sr}}$ & simulation error; state-restoration knowledge error (\cite{CDHZ25}) & \S\ref{sec:proofs:bcs}, \Cref{thm:cdhz} \\
$\varepsilon_{\mathrm{extract}}$, $\varepsilon_{\mathrm{offline}}$, $\varepsilon_{\mathrm{online}}$ & multi-extractability errors of Merkle commitments (\cite[Thm.~8.3]{CDHZ25}) & proof of \Cref{thm:cdhz} \\
$\mathfrak{m}_{\mathsf{FS}}$, $\mathfrak{m}_{\mathsf{VC}}$, $\mathfrak{m}_{\mathsf{MultiExtract}}$, $\mathfrak{m}$, $\mathfrak{m}'$ & total query masses (\cite[Def.~4.3]{CDHZ25}) & proof of \Cref{thm:cdhz} \\
$\Pi_i$, $\mathsf{vr}_i$, $\beta_i$ & IOR message; challenge bit string; its length & \S\ref{sec:proofs:bcs} \\
$\mathsf{KState}_\delta$, $\varepsilon^{\mathrm{rr}}_i$, $\varepsilon_{\mathrm{rr}}$ & knowledge state function; relaxed round-by-round errors; their maximum & \Cref{def:rrbr,thm:cdhz} \\
$\mathrm{cm}$, $\mathrm{td}$, $\mathsf{CR}$, $\widehat{\mathsf{CR}}$ & Merkle root; trapdoor; committed and relaxed committed relations & \Cref{def:committed} \\
$\mathsf{BCS}[\Pi]$, $\varepsilon_{\mathrm{ext}}$, $\varepsilon_{\mathrm{com}}$ & compiled reduction; extraction and commutativity errors & \Cref{thm:cdhz} \\
$l_{\max}$, $q_{\mathsf{V}}$, $q_{\widehat{\mathsf{CR}}}$, $q_s$, $T_1$ & parameters of the bound of \cref{thm:cdhz} & \Cref{thm:cdhz} \\
\midrule
\multicolumn{3}{l}{\emph{Kernel basis and folding}} \\
$K_y$, $K_b$, $K^{(n)}_b$ & kernel polynomial; Boolean kernel polynomial; with $n$ variables & \Cref{def:kernel} \\
$S_1 \otimes \dots \otimes S_n$, $S_{(k)}$ & Kronecker product; factor on bit $k$ & \Cref{def:kron} \\
$C$, $\mathbf{K}_n$, $\mathrm{coef}(U)$ & $2\times 2$ change of basis; coefficient matrix of $(K_b)$; coefficient vector & \S\ref{sec:kernel:basis} \\
$\lambda$, $\lambda(b)$ & kernel coordinates & \Cref{def:kernel-coords} \\
$\varphi$ & kernel coordinates of a low-degree polynomial with $t$ variables & \Cref{thm:lowdeg} \\
$\gamma_k(\zeta)$, $\psi_k(\zeta)$ & scalar and curve of the evaluation identity & \Cref{prop:eval} \\
$U_e$, $U_o$ & even and odd parts of $U$ & \Cref{lem:splits,prop:split} \\
$\mathrm{pfold}_r$, $U_j$ & kernel fold of a polynomial; iterated folds & \Cref{def:fold-poly,cor:iterated} \\
$\fold_r$ & kernel fold of a word & \Cref{def:fold-word,lem:fold-word} \\
$\mathrm{cfold}_\theta$, $\theta$ & classical FRI fold and its challenge & \Cref{def:cfold} \\
\midrule
\multicolumn{3}{l}{\emph{The commitment scheme}} \\
$R$, $\ell$, $\kappa$ & $\rho = 2^{-R}$; number of folding rounds; number of queries & \S\ref{sec:pcs:params} \\
$L_j$, $M_j$, $N_j$, $\Cc_j$ & $L^{2^j}$; $M/2^j$; $N/2^j$; $\RS_\F[L_j,N_j]$ & \S\ref{sec:pcs:params} \\
$U_f$ & kernel polynomial of a table & \Cref{def:encoding,lem:encoding} \\
$z$, $v$ & evaluation point; claimed value & \S\ref{sec:pcs:protocol} \\
$f_j$, $\chi_j$ & honest restricted table; restricted equality table & \S\ref{sec:pcs:protocol}, \Cref{lem:honest} \\
$w_0$; $w_j$ & commitment (instance oracle); prover oracles for $1 \le j < \ell$, and $w_\ell = \ev_{L_\ell}(G)$ & \S\ref{sec:pcs:protocol} \\
$g$, $G$ & final table; its kernel polynomial & \S\ref{sec:pcs:protocol} \\
$\omega$, $\iota$ & generator of $L$; generator of the extension field & \S\ref{sec:impl} \\
$\xi^{(i)}_0$, $\xi_j$ & query points; $\xi_0^{2^j}$ & \S\ref{sec:pcs:protocol} \\
$\Pieval$, $t_{\mathsf{H}}$ & evaluation protocol; cost of one hash & \S\ref{sec:pcs:protocol}, \S\ref{sec:pcs:costs} \\
\midrule
\multicolumn{3}{l}{\emph{Soundness}} \\
$\Delta^{\mathrm{fib}}_j$ & fibre distance at level $j$ & \Cref{def:fibre-dist} \\
$\mathrm{dec}_j(w)$, $f^*$ & decoded codeword; decoded table & \Cref{def:decoded} \\
$\mathsf{Good}_j(w,r)$ & $r$ is good for $w$ at level $j$ & \Cref{def:good} \\
$\Pp_j$, $\pi$ & passing sets; passing fraction & \S\ref{sec:sound:thm}, \Cref{lem:passing} \\
$c_j$ & codeword chain & \Cref{lem:chain} \\
$\varepsilon_{\mathrm{fold}}$ & folding error $M(1-2^{-\ell})/|\F|$ & \Cref{thm:sound} \\
$E_j$, $f^*_j$, $\chi^*_j$, $\sigma^*_j$ & bad events; honest tables and claims for $f^*$ & proof of \Cref{thm:sound} \\
$f^H$ & table decoded within Hamming distance $\delta$ & \Cref{prop:l0} \\
$\hat w_j$, $\mathcal{B}_j$ & $\fold_{r_j}(w_{j-1})$; points where it differs from $w_j$ & \S\ref{sec:sound:rbr} \\
$\lambda[c]$ & kernel coordinates of $P_c$ & \S\ref{sec:sound:rbr} \\
$\mathcal{W}_j(c)$, $\mathrm{dens}_j(c)$, $\sigma_j(c)$ & witness set; its density; claimed value for a codeword $c$ & \S\ref{sec:sound:rbr} \\
$s^c_j$ & honest round polynomial for a codeword $c$ & \Cref{lem:rbr-sumcheck} \\
$\delta^*$ & largest proximity parameter $(1-\rho)/2$ & \S\ref{sec:sound:pq} \\
$\xi_\eta$, $\mathrm{fib}_j$ & chosen square root of $\eta$; fibre string of a word over $\Sigma = \F^2$ & \S\ref{sec:sound:pq} \\
$\mathcal{R}_{\mathrm{eval}}$, $\Pieval^{\Sigma}$ & evaluation relation with implicit instance; $\Pieval$ as an IOR & \S\ref{sec:sound:pq} \\
$\beta$, $\phi$, $\mathrm{num}$, $\mathrm{pos}$ & challenge length; sampling map; enumerations of $\F$ and $L$ & \S\ref{sec:sound:pq}, \Cref{lem:sampling} \\
$\mathsf{E}$ & $\delta$-independent extractor & \Cref{def:kstate} \\
$\bar u$, $\mathcal{U}(u)$, $\mathcal{U}_i$ & filling of a partial string; its erased positions; erasures of $y$ and of the oracles & \S\ref{sec:sound:pq} \\
$\mathcal{W}^{\bot}_j$, $\mathrm{dens}^{\bot}_j$, $\Dd^{\bot}_\delta$ & witness sets, densities and doomed set with erasures & \S\ref{sec:sound:pq} \\
\bottomrule
\end{longtable}
\end{center}

\section{Lean formalisation}\label{sec:lean}

The companion Lean~4 project \texttt{KBFold} (Lean~4.23.0, Mathlib at the tag \texttt{v4.23.0}; about $7000$ lines in $16$ modules) formalises the mathematics of \cref{sec:prelim,sec:proofs,sec:kernel,sec:fold,sec:pcs,sec:sound}.
It contains no \texttt{sorry}. Its only axiom is \texttt{bciks\_unique}, the correlated-agreement theorem (\cref{thm:bciks}); \texttt{\#print axioms} shows that completeness, the fold dictionary, the basis theorem and the probability lemmas do not use it, and that soundness, binding and (relaxed) round-by-round knowledge soundness use it and the three standard axioms of Lean (\texttt{propext}, \texttt{Classical.choice}, \texttt{Quot.sound}) only.

\paragraph{Modelling.}
\begin{itemize}
  \item Bit vectors are functions $\mathtt{Fin}\ n \to \mathtt{Bool}$; bit $k$ of Lean is bit $k+1$ of the paper, and tables are functions on bit vectors. The multilinear extension is its evaluation formula (\cref{prop:mle}).
  \item One field $\F$ is used, and the smooth domain is a subgroup $L \le \F^\times$ of order $2^\mu$; the paper's setting $L \le \Fq^\times \subseteq \F^\times$ is a special case.
  \item Probabilities are rational numbers computed by counting over finite uniform spaces; rounds and challenges are indexed from $0$.
  \item A deterministic prover of $\Pieval$ is a family of messages indexed by the challenges, with an explicit causality hypothesis (message $j$ depends on $r_1,\dots,r_{j-1}$ only) and the hypothesis that round polynomials have degree at most $2$, which the paper ensures by the transmission of three values.
  \item Decoding is defined non-constructively, by uniqueness; the running times of extractors (\cref{rem:ext}) and \cref{fact:decode} are not formalised.
\end{itemize}

\paragraph{Correspondence.}
\begin{center}
\small
\begin{longtable}{@{}p{0.34\linewidth}p{0.18\linewidth}p{0.42\linewidth}@{}}
\toprule
Paper & Lean module & Main declarations \\
\midrule
\endhead
\Cref{lem:eq,prop:mle,cor:eq-mle,def:restriction,lem:restriction,lem:natural} & \texttt{Multilinear} & \texttt{eqv\_pt}, \texttt{mle\_pt}, \texttt{mle\_eqTable}, \texttt{mle\_res}, \texttt{map\_mle} \\
\Cref{lem:mobius} & \texttt{Mobius} & \texttt{mle\_eq\_sum\_mono}, \texttt{table\_coeff\_form} \\
\Cref{def:domain,lem:power,def:fibre,def:rs,lem:metric,lem:rs-params,lem:split-words} & \texttt{Domain} & \texttt{IsSmoothDomain}, \texttt{card\_powDom\_and\_ker}, \texttt{fibre}, \texttt{RS}, \texttt{RS\_agree\_lt}, \texttt{RS\_unique\_decoding}, \texttt{evenW}, \texttt{oddW} \\
\Cref{thm:bciks} & \texttt{BCIKS} & axiom \texttt{bciks\_unique} \\
\Cref{lem:sequential,lem:queries,lem:randomised} & \texttt{Probability} & \texttt{sequential\_challenges}, \texttt{independent\_queries}, \texttt{randomised\_prover\_le} \\
\Cref{def:iop-sound,def:rbr,def:rbr-knowledge,lem:rbr-overall,lem:sumcheck,lem:rbr-binding} & \texttt{RoundByRound} & \texttt{RBRSound}, \texttt{RBRKnowledge}, \texttt{rbr\_knowledge\_overall}, \texttt{sumcheck\_sound\_iop}, \texttt{rbr\_knowledge\_binding} \\
\Cref{def:kernel,thm:basis,def:kernel-coords} & \texttt{Kernel} & \texttt{kernel}, \texttt{ofCoords}, \texttt{ofCoordsLin\_bijective}, \texttt{exists\_unique\_coords} \\
\Cref{sec:kernel:kron,sec:kernel:lowdeg,sec:kernel:eval} & \texttt{LowDegree}, \texttt{Mobius} & \texttt{coeff\_kernel}, \texttt{kron\_mul}, \texttt{Kmat\_inv}, \texttt{coeff\_ofCoords}, \texttt{lowDegree\_iff}, \texttt{RS\_eq\_charCond}, \texttt{eval\_ofCoords} \\
\Cref{prop:split,def:fold-poly,thm:fold-dictionary,cor:iterated} & \texttt{Kernel}, \texttt{WordFold} & \texttt{pfold}, \texttt{fold\_dictionary}, \texttt{pfoldList\_ofCoords}, \texttt{pfoldSeq\_ofCoords} \\
\Cref{def:fold-word,lem:fold-word,lem:line,prop:fold-consistency,cor:codeword-fold,def:cfold,prop:classical} & \texttt{WordFold} & \texttt{wfold}, \texttt{wfold\_line}, \texttt{wfold\_ev}, \texttt{codeword\_fold}, \texttt{cfold\_eq\_pfold} \\
\Cref{lem:encoding,lem:honest,thm:complete} & \texttt{Fibre}, \texttt{Protocol}, \texttt{Completeness} & \texttt{Enc\_injective}, \texttt{honest\_table}, \texttt{honS\_eval}, \texttt{completeness} \\
\Cref{def:fibre-dist,lem:fibre-dist,lem:unique,def:decoded,lem:far,lem:survive,def:good,lem:good,lem:step} & \texttt{Fibre} & \texttt{fibDist}, \texttt{fib\_unique}, \texttt{dec}, \texttt{decTable}, \texttt{far\_fold}, \texttt{survive}, \texttt{Good}, \texttt{ncard\_not\_good\_le}, \texttt{one\_step} \\
\Cref{lem:passing,lem:chain,thm:sound,cor:binding,prop:l0} & \texttt{Soundness} & \texttt{Pass\_card}, \texttt{chain}, \texttt{chain\_final}, \texttt{soundness\_far}, \texttt{soundness\_close}, \texttt{eval\_binding}, \texttt{no\_folding} \\
\Cref{def:state,lem:rbr-step,lem:rbr-sumcheck,thm:rbr,prop:rbr-l0} & \texttt{RBR}, \texttt{RBRGeneric} & \texttt{Doomed}, \texttt{rbr\_step}, \texttt{sCw\_eval}, \texttt{rbr\_round\_one}, \texttt{rbr\_round}, \texttt{rbr\_final}, \texttt{rbr\_knowledge}, \texttt{rbr\_l0} \\
\Cref{lem:sampling,def:kstate,lem:rbr-erasures,thm:rrbr,lem:cr-unique} & \texttt{NonInteractive} & \texttt{prob\_sample\_le}, \texttt{KState}, \texttt{rbr\_stepE}, \texttt{rbr\_erasures\_one}, \texttt{rrbr\_round}, \texttt{rrbr\_final}, \texttt{cr\_unique\_core} \\
\bottomrule
\end{longtable}
\end{center}

\paragraph{Not formalised.}\leavevmode\\
\Cref{thm:cdhz} and \cref{cor:pq}, which rest on the quantum analysis of Chiesa, Di, Hu and Zheng~\cite{CDHZ25}; the Merkle-tree statements (\cref{lem:merkle} and the collision step of \cref{lem:cr-unique}); Facts~\ref{fact:cyclic} and~\ref{fact:decode}, \cref{lem:descent} and \cref{rem:det}; running times and operation counts; \cref{sec:impl,sec:exp}.
In \cref{lem:rbr-erasures,thm:rrbr} the Lean statements give the per-round bounds for a model of transcripts by stages (empty, after $j$ rounds, full); they do not build a general interactive-oracle-reduction structure.

\paragraph{Findings of the formalisation.}
The formalisation found no false statement. It showed that \cref{lem:rbr-overall} (and hence \cref{lem:rbr-binding}) needs non-negative round errors, now stated in \cref{def:rbr,def:rbr-knowledge}; that case 2 of \cref{thm:sound} does not need its hypothesis $\Delta^{\mathrm{fib}}_0(w_0, \Cc_0) \le \delta$; and several hypotheses that are implied by the parameters (odd characteristic, $\ell \le n + R$, $N_j \le M_j$).


\end{document}
